% Les paratiques agropastoraleset artisanales traditionnelles — Géographie Tle Tech — Cours signé OnBuch+
% Programme officiel MINESEC. Compiler avec : tectonic N02-paratiques-agropastoraleset-artisanales-traditionnelles.tex
\newcommand{\DOCMATIERE}{Géographie}
\newcommand{\DOCNIVEAU}{Tle Tech}
\newcommand{\DOCMODULE}{Géographie – classe de Terminale technique}
\newcommand{\DOCLECON}{N02}
\newcommand{\DOCTITRE}{Les paratiques agropastoraleset artisanales traditionnelles}
% OnBuch+ — préambule « cours signés » ENRICHI (compilé avec tectonic / XeLaTeX)
% Système visuel « Pop » d'OnBuch+ : fond crème (jamais blanc), bordures encre
% épaisses, ombres « dures » décalées, titres Archivo Black en capitales.
% Version MATHÉMATIQUES : graphes (pgfplots/tikz), boîtes pédagogiques
% (définition, propriété, méthode, attention, le savais-tu, exercice résolu,
% exercice, TP, bilan), page de garde et signature OnBuch+ ancrée en bas.
%
% Macros attendues dans main.tex AVANT \input{preamble} :
%   \DOCMATIERE  \DOCNIVEAU  \DOCMODULE  \DOCLECON (numéro)  \DOCTITRE
\documentclass[11pt]{article}

\usepackage{fontspec}
\usepackage[french]{babel}
\usepackage[a4paper,margin=2cm,top=3.4cm,bottom=2.8cm,headheight=1.4cm,footskip=1.3cm]{geometry}
\usepackage{amsmath,amssymb,amsfonts}
\usepackage{mathtools} % \xmapsto, \coloneqq... souvent utilisés par les modèles
\usepackage{cancel} % simplifications barrées (bilans de forces)
% \abs{z} (module, valeur absolue) et \norm{u} (norme), utilisés par les modèles
\providecommand{\abs}{}\let\abs\relax\DeclarePairedDelimiter{\abs}{\lvert}{\rvert}
\providecommand{\norm}{}\let\norm\relax\DeclarePairedDelimiter{\norm}{\lVert}{\rVert}
\usepackage[table]{xcolor} % [table] : fournit aussi \rowcolors (alternance de lignes)
\usepackage{pagecolor}
\usepackage{tikz}
\usetikzlibrary{arrows.meta,positioning,calc,shapes.geometric,shapes.misc,shapes.arrows,decorations.pathmorphing,decorations.pathreplacing,decorations.markings,patterns,shadows,mindmap,trees,fit,backgrounds,matrix,intersections,angles,quotes}
\usepackage{pgfplots}
\usepackage[version=4]{mhchem} % équations nucléaires \ce{^{235}_{92}U}
\usepackage[european,siunitx]{circuitikz} % schémas électriques
\pgfplotsset{compat=1.18}
\usepgfplotslibrary{fillbetween,groupplots} % aires entre courbes (calcul intégral)
\usepackage{enumitem}
\usepackage{fancyhdr}
% [most] charge toutes les bibliothèques tcolorbox (listings, minted, raster,
% documentation...) et gonfle inutilement la mémoire TeX sur un document long
% (~300 boîtes) : on ne charge que ce qui est réellement utilisé.
\usepackage[skins,breakable]{tcolorbox}
\usepackage{graphicx}
\usepackage{colortbl}
\usepackage{booktabs}
\usepackage{tabularx}
\usepackage{multirow}
\usepackage{diagbox} % case d'en-tête diagonale des tableaux à double entrée
% Colonnes extensibles centrée / gauche / droite (C, L, R), souvent utilisées par les modèles
\newcolumntype{C}{>{\centering\arraybackslash}X}
\newcolumntype{L}{>{\raggedright\arraybackslash}X}
\newcolumntype{R}{>{\raggedleft\arraybackslash}X}
\usepackage{array}
\usepackage{multicol}
\usepackage{siunitx}
% parse-numbers=false : accepte aussi \SI{2,5 \times 10^{-3}}{...}
\sisetup{locale=FR,output-decimal-marker={,},group-minimum-digits=4,parse-numbers=false}
\DeclareSIUnit\ohm{\ensuremath{\symup{\Omega}}} % U+2126 absent de la police
\DeclareSIUnit\division{div} % réglages d'oscilloscope (ms/div, V/div)
\DeclareSIUnit\tr{tr} % tours (vitesse de rotation, ex. \SI{1500}{\tr\per\minute})
\DeclareSIUnit\molar{M} % concentration molaire (ex. \SI{3}{\molar}, \SI{1}{\micro\molar})
\DeclareSIUnit\an{an} % année (le modèle confond parfois avec \year, un compteur TeX) — ex. \SI{5}{\kilo\meter\per\an}
\DeclareSIUnit\FCFA{FCFA} % franc CFA (ex. \SI{500}{\FCFA\per\kilo\gram})
\DeclareSIUnit\degreeBrix{\textdegree Brix} % teneur en sucre d'un sirop/jus (réfractométrie)
\DeclareSIUnit\month{mois} % le modèle utilise parfois \month, un compteur TeX, comme unité de durée
\DeclareSIUnit\microliter{\micro\liter} % le modèle écrit parfois \microliter en un seul token
\DeclareSIUnit\million{millions} % dénombrement (ex. \SI{15}{\million\per\mL} spermatozoïdes)
\usepackage{qrcode}
% \degree : commande fréquemment utilisée par les modèles pour un angle ou une
% température en dehors de \SI{}{} (non fournie par les paquets chargés).
\providecommand{\degree}{^{\circ}}
% \qed (fin de démonstration), utilisé par les modèles sans amsthm
\providecommand{\qed}{\hfill$\square$}
% \barre : double barre des valeurs interdites dans un tableau de variations (tkz-tab)
\providecommand{\barre}{\ensuremath{\|}}
% environnement proof (amsthm non chargé) utilisé par les modèles
\ifdefined\proof\else\newenvironment{proof}[1][Démonstration]{\par\noindent\textit{#1.}\ }{\qed\par}\fi
\usepackage{lastpage}
\usepackage{hyperref}
\hypersetup{hidelinks}

% ============================================================
% Palette « Pop » OnBuch+ — mêmes hex que la classe P (plus_ui.dart)
% ============================================================
\definecolor{poporange}{HTML}{F59321}
\definecolor{popdark}{HTML}{E07A0C}
\definecolor{popcream}{HTML}{FFF6E8}
\definecolor{popink}{HTML}{1C1714}
\definecolor{poppurple}{HTML}{7A5AE0}
\definecolor{popgreen}{HTML}{2E9E6A}
\definecolor{poppink}{HTML}{E0457B}
\definecolor{popblue}{HTML}{2F7DE1}
\definecolor{popgold}{HTML}{C98A12}
\definecolor{popmuted}{HTML}{8A7F76}
\definecolor{popline}{HTML}{EADFD2}
\definecolor{popgray}{HTML}{8A7F76}
\colorlet{popgrey}{popgray}
% Alias tolérants (le modèle nomme parfois les couleurs différemment)
\colorlet{popwhite}{white}
\colorlet{popblack}{popink}
\colorlet{popred}{poppink}
\colorlet{popyellow}{popgold}
% Teintes claires (fonds de tableaux/encadrés) — le modèle nomme parfois
% les couleurs avec un suffixe L (ex. popblueL) sans les avoir définies.
\colorlet{poporangeL}{poporange!20}
\colorlet{popdarkL}{popdark!20}
\colorlet{poppurpleL}{poppurple!20}
\colorlet{popgreenL}{popgreen!20}
\colorlet{poppinkL}{poppink!20}
\colorlet{popblueL}{popblue!20}
\colorlet{popgoldL}{popgold!20}
\colorlet{popredL}{popred!20}
\colorlet{popgrayL}{popgray!20}
% Teintes claires pour fonds de tableaux / zones
\colorlet{popblueL}{popblue!10!white}
\colorlet{poporangeL}{poporange!14!white}
\colorlet{popgreenL}{popgreen!12!white}
\colorlet{poppurpleL}{poppurple!12!white}
\colorlet{poppinkL}{poppink!12!white}
\colorlet{popgoldL}{popgold!14!white}

\pagecolor{popcream}

% ============================================================
% Typographie
% ============================================================
\defaultfontfeatures{Ligatures=TeX}
\setmainfont{PlusJakartaSans-Medium.ttf}[Path=../../../fonts/,BoldFont=PlusJakartaSans-Bold.ttf]
% Maths en sans-serif (Fira Math) avec chiffres et lettres droites de Plus
% Jakarta, pour que les formules se fondent dans le texte.
\usepackage[math-style=ISO,bold-style=ISO]{unicode-math}
\setmathfont{FiraMath-Regular.otf}[Path=../../../fonts/]
\setmathfont{PlusJakartaSans-Medium.ttf}[Path=../../../fonts/,range={up/num,up/{Latin,latin},\mathup}]
\usepackage{icomma} % virgule décimale sans espace en mode math
% Caractères mathématiques Unicode absents des polices de texte (Plus Jakarta,
% Archivo Black) : rendus par leur équivalent mathématique, en texte comme en
% formule — sinon ils s'affichent en carré vide (ex. « Arithmétique dans ℤ »).
\usepackage{newunicodechar}
\newunicodechar{ℝ}{\ensuremath{\mathbb{R}}}
\newunicodechar{ℤ}{\ensuremath{\mathbb{Z}}}
\newunicodechar{ℂ}{\ensuremath{\mathbb{C}}}
\newunicodechar{ℕ}{\ensuremath{\mathbb{N}}}
\newunicodechar{ℚ}{\ensuremath{\mathbb{Q}}}
\newunicodechar{𝔻}{\ensuremath{\mathbb{D}}}
\newunicodechar{α}{\ensuremath{\alpha}}
\newunicodechar{β}{\ensuremath{\beta}}
\newunicodechar{γ}{\ensuremath{\gamma}}
\newunicodechar{δ}{\ensuremath{\delta}}
\newunicodechar{ε}{\ensuremath{\varepsilon}}
\newunicodechar{θ}{\ensuremath{\theta}}
\newunicodechar{λ}{\ensuremath{\lambda}}
\newunicodechar{μ}{\ensuremath{\mu}}
\newunicodechar{σ}{\ensuremath{\sigma}}
\newunicodechar{τ}{\ensuremath{\tau}}
\newunicodechar{φ}{\ensuremath{\varphi}}
\newunicodechar{ψ}{\ensuremath{\psi}}
\newunicodechar{ω}{\ensuremath{\omega}}
\newunicodechar{Σ}{\ensuremath{\Sigma}}
% majuscules produites par \MakeUppercase dans les titres (α-aminés -> Α)
\newunicodechar{Α}{\ensuremath{\alpha}}
\newunicodechar{Β}{\ensuremath{\beta}}
\newunicodechar{Γ}{\ensuremath{\Gamma}}
\newunicodechar{Δ}{\ensuremath{\Delta}}
\newunicodechar{Ε}{\ensuremath{\varepsilon}}
\newunicodechar{Θ}{\ensuremath{\Theta}}
\newunicodechar{Λ}{\ensuremath{\Lambda}}
\newunicodechar{Μ}{\ensuremath{\mu}}
\newunicodechar{Τ}{\ensuremath{\tau}}
\newunicodechar{Φ}{\ensuremath{\Phi}}
\newunicodechar{Ψ}{\ensuremath{\Psi}}
\newunicodechar{Ω}{\ensuremath{\Omega}}
\newunicodechar{↦}{\ensuremath{\mapsto}}
\newunicodechar{∈}{\ensuremath{\in}}
\newunicodechar{∉}{\ensuremath{\notin}}
\newunicodechar{∧}{\ensuremath{\wedge}}
\newunicodechar{⇔}{\ensuremath{\Leftrightarrow}}
\newunicodechar{⟺}{\ensuremath{\Longleftrightarrow}}
\newunicodechar{⇒}{\ensuremath{\Rightarrow}}
\newunicodechar{⟹}{\ensuremath{\Longrightarrow}}
\newunicodechar{∘}{\ensuremath{\circ}}
\newunicodechar{∪}{\ensuremath{\cup}}
\newunicodechar{∩}{\ensuremath{\cap}}
\newunicodechar{≡}{\ensuremath{\equiv}}
\newunicodechar{⊂}{\ensuremath{\subset}}
\newunicodechar{⊥}{\ensuremath{\perp}}
\newunicodechar{∃}{\ensuremath{\exists}}
\newunicodechar{∀}{\ensuremath{\forall}}
\newunicodechar{∛}{\ensuremath{\sqrt[3]{\,}}}
\newunicodechar{∜}{\ensuremath{\sqrt[4]{\,}}}
\newunicodechar{⃗}{\ensuremath{{}^{\to}}}
\newunicodechar{ⁿ}{\ifmmode^{n}\else\textsuperscript{\ensuremath{n}}\fi}
\newunicodechar{ˣ}{\ifmmode^{x}\else\textsuperscript{\ensuremath{x}}\fi}
\newunicodechar{ᵇ}{\ifmmode^{b}\else\textsuperscript{\ensuremath{b}}\fi}
\newunicodechar{ʳ}{\ifmmode^{r}\else\textsuperscript{\ensuremath{r}}\fi}
\newunicodechar{ᵘ}{\ifmmode^{u}\else\textsuperscript{\ensuremath{u}}\fi}
\newunicodechar{ʸ}{\ifmmode^{y}\else\textsuperscript{\ensuremath{y}}\fi}
\newunicodechar{ᵃ}{\ifmmode^{a}\else\textsuperscript{\ensuremath{a}}\fi}
\newunicodechar{ᵉ}{\ifmmode^{e}\else\textsuperscript{\ensuremath{e}}\fi}
\newunicodechar{ᵏ}{\ifmmode^{k}\else\textsuperscript{\ensuremath{k}}\fi}
\newunicodechar{ᵖ}{\ifmmode^{p}\else\textsuperscript{\ensuremath{p}}\fi}
\newunicodechar{ᴺ}{\ifmmode^{N}\else\textsuperscript{\ensuremath{N}}\fi}
\newunicodechar{ᶜ}{\ifmmode^{c}\else\textsuperscript{\ensuremath{c}}\fi}
\newunicodechar{ʰ}{\ifmmode^{h}\else\textsuperscript{\ensuremath{h}}\fi}
\newunicodechar{ᵠ}{\ifmmode^{q}\else\textsuperscript{\ensuremath{q}}\fi}
\newunicodechar{ₙ}{\ifmmode_{n}\else\textsubscript{\ensuremath{n}}\fi}
\newunicodechar{ₐ}{\ifmmode_{a}\else\textsubscript{\ensuremath{a}}\fi}
\newunicodechar{ᵢ}{\ifmmode_{i}\else\textsubscript{\ensuremath{i}}\fi}
\newunicodechar{ᵣ}{\ifmmode_{r}\else\textsubscript{\ensuremath{r}}\fi}
\newunicodechar{ₖ}{\ifmmode_{k}\else\textsubscript{\ensuremath{k}}\fi}
\newunicodechar{ᵦ}{\ifmmode_{\beta}\else\textsubscript{\ensuremath{\beta}}\fi}
\newunicodechar{ₓ}{\ifmmode_{x}\else\textsubscript{\ensuremath{x}}\fi}
\newfontfamily\popdisplay{ArchivoBlack-Regular.ttf}[Path=../../../fonts/]
\newfontfamily\poplogo{SpaceGrotesk-Bold.ttf}[Path=../../../fonts/]
\newfontfamily\popsemi{PlusJakartaSans-SemiBold.ttf}[Path=../../../fonts/]
\newfontfamily\popxbold{PlusJakartaSans-ExtraBold.ttf}[Path=../../../fonts/]
\newfontfamily\popxboldls{PlusJakartaSans-ExtraBold.ttf}[Path=../../../fonts/,LetterSpace=3.0]

\setlist[enumerate]{leftmargin=1.7em,itemsep=5pt,topsep=4pt}
\setlist[itemize]{leftmargin=1.7em,itemsep=4pt,topsep=4pt,label={\color{poporange}\textbullet}}
\setlength{\parindent}{0pt}
\setlength{\parskip}{5pt}
\renewcommand{\arraystretch}{1.25}

% pgfplots : style Pop par défaut
\pgfplotsset{
  every axis/.append style={axis line style={popink,line width=1pt},tick style={popink},
    grid style={popline,line width=0.5pt},label style={font=\small},tick label style={font=\footnotesize}},
  popcurve/.style={poporange,line width=1.8pt,no markers,smooth},
}
% Même style disponible dans un \draw TikZ simple (hors pgfplots)
\tikzset{popcurve/.style={poporange,line width=1.8pt}}
\tikzset{popgrid/.style={popline,very thin}}
\colorlet{popcurve}{poporange} % le nom du style est parfois utilisé comme couleur

% ============================================================
% Pill / squiggle
% ============================================================
\newcommand{\poppill}[3]{%
  \tikz[baseline=-0.55ex]\node[draw=popink,line width=1.2pt,rounded corners=2.6pt,%
    fill=#2,inner xsep=6.5pt,inner ysep=3pt]{%
    \popxboldls\fontsize{7}{7}\selectfont\color{#3}\MakeUppercase{#1}};%
}
\newcommand{\popsquiggle}[1]{%
  \tikz[baseline=-0.4ex]{
    \draw[#1,line width=2.6pt,line cap=round]
      plot [smooth] coordinates {(0,0.09) (0.34,0.01) (0.68,0.21) (1.02,0.01) (1.36,0.10)};
  }%
}
% Mot-clé surligné (marqueur orange sous le texte)
\newcommand{\cle}[1]{{\popxbold\color{popdark}#1}}

% ============================================================
% En-tête / pied
% ============================================================
\pagestyle{fancy}
\fancyhf{}
\renewcommand{\headrulewidth}{0pt}
\renewcommand{\footrulewidth}{0pt}
\lhead{\raisebox{-0.15ex}{\poplogo\fontsize{13}{13}\selectfont\color{popink}OnBuch\color{poporange}+}}
\rhead{\poppill{\DOCMATIERE\ \textbullet\ \DOCNIVEAU\ \textbullet\ Leçon \DOCLECON}{white}{popink}}
\lfoot{\popxbold\fontsize{7.6}{7.6}\selectfont\color{popmuted}\MakeUppercase{Cours signé OnBuch+}\ \textbullet\ {\popsemi\DOCTITRE}}
\rfoot{\tikz[baseline=-0.6ex]\node[rounded corners=8.5pt,fill=popink,minimum height=17pt,minimum width=17pt,inner xsep=5pt,inner sep=0]{\popxbold\fontsize{8}{8}\selectfont\color{popcream}\thepage\,/\,\pageref*{LastPage}};}

% ============================================================
% Titres
% ============================================================
\newcommand{\chaptitre}[2]{%
  \begin{tcolorbox}[enhanced,colback=poporange,colframe=popink,boxrule=2.6pt,arc=11pt,
      drop shadow=popink,left=15pt,right=15pt,top=12pt,bottom=11pt,before skip=3pt,after skip=18pt]
    \poppill{#2}{popink}{popcream}\par\vspace{9pt}
    {\raggedright\hyphenpenalty=10000\exhyphenpenalty=10000\popdisplay\fontsize{22}{24}\selectfont\color{popcream}\MakeUppercase{#1}\par}
  \end{tcolorbox}
}
% Section numérotée : pastille orange avec le numéro + titre Archivo
\newcounter{popsec}
\newcommand{\coursec}[1]{%
  \stepcounter{popsec}\setcounter{popsub}{0}%
  \par\vspace{16pt}\noindent
  \begin{minipage}{\linewidth}
  \tikz[baseline=-0.7ex]\node[draw=popink,line width=1.6pt,fill=poporange,rounded corners=4pt,minimum size=22pt,inner sep=0]{\popdisplay\fontsize{12}{12}\selectfont\color{popcream}\thepopsec};%
  \hspace{8pt}{\popdisplay\fontsize{15}{17}\selectfont\color{popink}\MakeUppercase{#1}}\par
  \vspace{4pt}\noindent{\color{poporange}\rule{36pt}{3.6pt}}
  \end{minipage}\par\nopagebreak\vspace{8pt}
}
\newcounter{popsub}
\newcommand{\courssub}[1]{%
  \stepcounter{popsub}%
  \par\vspace{9pt}\noindent{\popxbold\fontsize{11.5}{13}\selectfont\color{popink}{\color{poporange}\thepopsec.\thepopsub}\hspace{6pt}#1}\par\nopagebreak\vspace{4pt}
}

% ============================================================
% Boîtes pédagogiques — même recette : fond blanc, bordure épaisse colorée,
% ombre dure encre, badge plein. Titre optionnel [#1].
% ============================================================
\tcbset{popbase/.style={breakable,enhanced,colback=white,boxrule=2.3pt,arc=9pt,
  shadow={3pt}{-3pt}{0pt}{popink},left=14pt,right=14pt,top=11pt,bottom=11pt,before skip=11pt,after skip=15pt,
  fontupper=\popsemi\fontsize{10.6}{14.5}\selectfont\color{popink},
  fontlower=\popsemi\fontsize{10.6}{14.5}\selectfont\color{popink}}}
% Coche dessinée (le glyphe ✓ n'existe pas dans les polices du document)
\newcommand{\popcheck}[1]{\tikz[baseline=-0.5ex]\draw[#1,line width=1.6pt,line cap=round,line join=round](-0.13,0.0)--(-0.04,-0.09)--(0.13,0.1);}
\providecommand{\coche}{\popcheck{popgreen}}
\providecommand{\resultat}[1]{\par\smallskip\noindent\fcolorbox{popgreen}{popgreenL}{\parbox{\dimexpr\linewidth-2\fboxsep-2\fboxrule\relax}{\textbf{Résultat.} #1}}\par\smallskip}
\newcommand{\popboxhead}[3]{\poppill{#1}{#2}{white}\ifx\relax#3\relax\else\hspace{6pt}{\popxbold\fontsize{10.6}{12}\selectfont\color{popink}#3}\fi\par\vspace{7pt}}

\newtcolorbox{aretenir}[1][]{popbase,colframe=popgreen,before upper={\popboxhead{À retenir}{popgreen}{#1}}}
\newtcolorbox{exemplebox}[1][]{popbase,colframe=poporange,before upper={\popboxhead{Exemple}{poporange}{#1}}}
\newtcolorbox{definition}[1][]{popbase,colframe=popblue,before upper={\popboxhead{Définition}{popblue}{#1}}}
\newtcolorbox{propriete}[1][]{popbase,colframe=poppurple,before upper={\popboxhead{Propriété}{poppurple}{#1}}}
\newtcolorbox{methode}[1][]{popbase,colframe=popgold,before upper={\popboxhead{Méthode}{popgold}{#1}}}
\newtcolorbox{attention}[1][]{popbase,colframe=poppink,before upper={\popboxhead{Attention}{poppink}{#1}}}
\newtcolorbox{experience}[1][]{popbase,colframe=popblue,colback=popblueL,before upper={\popboxhead{Expérience}{popblue}{#1}}}
% « Le savais-tu ? » : carte violette pleine (contexte, culture, Cameroun)
\newtcolorbox{savaistu}[1][]{popbase,colback=poppurple,colframe=popink,
  fontupper=\popsemi\fontsize{10.4}{14.2}\selectfont\color{white},
  before upper={\poppill{Le savais-tu\,?}{white}{poppurple}\ifx\relax#1\relax\else\hspace{6pt}{\popxbold\color{white}#1}\fi\par\vspace{7pt}}}
% Exercice résolu : énoncé en haut, \tcblower puis la solution sur fond crème
\newcounter{popexo}\newcounter{popexr}
\newtcolorbox{exoresolu}[1][]{popbase,colframe=poporange,colbacklower=popcream,
  segmentation style={popink,line width=1.2pt,dashed},
  before upper={\stepcounter{popexr}\popboxhead{Exercice résolu \thepopexr}{poporange}{#1}},
  before lower={\poppill{Solution}{popink}{popcream}\par\vspace{6pt}}}
% Exercice (non résolu en place) — la correction est dans la partie Corrigés
\newtcolorbox{exercice}[1][]{popbase,colframe=popink,
  before upper={\stepcounter{popexo}\popboxhead{Exercice \thepopexo}{popink}{#1}}}
% Corrigé
\newtcolorbox{corrige}[1][]{popbase,colframe=popgreen,colback=popgreenL,
  before upper={\popboxhead{Corrigé}{popgreen}{#1}}}
% Objectifs / prérequis / bilan
\newtcolorbox{objectifs}{popbase,colframe=popink,colback=poporangeL,
  before upper={\popboxhead{Objectifs de la leçon}{popink}{}}}
\newtcolorbox{prerequis}{popbase,colframe=popink,colback=popblueL,
  before upper={\popboxhead{Prérequis}{popblue}{}}}
\newtcolorbox{bilan}{popbase,colframe=popink,colback=white,boxrule=2.8pt,
  before upper={\popboxhead{Fiche bilan}{poporange}{L'essentiel à connaître pour l'examen}}}
% Figure encadrée avec légende
\newtcolorbox{popfigure}[1][]{enhanced,breakable=false,colback=white,colframe=popline,boxrule=1.4pt,arc=8pt,
  left=10pt,right=10pt,top=10pt,bottom=8pt,before skip=10pt,after skip=12pt,halign=center,
  lower separated=false,halign lower=center,
  fontlower=\popsemi\fontsize{9}{11.5}\selectfont\color{popmuted},
  #1}
\newcommand{\legende}[1]{\par\vspace{6pt}{\popsemi\fontsize{9}{11.5}\selectfont\color{popmuted}#1}}

% ============================================================
% Signature OnBuch+
% ============================================================
\newcommand{\poplogoinline}[1]{{\poplogo\fontsize{#1}{#1}\selectfont\color{popink}OnBuch\color{poporange}+}}
\newcommand{\signatureonbuch}{%
  \begin{tcolorbox}[enhanced,colback=popink,colframe=popink,boxrule=2.2pt,arc=10pt,
      drop shadow=poporange,left=14pt,right=14pt,top=10pt,bottom=10pt,before skip=0pt,after skip=0pt]
    \tikz[baseline=-0.7ex]\node[circle,fill=popgreen,draw=popcream,line width=1.3pt,minimum size=22pt,inner sep=0]{\popcheck{white}};%
    \hspace{9pt}\begin{minipage}[c]{0.62\linewidth}
      {\popxbold\fontsize{12}{13}\selectfont\color{popcream}Signé\ \textbullet\ {\poplogo OnBuch\color{poporange}+}}\par\vspace{2pt}
      {\raggedright\popsemi\fontsize{8}{10.5}\selectfont\color{popline}\MakeUppercase{Équipe pédagogique OnBuch+}\\\MakeUppercase{Conforme au programme officiel MINESEC}\par}
    \end{minipage}\hfill
    {\poplogo\fontsize{22}{22}\selectfont\color{popcream}OnBuch\color{poporange}+}
  \end{tcolorbox}%
}

% ============================================================
% Page de garde
% #1 titre · #2 matière · #3 niveau · #4 module · #5 n° leçon · #6 durée ·
% #7 description / objectifs (texte ou itemize)
% ============================================================
\newcommand{\pagegarde}[7]{%
  \thispagestyle{empty}
  \newgeometry{a4paper,margin=1.7cm,top=1.6cm,bottom=1.4cm}%
  \begin{tikzpicture}[remember picture,overlay]
    \fill[poporange!18] ([xshift=-3.2cm,yshift=-3.6cm]current page.north east) circle (4.6cm);
    \fill[poppurple!14] ([xshift=2.4cm,yshift=7.6cm]current page.south west) circle (3.8cm);
  \end{tikzpicture}%
  {\poplogo\fontsize{30}{30}\selectfont\color{popink}OnBuch\color{poporange}+}\hfill
  \raisebox{0.6ex}{\poppill{Cours signé \textbullet\ Premium}{popink}{popcream}}\par
  \vspace{1pt}\popsquiggle{poppurple}\par
  \vspace{8pt}
  \begin{tcolorbox}[enhanced,colback=poporange,colframe=popink,boxrule=2.8pt,arc=13pt,
      drop shadow=popink,left=18pt,right=18pt,top=12pt,bottom=13pt,after skip=10pt]
    \poppill{#2 \textbullet\ #3}{popink}{popcream}\hspace{5pt}\poppill{#4}{white}{popink}\par\vspace{8pt}
    {\popdisplay\fontsize{14}{14}\selectfont\color{popink}LEÇON #5}\par\vspace{4pt}
    {\raggedright\hyphenpenalty=10000\exhyphenpenalty=10000\popdisplay\fontsize{25}{27}\selectfont\color{popcream}\MakeUppercase{#1}\par}
    \vspace{8pt}
    {\popxbold\fontsize{9.5}{9.5}\selectfont\color{popink}\MakeUppercase{Durée conseillée : #6}}
  \end{tcolorbox}
  \begin{tcolorbox}[enhanced,colback=white,colframe=popink,boxrule=2.2pt,arc=10pt,
      drop shadow=popink,left=16pt,right=16pt,top=10pt,bottom=10pt,after skip=8pt]
    \poppill{Dans cette leçon}{poppurple}{white}\par\vspace{5pt}
    \popsemi\fontsize{9.3}{12.6}\selectfont\color{popink}#7
  \end{tcolorbox}
  \noindent\begin{minipage}[t]{0.487\linewidth}
    \begin{tcolorbox}[enhanced,colback=popgreen,colframe=popink,boxrule=2.2pt,arc=10pt,
        drop shadow=popink,left=11pt,right=11pt,top=10pt,bottom=10pt,height=3.1cm,valign=center]
      \tcbox[colback=white,colframe=popink,boxrule=1.3pt,arc=5pt,boxsep=0pt,nobeforeafter,
        left=3pt,right=3pt,top=3pt,bottom=3pt,valign=center]{\qrcode[height=1.9cm]{https://play.google.com/store/apps/details?id=com.luvvix.onbuch&hl=fr}}%
      \hspace{8pt}\begin{minipage}[c]{0.5\linewidth}
        {\popdisplay\fontsize{11}{12.5}\selectfont\color{white}\MakeUppercase{Télécharge OnBuch}}\par\vspace{4pt}
        {\raggedright\popsemi\fontsize{8.4}{11}\selectfont\color{white}Gratuit \textbullet\ Google Play\\Tuteur IA, annales, concours, résultats\par}
      \end{minipage}
    \end{tcolorbox}
  \end{minipage}\hfill
  \begin{minipage}[t]{0.487\linewidth}
    \begin{tcolorbox}[enhanced,colback=poppurple,colframe=popink,boxrule=2.2pt,arc=10pt,
        drop shadow=popink,left=11pt,right=11pt,top=10pt,bottom=10pt,height=3.1cm,valign=center]
      \tikz[baseline=-0.7ex]\node[circle,fill=white,draw=popink,line width=1.3pt,minimum size=1.6cm,inner sep=0]{\popdisplay\fontsize{24}{24}\selectfont\color{poppurple}+};%
      \hspace{8pt}\begin{minipage}[c]{0.6\linewidth}
        {\popdisplay\fontsize{11}{12.5}\selectfont\color{white}\MakeUppercase{Passe à OnBuch+}}\par\vspace{4pt}
        {\raggedright\popsemi\fontsize{8.4}{11}\selectfont\color{white}Tous les cours signés, Tuteur IA illimité, fiches PDF — onglet OnBuch+ de l'app\par}
      \end{minipage}
    \end{tcolorbox}
  \end{minipage}
  \vspace{8pt}
  \popcontenu
  \vfill
  \signatureonbuch
  \restoregeometry
  \newpage
}

% Rangée « Ce cours contient » (page de garde)
\newcommand{\poptile}[3]{% #1 couleur #2 gros texte #3 légende
  \begin{tcolorbox}[enhanced,colback=#1,colframe=popink,boxrule=2pt,arc=9pt,shadow={2.5pt}{-2.5pt}{0pt}{popink},
      left=6pt,right=6pt,top=9pt,bottom=9pt,halign=center,valign=center,height=1.75cm,width=\linewidth,nobeforeafter]
    {\popdisplay\fontsize{12.5}{13}\selectfont\color{white}\MakeUppercase{#2}}\par\vspace{3pt}
    {\centering\popsemi\fontsize{7.8}{9.5}\selectfont\color{white}#3\par}
  \end{tcolorbox}}
\newcommand{\popcontenu}{%
  \par\noindent{\popxboldls\fontsize{7.5}{7.5}\selectfont\color{popmuted}CE COURS SIGNÉ CONTIENT}\par\vspace{7pt}
  \noindent\begin{minipage}[t]{0.235\linewidth}\poptile{popblue}{Cours}{complet, illustré, pas à pas}\end{minipage}\hfill
  \begin{minipage}[t]{0.235\linewidth}\poptile{poporange}{Méthodes}{et exercices résolus}\end{minipage}\hfill
  \begin{minipage}[t]{0.235\linewidth}\poptile{poppink}{Exercices}{type examen + corrigés}\end{minipage}\hfill
  \begin{minipage}[t]{0.235\linewidth}\poptile{popgreen}{Fiche bilan}{l'essentiel à retenir}\end{minipage}\par}

% Sommaire
\newtcolorbox{sommaire}{enhanced,colback=white,colframe=popink,boxrule=2.2pt,arc=10pt,
  drop shadow=popink,left=18pt,right=18pt,top=13pt,bottom=13pt,before skip=6pt,after skip=20pt,
  before upper={\poppill{Sommaire}{poppurple}{white}\par\vspace{9pt}\popsemi\fontsize{10.6}{14.5}\selectfont\color{popink}}}

% Signature de fin de cours — poussée tout en bas de la dernière page
\newcommand{\findecours}{%
  \par\vfill\nopagebreak
  \begin{center}\popsquiggle{poporange}\end{center}\vspace{4pt}
  \signatureonbuch
}
\AtBeginDocument{\def\llbracket{\mathopen{[\mkern-3.5mu[}}\def\rrbracket{\mathclose{]\mkern-3.5mu]}}} % intervalles d'entiers (glyphe ⟦ absent de la police)
% Glyphes absents de Fira Math : redéfinis à partir de glyphes disponibles
\AtBeginDocument{%
  \def\setminus{\mathbin{\backslash}}\let\smallsetminus\setminus
  \def\vdots{\mathord{\vbox{\baselineskip3.2pt\lineskiplimit0pt\kern4pt\hbox{.}\hbox{.}\hbox{.}}}}%
  \def\ddots{\mathinner{\mkern1mu\raise6.5pt\hbox{.}\mkern2mu\raise3.6pt\hbox{.}\mkern2mu\raise0.7pt\hbox{.}\mkern1mu}}%
  \def\triangle{\mathord{\scalebox{0.9}{$\symup{\Delta}$}}}%
  \def\nabla{\mathord{\raisebox{\depth}{\scalebox{1}[-1]{$\symup{\Delta}$}}}}%
  \def\checkmark{\popcheck{popdark}}%
  \def\star{\mathbin{\ast}}\def\bigstar{\mathord{\ast}}%
  \def\diamond{\mathbin{\scalebox{0.6}{$\Diamond$}}}%
  \def\bigcup{\mathop{\mathchoice{\raisebox{-0.3em}{\scalebox{1.7}{$\cup$}}}{\scalebox{1.25}{$\cup$}}{\cup}{\cup}}\displaylimits}%
  \def\bigcap{\mathop{\mathchoice{\raisebox{-0.3em}{\scalebox{1.7}{$\cap$}}}{\scalebox{1.25}{$\cap$}}{\cap}{\cap}}\displaylimits}%
  \def\lesseqqgtr{\mathrel{\lessgtr}}\def\bigoplus{\mathop{\oplus}\displaylimits}%
}
\providecommand{\ssi}{si et seulement si} % abréviation fréquente
\AtBeginDocument{\providecommand{\wideparen}{\overparen}} % arc de cercle

%% FIGFIX-BEGIN v3 — correctifs globaux des figures (docs/onbuch-generation/tools/figfix)
\usepackage{adjustbox}\usepackage{etoolbox}
% toute figure TikZ/pgfplots plus large que la ligne est réduite à la largeur disponible
\BeforeBeginEnvironment{tikzpicture}{\begin{adjustbox}{max width=\linewidth}}
\AfterEndEnvironment{tikzpicture}{\end{adjustbox}}
% légendes pgfplots lisibles par-dessus les courbes
\pgfplotsset{every axis/.append style={legend style={fill=white,fill opacity=0.92,text opacity=1,draw=black!15,inner sep=3pt}}}
% étiquettes d'arêtes (syntaxe edge["..."]) avec fond blanc
\tikzset{every edge quotes/.append style={fill=white,inner sep=1.2pt}}
% bulles de cartes mentales : texte à la ligne dans la bulle, bulle un peu plus grande
\tikzset{every concept/.append style={text width=2.0cm,align=center,minimum size=2.3cm,inner sep=2pt}}
%% FIGFIX-END
\begin{document}

\pagegarde{Les paratiques agropastoraleset artisanales traditionnelles}{Géographie}{Tle Tech}{Géographie – classe de Terminale technique}{N02}{2 séances (fiche de progression harmonisée)}{Cette leçon analyse les mutations des systèmes agropastoraux et artisanaux traditionnels au Cameroun, en mobilisant la cartographie thématique et l'analyse de documents pour comprendre l'adaptation des acteurs aux contraintes modernes.
\vspace{4pt}
\begin{multicols}{2}\small
\begin{itemize}
\item À la fin de cette leçon, je sais caractériser les systèmes d'élevage et de pêche traditionnels au Cameroun (types, itinéraires, techniques).
\item À la fin de cette leçon, je sais expliquer les facteurs de transformation de l'élevage traditionnel (sédentarisation, conflits agro-pastoraux, modernisation).
\item À la fin de cette leçon, je sais analyser l'évolution de la pêche artisanale face à la motorisation, la surpêche et la réglementation.
\item À la fin de cette leçon, je sais identifier les principales filières de l'artisanat traditionnel et leur ancrage culturel et spatial.
\item À la fin de cette leçon, je sais évaluer l'impact de la commercialisation et du tourisme sur l'authenticité et la transmission des savoir-faire artisanaux.
\item À la fin de cette leçon, je sais construire un croquis légendé de la répartition des zones d'élevage, de pêche et d'artisanat au Cameroun.
\end{itemize}\end{multicols}}

\begin{sommaire}
\begin{multicols}{2}
\begin{enumerate}
\item 1. L'élevage traditionnel : diversité des systèmes et organisation de l'espace
\item 2. Transformations de l'élevage traditionnel : modernisation, conflits et politiques publiques
\item 3. La pêche traditionnelle : techniques, espaces et mutations
\item 4. L'artisanat traditionnel : diversité, organisation et évolution
\item 5. Synthèse : Enjeux transversaux et développement durable des pratiques traditionnelles
\item Activité d'intégration
\item Méthodes et exercices résolus
\item Exercices
\item Corrigés des exercices
\item Fiche bilan
\end{enumerate}
\end{multicols}
\end{sommaire}

\begin{objectifs}
\begin{itemize}
\item À la fin de cette leçon, je sais caractériser les systèmes d'élevage et de pêche traditionnels au Cameroun (types, itinéraires, techniques).
\item À la fin de cette leçon, je sais expliquer les facteurs de transformation de l'élevage traditionnel (sédentarisation, conflits agro-pastoraux, modernisation).
\item À la fin de cette leçon, je sais analyser l'évolution de la pêche artisanale face à la motorisation, la surpêche et la réglementation.
\item À la fin de cette leçon, je sais identifier les principales filières de l'artisanat traditionnel et leur ancrage culturel et spatial.
\item À la fin de cette leçon, je sais évaluer l'impact de la commercialisation et du tourisme sur l'authenticité et la transmission des savoir-faire artisanaux.
\item À la fin de cette leçon, je sais construire un croquis légendé de la répartition des zones d'élevage, de pêche et d'artisanat au Cameroun.
\item À la fin de cette leçon, je sais rédiger une argumentation structurée justifiant l'intégration des pratiques traditionnelles dans les politiques de développement durable.
\end{itemize}
\end{objectifs}

\begin{prerequis}
\begin{itemize}
\item Notion de milieu physique et d'adaptation de l'homme (classe de 4ème / 3ème).
\item Repérage des grandes régions naturelles du Cameroun (Soudan, Guinée, Équateur) et de leurs potentialités.
\item Initiation à la lecture de cartes thématiques (choroplèthes, symboles proportionnels) et à la construction de croquis.
\item Vocabulaire de base : transhumance, nomadisme, sédentarisation, artisanat d'art, chaîne opératoire.
\end{itemize}
\end{prerequis}

\newpage

\coursec{L'élevage traditionnel : diversité des systèmes et organisation de l'espace}

Dans la cour de la chefferie de Bafut, un éleveur mbororo négocie le passage de son troupeau avec un cultivateur dont le champ de maïs borde le parcours. Plus au sud, à Kribi, un pêcheur artisanal rentre au port avec une pirogue motorisée, tandis que son père lui montre comment tresser les nasses en rotin comme au temps jadis. À Foumban, un artisan bronzier façonne une pipe royale sous le regard de touristes qui préfèrent parfois les copies en résine importées. Comment ces pratiques séculaires s'adaptent-elles aux pressions foncières, technologiques et marchandes d'aujourd'hui ?

\textbf{Problématique :} comment l'élevage traditionnel camerounais organise-t-il l'espace pour faire vivre des millions d'animaux, et quelles sont les formes, les richesses et les limites de ces systèmes pastoraux ?

Pour répondre à cette question, nous étudierons d'abord la typologie des systèmes d'élevage, puis le cheptel national et sa répartition, ensuite l'organisation spatiale du pastoralisme, enfin le rôle socio-économique de l'élevage et ses contraintes naturelles.

\courssub{Typologie des élevages traditionnels au Cameroun}

\textbf{Intuition.} Imaginez un troupeau de cent zébus. Où trouvera-t-il de l'herbe en pleine saison sèche, quand tout est brûlé dans les savanes du Nord ? L'éleveur n'a que deux solutions : se déplacer avec ses bêtes, ou rester sur place en nourrissant le troupeau autrement. De cette alternative naissent les grands types d'élevage traditionnel.

\begin{definition}[Systèmes d'élevage traditionnel]
On distingue trois grands systèmes pastoraux au Cameroun :
\begin{itemize}
\item le \cle{nomadisme} : déplacement quasi permanent du berger et de sa famille avec le troupeau, sans base fixe, au gré de l'eau et des pâturages ;
\item la \cle{transhumance} : migration saisonnière, cyclique et organisée du troupeau entre deux aires géographiques complémentaires (pâturage de saison sèche / pâturage de saison humide), avec une base fixe où reste une partie de la famille ;
\item l'\cle{élevage sédentaire de case} : les animaux (petits ruminants, quelques bovins) vivent attachés ou parqués près de la concession et sont nourris de résidus de récolte, de fanes et d'herbe coupée.
\end{itemize}
\end{definition}

\textbf{Le nomadisme pur} est pratiqué par certains éleveurs \textbf{Bororo} (Wodaabe) dans l'Extrême-Nord, autour du lac Tchad et de la plaine du Logone. Le campement, fait de cases légères démontables, suit le troupeau toute l'année. C'est un système en forte régression : sécheresses récurrentes, insécurité et pression agricole réduisent les espaces de mobilité.

\textbf{La transhumance nord-sud} est le système dominant. Les éleveurs \textbf{Mbororo} quittent les savanes soudaniennes du Nord et de l'Extrême-Nord à la fin des pluies (octobre-novembre), descendent vers les pâturages de saison sèche (plaines inondables, plateaux de l'Adamaoua), puis remontent au début des pluies (avril-mai) pour échapper aux eaux stagnantes et aux mouches. Certains flux atteignent l'Ouest et même le Nigeria voisin.

\textbf{L'élevage sédentaire de case} domine dans les zones de haute altitude et de forte densité humaine : \textbf{Adamaoua} (autour des concessions peulhs), \textbf{hauts plateaux de l'Ouest} (pays bamiléké, bamoun) et \textbf{Nord-Ouest}. On y trouve surtout des ovins, des caprins et quelques bovins taurins résistants, nourris à l'auge.

\begin{attention}[Nomadisme ou transhumance ?]
Erreur fréquente à l'examen : confondre \cle{nomadisme} (déplacement constant, pas de base fixe) et \cle{transhumance} (base fixe, aller-retour saisonnier entre deux aires complémentaires). Au Cameroun, c'est quasi-exclusivement de la \textbf{transhumance} ; le nomadisme pur ne concerne qu'une minorité de Bororo de l'Extrême-Nord. Ne jamais écrire « les éleveurs nomades du Cameroun » sans nuance.
\end{attention}

\courssub{Le cheptel national : composition et répartition régionale}

\textbf{Observation.} Sur un marché à bétail de Figuil ou de Ngaoundéré, on voit défiler des milliers de zébus à bosse. Mais dans les villages de l'Ouest, ce sont des chèvres et des moutons qui entourent les cases. Le cheptel camerounais est donc à la fois abondant et inégalement réparti.

\begin{exemplebox}[Un cheptel dominé par les bovins]
En 2022, le cheptel bovin camerounais est estimé à environ \textbf{6,5 millions de têtes}, auxquels s'ajoutent environ \textbf{6 millions d'ovins} et \textbf{8 millions de caprins} (Source : Annuaire statistique MINADER / FAO). Deux types de races bovines dominent :
\begin{itemize}
\item les \cle{zébus} (zébu Peul ou Goudali, zébu Djafoun), animaux à bosse des savanes, rustiques et résistants à la chaleur ;
\item les \cle{taurins} (Ndama, Muturu), sans bosse, tolérant mieux la trypanosomiase, présents dans les zones humides du sud.
\end{itemize}
L'Adamaoua et le Nord concentrent à eux deux environ \textbf{60 \%} des effectifs bovins nationaux. L'Adamaoua, avec ses vastes pâturages de plateau et sa faible densité humaine, est la première région d'élevage du pays.
\end{exemplebox}

\begin{center}
\begin{tabularx}{\linewidth}{|l|X|X|}
\hline
\rowcolor{popblueL} \textbf{Région} & \textbf{Système dominant} & \textbf{Importance du cheptel} \\
\hline
Adamaoua & Transhumance + sédentaire & Très forte (1\textsuperscript{re} région bovine) \\
\hline
Nord & Transhumance & Forte \\
\hline
Extrême-Nord & Nomadisme + transhumance & Moyenne (ovins, caprins, bovins) \\
\hline
Nord-Ouest / Ouest & Élevage de case & Moyenne (petits ruminants, taurins) \\
\hline
Sud forestier & Élevage de case & Faible (barrière de la mouche tsé-tsé) \\
\hline
\end{tabularx}
\end{center}

\begin{popfigure}
\begin{tikzpicture}[scale=0.95, every node/.style={font=\small}]
% Titre
\node[font=\bfseries] at (4,10.6) {Densité du cheptel bovin et axes de transhumance (Cameroun, schéma)} ;
% Contour très simplifié du Cameroun
\draw[thick, popdark] (1.5,9.5) -- (3.5,10.2) -- (5.5,9.8) -- (6.2,8.5) -- (6.8,7.2) -- (6.2,5.5) -- (6.5,3.5) -- (5.5,1.5) -- (3.5,0.8) -- (2.0,1.5) -- (1.2,3.0) -- (0.8,5.0) -- (1.0,7.0) -- (1.5,9.5) -- cycle;
% Zones colorées (densité)
\fill[poporange!55] (2.2,8.2) -- (5.2,8.6) -- (5.6,7.0) -- (2.0,6.8) -- cycle; % Nord / Extrême-Nord
\fill[poporange!80] (2.0,6.8) -- (5.6,7.0) -- (5.8,5.2) -- (2.4,4.9) -- cycle; % Adamaoua
\fill[popgreen!35] (1.2,4.9) -- (2.4,4.9) -- (2.2,2.6) -- (1.0,2.8) -- cycle; % Ouest/Nord-Ouest
\fill[popgreen!15] (2.4,4.9) -- (5.8,5.2) -- (5.5,1.6) -- (3.5,1.0) -- (2.2,2.6) -- cycle; % Sud forestier
% Villes
\fill[popink] (3.4,9.0) circle (1.5pt) node[above left]{Maroua};
\fill[popink] (3.6,7.6) circle (1.5pt) node[above left]{Garoua};
\fill[popink] (4.0,5.9) circle (1.5pt) node[right]{Ngaoundéré};
\fill[popink] (1.7,3.6) circle (1.5pt) node[left]{Bafoussam};
\fill[popink] (3.4,2.2) circle (1.5pt) node[below]{Yaoundé};
% Flèches de transhumance (épaisseur proportionnelle)
\draw[-{Stealth[length=3mm]}, line width=2.5pt, poppurple] (3.2,8.6) .. controls (3.0,7.6) .. (3.8,6.4);
\draw[-{Stealth[length=3mm]}, line width=2pt, poppurple] (4.6,8.2) .. controls (4.8,7.2) .. (4.4,6.2);
\draw[-{Stealth[length=2.5mm]}, line width=1.2pt, poppurple] (3.0,6.0) .. controls (2.2,5.2) .. (1.8,4.2);
% Nord
\draw[-{Stealth}, popdark] (0.6,10.2) -- (0.6,10.9) node[above]{N};
\end{tikzpicture}
\legende{1 : zone de très forte densité bovine (Adamaoua, orange foncé) ; 2 : zone de forte densité (Nord et Extrême-Nord, orange clair) ; 3 : élevage de case des hauts plateaux (vert) ; 4 : zone forestière à faible cheptel (vert pâle) ; flèches violettes : grands axes de transhumance nord-sud (épaisseur proportionnelle aux flux). Croquis schématique, échelle indicative. Source : données MINADER 2022.}
\end{popfigure}

\courssub{L'organisation spatiale du pastoralisme}

\textbf{Intuition.} Un troupeau qui se déplace a besoin de trois choses : un itinéraire sûr, de l'eau régulière, et de l'herbe à l'arrivée. Toute la géographie pastorale du Cameroun s'organise autour de ces trois éléments.

\begin{definition}[Couloir de transhumance]
Un \cle{couloir de transhumance} est un itinéraire balisé et juridiquement protégé que suivent les troupeaux entre leurs pâturages de saison humide et de saison sèche. Les grands couloirs camerounais sont le couloir du \textbf{Mayo-Rey} (Nord vers Adamaoua), les couloirs du \textbf{Logone} (Extrême-Nord) et ceux de la \textbf{Bénoué}. Leur largeur réglementaire est de 50 à 100 mètres.
\end{definition}

L'espace pastoral se structure ainsi en trois éléments complémentaires :
\begin{itemize}
\item les \textbf{pâturages de saison humide} : savanes soudaniennes du Nord et de l'Extrême-Nord, riches en herbes de mai à octobre, mais envahies par les eaux et les parasites en pleine saison des pluies ;
\item les \textbf{pâturages de saison sèche} : les \cle{yaérés}, plaines inondables du Logone et du Chari dont la décrue laisse des herbages verts en pleine saison sèche, véritables « greniers à herbe » convoités par des dizaines de milliers de têtes ;
\item les \textbf{points d'eau} : mares, forages pastoraux, retenues collinaires et rivières pérennes, qui rythment les étapes quotidiennes du troupeau (un bovin ne s'éloigne guère de plus de 15 km d'un point d'eau).
\end{itemize}

\begin{popfigure}
\begin{tikzpicture}[scale=1, every node/.style={font=\small}]
% Coupe topographique Nord-Sud
\draw[thick, popdark] (0,3.4) -- (2,3.6) -- (3.5,3.0) -- (5,4.2) -- (7,4.0) -- (8.5,3.2) -- (10,4.8) -- (11.5,4.6);
% Niveau de référence
\draw[dashed, popmuted] (0,1.2) -- (11.5,1.2) node[right, font=\footnotesize]{niveau de référence};
% Zones
\node[popblue] at (1,2.6) {Yaérés (Logone)};
\node[font=\footnotesize, popblue] at (1,2.2) {pâturage de saison sèche};
\node[poporange] at (4.2,3.5) {Savanes du Nord};
\node[font=\footnotesize, poporange] at (4.2,3.1) {pâturage de saison humide};
\node[popgreen] at (6.8,4.6) {Plateaux de l'Adamaoua};
\node[font=\footnotesize, popgreen] at (6.8,4.2) {(\SI{1000}{m} -- \SI{1200}{m})};
\node[poppurple] at (10.6,5.3) {Hauts plateaux de l'Ouest};
\node[font=\footnotesize, poppurple] at (10.6,4.9) {(\SI{1500}{m} et plus)};
% Flèches de migration
\draw[-{Stealth[length=3mm]}, line width=1.6pt, poppink] (3.2,4.4) .. controls (2.0,5.2) .. (1.2,4.0) node[midway, above, font=\footnotesize]{saison sèche};
\draw[-{Stealth[length=3mm]}, line width=1.6pt, popgold] (1.4,3.9) .. controls (2.4,2.6) .. (3.4,3.6) node[midway, below, font=\footnotesize]{saison des pluies};
\draw[-{Stealth[length=2.5mm]}, line width=1.2pt, poppink] (6.5,5.2) .. controls (8.0,6.0) .. (9.6,5.6) node[midway, above, font=\footnotesize]{flux vers l'Ouest};
% Axe
\draw[-{Stealth}] (0,0.6) -- (11.5,0.6) node[right]{Sud};
\node[left] at (0,0.6) {Nord};
\end{tikzpicture}
\legende{Coupe topographique schématique Nord-Sud : étagement des pâturages entre les yaérés du Logone (saison sèche), les savanes du Nord (saison humide), les plateaux de l'Adamaoua et les hauts plateaux de l'Ouest. Les flèches montrent les migrations saisonnières du troupeau (aller-retour) et les flux vers l'Ouest. Altitudes indicatives.}
\end{popfigure}

\begin{methode}[Construire un croquis de l'élevage au Cameroun]
Pour réussir un croquis géographique sur l'élevage traditionnel, procéder en étapes :
\begin{enumerate}
\item tracer le contour simplifié du pays, placer la flèche du nord et le titre ;
\item figurer les \textbf{surfaces} : hachurer ou colorer les grandes zones d'élevage (Adamaoua, Nord, yaérés) ;
\item figurer les \textbf{flèches} : dessiner les itinéraires transhumants nord-sud, avec une épaisseur proportionnelle à l'importance des flux ;
\item figurer les \textbf{points} : localiser par des symboles les points d'eau, les marchés à bétail (Ngaoundéré, Figuil) et les parcs de vaccination ;
\item rédiger une légende numérotée qui reprend chaque symbole dans l'ordre.
\end{enumerate}
\end{methode}

\courssub{Rôle socio-économique de l'élevage traditionnel}

L'élevage n'est pas seulement une production de viande : c'est un pilier de la vie rurale camerounaise.
\begin{itemize}
\item \textbf{Apport alimentaire} : il fournit l'essentiel des \cle{protéines animales} (viande, lait) des populations du septentrion et des villes ; le lait caillé et le beurre sont des produits de base chez les Peulhs.
\item \textbf{Force de travail} : les bœufs de trait servent au labour dans les savanes du Nord, et les ânes au transport.
\item \textbf{Fumure} : les contrats de fumure entre éleveurs et agriculteurs (le troupeau parqué sur un champ après récolte y dépose son fumier) lient élevage et agriculture.
\item \textbf{Capital d'épargne} : le troupeau est une « banque sur pattes » que l'on vend en cas de besoin (école, santé, cérémonies).
\item \textbf{Prestige social} : la richesse d'un homme se mesure au nombre de têtes ; le bétail sert au paiement de la \cle{dot} et aux sacrifices rituels.
\end{itemize}

\begin{exoresolu}[Calculer la part d'une région dans le cheptel national]
Énoncé : le cheptel bovin camerounais est estimé à 6,5 millions de têtes en 2022. L'Adamaoua et le Nord en concentrent environ 60 \%. Calculer l'effectif détenu par ces deux régions, puis l'effectif du reste du pays.
\tcblower
Solution détaillée :
\begin{align*}
\text{Effectif Adamaoua + Nord} &= 6{,}5 \times \frac{60}{100} = 3{,}9 \text{ millions de têtes} \\
\text{Effectif du reste du pays} &= 6{,}5 - 3{,}9 = 2{,}6 \text{ millions de têtes}
\end{align*}
Conclusion : près de 4 millions de bovins sur 6,5 millions vivent dans seulement deux régions, ce qui illustre la très forte concentration septentrionale de l'élevage camerounais et explique l'intensité des flux de transhumance vers le sud.
\end{exoresolu}

\courssub{Les contraintes naturelles de l'élevage traditionnel}

Le pastoralisme camerounais bute sur quatre contraintes majeures :
\begin{enumerate}
\item \textbf{La saisonnalité des ressources} : en saison sèche, les savanes du Nord sont brûlées ou épuisées et les mares tarissent ; d'où l'obligation de migrer vers les yaérés et les plateaux mieux arrosés.
\item \textbf{Les épizooties} : la \cle{péripneumonie contagieuse bovine} et la \cle{dermatose nodulaire contagieuse} déciment régulièrement les troupeaux ; les campagnes de vaccination restent insuffisantes.
\item \textbf{La mouche tsé-tsé} : la glossine, vecteur de la trypanosomiase animale (ou « nagana »), forme une véritable \textbf{barrière sanitaire au sud du 6\textsuperscript{e} parallèle} : elle interdit l'élevage des zébus dans la zone forestière humide, où seuls les taurins trypanotolérants survivent.
\item \textbf{Les feux de brousse} tardifs et l'avancée du front agricole, qui réduisent chaque année les pâturages disponibles.
\end{enumerate}

\begin{savaistu}[Des couloirs protégés par la loi]
Les couloirs de transhumance sont juridiquement protégés par la \textbf{loi n° 74/1 du 6 juillet 1974} fixant le régime foncier et domanial au Cameroun. Leur largeur réglementaire (50 à 100 mètres) est pourtant de plus en plus empiétée par les champs de maïs et de coton, ce qui est l'une des causes directes des conflits agriculteurs-éleveurs que nous étudierons dans la prochaine section.
\end{savaistu}

\begin{aretenir}[L'essentiel de la section]
\begin{itemize}
\item Trois systèmes : nomadisme (Bororo, marginal), transhumance (Mbororo, dominant), élevage de case (hautes terres).
\item Cheptel 2022 : environ 6,5 millions de bovins ; l'Adamaoua et le Nord en détiennent 60 \%.
\item L'espace pastoral s'organise autour des couloirs de transhumance, des points d'eau et des pâturages complémentaires (yaérés en saison sèche, savanes du Nord en saison humide).
\item L'élevage fournit protéines, travail, fumure, épargne et prestige social.
\item Ses limites : saisonnalité, épizooties, feux de brousse et barrière de la tsé-tsé au sud du 6\textsuperscript{e} parallèle.
\end{itemize}
\end{aretenir}

\begin{exercice}[Pour s'entraîner]
\begin{enumerate}
\item Définir la transhumance et la distinguer du nomadisme en deux phrases.
\item Sur un croquis du Cameroun, placer : les yaérés, l'Adamaoua, deux couloirs de transhumance et un marché à bétail.
\item Expliquer pourquoi l'élevage des zébus est quasi impossible dans la région du Sud.
\end{enumerate}
\end{exercice}

\begin{corrige}[Exercice]
\begin{enumerate}
\item La transhumance est une migration saisonnière, cyclique et organisée du troupeau entre deux aires complémentaires, avec une base fixe. Le nomadisme est un déplacement quasi permanent sans base fixe ; au Cameroun, il ne concerne qu'une minorité de Bororo.
\item Yaérés : plaine du Logone (Extrême-Nord) ; Adamaoua : centre du pays ; couloirs : Mayo-Rey et Logone, flèches nord-sud ; marché : Ngaoundéré ou Figuil.
\item Au sud du 6\textsuperscript{e} parallèle, la mouche tsé-tsé transmet la trypanosomiase animale (nagana), mortelle pour les zébus ; le milieu forestier humide entretient la glossine, et seuls les taurins trypanotolérants peuvent y être élevés.
\end{enumerate}
\end{corrige}

\coursec{Transformations de l'élevage traditionnel : modernisation, conflits et politiques publiques}

Nous avons vu dans la section précédente que l'élevage traditionnel camerounais repose sur des systèmes variés (transhumance, nomadisme, élevage extensif villageois) et sur une organisation de l'espace fondée sur la \cle{mobilité}. Or, depuis les indépendances et surtout depuis les années 1980, ce monde pastoral se transforme en profondeur : l'État pousse à la \cle{sédentarisation}, les terres se raréfient sous la pression démographique et agricole, les techniques se modernisent et de nouveaux acteurs urbains investissent dans le bétail. Cette section étudie ces mutations : la sédentarisation et les ranchs, les conflits agro-pastoraux qui en résultent, la modernisation technique et sanitaire, les politiques publiques d'appui, et enfin l'émergence des éleveurs « citadins ».

\courssub{La sédentarisation de l'élevage : ranchs, parcs améliorés et cultures fourragères}

\textbf{Intuition.} Un éleveur transhumant déplace son troupeau sur des centaines de kilomètres pour suivre l'herbe et l'eau. Mais que se passe-t-il si l'on clôture un vaste espace, qu'on y plante de l'herbe améliorée, qu'on y creuse des puits et qu'on y construit des bâtiments ? Le troupeau n'a plus besoin de bouger : c'est le principe du \cle{ranching}, cœur de la politique de sédentarisation menée au Cameroun depuis les années 1970.

\begin{definition}[Le ranching]
Le \cle{ranching} est un mode d'exploitation d'élevage \textbf{sédentaire}, pratiqué sur un espace \textbf{clôturé} de grande surface, géré de façon \textbf{extensive ou semi-intensive} : le troupeau y est contenu, l'alimentation y est améliorée (cultures fourragères, compléments) et les infrastructures (bâtiments, abreuvoirs, couloirs de vaccination) y sont fixes.
\end{definition}

\textbf{Pourquoi sédentariser ?} Les pouvoirs publics avancent plusieurs arguments : fixer les populations pastorales pour mieux les recenser, les scolariser et les soigner ; contrôler la santé animale (vaccinations, contrôles vétérinaires) ; intensifier la production de viande et de lait pour approvisionner les villes ; et réduire les conflits liés à la mobilité des troupeaux.

\textbf{Les réalisations concrètes.} Au Cameroun, la sédentarisation s'est traduite par trois types d'aménagements :

\begin{itemize}
\item La \textbf{création de ranchs} d'État ou parapublics, notamment dans la région de l'Adamaoua : le ranch de \cle{Ngaoundéré} et celui de \cle{Koundini} (situé dans le département de la Vina) sont les exemples de référence. On y trouve des bâtiments d'exploitation, des paddocks clôturés, des abreuvoirs permanents et des parcelles de fourrage cultivé.
\item Les \textbf{parcs à bétail améliorés} : espaces délimités où les éleveurs sédentarisés bénéficient d'infrastructures collectives (points d'eau, bains parasiticides, couloirs de contention).
\item Les \textbf{cultures fourragères} : introduction de graminées et de légumineuses fourragères améliorées comme le \cle{Brachiaria} (graminée très productive) et le \cle{Stylosanthes} (légumineuse riche en protéines), qui permettent de nourrir le bétail sur place, y compris en saison sèche.
\end{itemize}

\begin{propriete}[Bilan contrasté de la sédentarisation]
La sédentarisation \textbf{réduit la mobilité stratégique} qui permettait au pasteur de s'adapter aux aléas climatiques (suivre la pluie, fuir la sécheresse). En contrepartie, elle \textbf{permet} : le contrôle sanitaire du troupeau (vaccination, traitements), la scolarisation des enfants des éleveurs, l'accès aux services (santé, marchés) et la \textbf{fumure organique des champs} par les déjections du bétail parqué, ce qui favorise l'intégration agriculture-élevage.
\end{propriete}

\begin{attention}[Ne pas idéaliser le ranch]
La sédentarisation n'est pas toujours « choisie » : elle est souvent \textbf{forcée} par le manque de terres, la pression agricole ou les décisions administratives. De plus, un ranch exige des investissements (clôtures, forages, semences fourragères) inaccessibles au petit éleveur : les ranchs bénéficient surtout aux gros exploitants et aux investisseurs urbains. Enfin, sur un espace fixe, le surpâturage peut dégrader rapidement le sol autour des points d'eau.
\end{attention}

\courssub{Les conflits agro-pastoraux : causes, manifestations et mécanismes de résolution}

\textbf{Intuition et constat.} Lorsque les champs avancent sur les parcours et que les troupeaux, faute d'espace, traversent les cultures, l'affrontement devient presque inévitable. Le conflit agro-pastoral oppose un \textbf{agriculteur} (qui veut protéger ses champs) et un \textbf{éleveur} (qui veut faire paître et abreuver son troupeau) pour l'usage d'un même espace.

\textbf{Les causes profondes :}
\begin{itemize}
\item la \textbf{pression démographique} et l'extension des champs sur les anciens parcours et zones de pâturage ;
\item l'\textbf{empiètement des cultures sur les couloirs de transhumance} officiellement reconnus, qui bloque le passage des troupeaux ;
\item la \textbf{rareté de l'eau} en saison sèche, qui concentre hommes et animaux autour des mêmes points d'eau ;
\item le \textbf{non-respect du calendrier} : troupeaux lâchés avant la fin des récoltes.
\end{itemize}

\textbf{Les manifestations :} destruction des cultures par le bétail (champs de maïs, de coton ou de manioc piétinés et broutés), blessures et parfois morts d'hommes lors d'affrontements, abattage de bétail par représailles, incendies de campements ou de greniers.

\textbf{Les mécanismes de résolution :}
\begin{itemize}
\item la \textbf{médiation traditionnelle} : chefs de village, lamibé (chefs traditionnels du Nord) et chefs de troupeaux négocient des compensations (indemnisation en nature ou en argent) ;
\item les \textbf{commissions mixtes} agro-pastorales, instituées au niveau des sous-préfectures, qui réunissent agriculteurs, éleveurs, autorités administratives et services techniques pour juger les litiges et délimiter les espaces ;
\item le cadre juridique : la \textbf{loi n° 74/1 du 6 juillet 1974} fixant le régime foncier et domanial, complétée par les textes sur le pastoralisme, qui reconnaît l'existence de terres de parcours et encadre les déplacements du bétail.
\end{itemize}

\begin{methode}[Analyser un conflit agro-pastoral : la grille A.E.E.T.R.]
Pour étudier rigoureusement un conflit agro-pastoral (au brouillon comme à l'examen), appliquer la grille en cinq entrées :
\begin{enumerate}
\item \textbf{ACTEURS} : qui est en présence ? (agriculteurs, éleveurs peuls/mbororo, chefs traditionnels, autorités administratives, ONG, services du MINADER et du MINEPIA) ;
\item \textbf{ENJEUX} : que dispute-t-on ? (terre, eau, couloirs de passage, revenus, sécurité) ;
\item \textbf{ESPACE} : où se situe le conflit ? (parcours, champs, points d'eau, zones frontalières) ;
\item \textbf{TEMPS} : quand éclate-t-il ? (saison sèche = transhumance ; période des récoltes = dégâts de cultures) ;
\item \textbf{RÈGLEMENTATION} : quels outils de règlement ? (loi 74/1, commissions mixtes, médiation traditionnelle, arrêtés préfectoraux).
\end{enumerate}
\end{methode}

\begin{popfigure}
\begin{center}
\begin{tikzpicture}
\begin{axis}[
ybar, bar width=7pt,
width=\linewidth, height=7.5cm,
ylabel={Nombre de conflits déclarés},
symbolic x coords={Adamaoua,Nord,Extrême-Nord,Nord-Ouest},
xtick=data,
ymin=0, ymax=160,
legend style={at={(0.5,-0.18)},anchor=north,legend columns=2},
nodes near coords, every node near coord/.append style={font=\scriptsize},
]
\addplot[fill=popblue] coordinates {(Adamaoua,120) (Nord,95) (Extrême-Nord,70) (Nord-Ouest,60)};
\addplot[fill=poporange] coordinates {(Adamaoua,140) (Nord,110) (Extrême-Nord,85) (Nord-Ouest,55)};
\legend{2015, 2023}
\end{axis}
\end{tikzpicture}
\end{center}
\legende{Nombre de conflits agro-pastoraux déclarés par région en 2015 et 2023 (données indicatives reconstituées pour l'enseignement, d'après rapports MINADER et ONG locales). Lecture : la hausse est générale, surtout dans l'Adamaoua (+20 conflits, soit environ +17 \%), liée à la pression foncière et à l'arrivée de troupeaux en provenance des pays voisins ; seule la région du Nord-Ouest connaît une légère baisse grâce aux commissions mixtes.}
\end{popfigure}

\courssub{La modernisation technique : traction animale, amélioration génétique et santé vétérinaire}

\textbf{La traction animale.} Dans les régions du Nord (Grand-Nord, ceinture cotonnière), les bœufs de labour attelés à des \textbf{charrues} et des \textbf{charrettes} remplacent la houe : on laboure plus vite, on cultive plus de surface, on transporte récoltes et coton. C'est un maillon essentiel de l'\cle{intégration agriculture-élevage} : l'animal travaille le champ et le fume, le champ fournit résidus de récolte et fourrages.

\begin{attention}[Ne pas généraliser la traction animale]
La traction animale se développe au \textbf{Nord} (plaine soudanienne, ceinture cotonnière, climat sec sans maladie du sommeil) mais reste \textbf{marginale au Sud} : la forêt dense, le relief accidenté et surtout la \textbf{mouche tsé-tsé} (vecteur de la trypanosomiase animale) y empêchent le maintien de gros bovins de trait. À l'examen, ne jamais écrire que la traction animale est répandue « dans tout le Cameroun ».
\end{attention}

\textbf{L'amélioration génétique.} Pour augmenter la production de lait et de viande, les services vétérinaires pratiquent :
\begin{itemize}
\item l'\textbf{insémination artificielle} (IA), qui introduit des gènes de races performantes sans déplacer les reproducteurs ;
\item le \textbf{croisement} entre le \cle{zébu} local (rustique, résistant à la chaleur) et des \textbf{taurins laitiers} importés, pour obtenir des animaux métis plus productifs.
\end{itemize}

\begin{savaistu}[Zébus améliorés : Goudali et Bokolodji]
Le zébu \cle{Goudali} (race laitière de l'Adamaoua) et le \cle{Bokolodji} (race à viande) améliorent nettement la productivité des troupeaux camerounais. Mais ces animaux performants sont aussi plus \textbf{exigeants} : ils réclament des compléments alimentaires, des soins vétérinaires réguliers et un abreuvement de qualité, autant d'intrants souvent \textbf{inaccessibles aux éleveurs pauvres}. Le progrès génétique peut donc creuser les écarts entre gros et petits éleveurs.
\end{savaistu}

\textbf{La santé vétérinaire.} L'État et ses partenaires organisent des \textbf{campagnes de vaccination} (contre la péripneumonie contagieuse bovine, la pasteurellose, le charbon symptomatique) et installent des bains parasiticides. Le contrôle sanitaire, impossible sur un troupeau nomade dispersé, devient réalisable dès lors que les animaux sont sédentarisés ou rassemblés aux points d'eau aménagés.

\courssub{Les politiques publiques d'appui à l'élevage}

L'État camerounais encadre la transformation de l'élevage par plusieurs instruments :

\begin{itemize}
\item le \textbf{PNDP} (Programme National de Développement Participatif), qui finance les infrastructures rurales de base (points d'eau, pistes, marchés à bétail) décidées avec les communes ;
\item le \textbf{PEAJ} (Programme d'Appui à l'Entrepreneuriat des Adolescents et Jeunes), qui aide les jeunes à s'installer dans les métiers de l'élevage ;
\item le \textbf{PRODEL} (Projet de Développement de l'Élevage), projet phare du MINEPIA, qui appuie la modernisation des élevages familiaux et pastoraux.
\end{itemize}

\begin{exemplebox}[Le PRODEL en chiffres]
Entre 2018 et 2023, le PRODEL a financé la réhabilitation de \textbf{plus de 500 km de couloirs de transhumance} et la modernisation de \textbf{200 points d'eau} pastoraux, tout en appuyant la formation des éleveurs et l'équipement des groupements. Objectif : sécuriser les déplacements du bétail et réduire les conflits en matérialisant les parcours.
\end{exemplebox}

Parallèlement, l'État encourage l'organisation des éleveurs en \textbf{GIC} (Groupes d'Initiative Commune) et en \textbf{coopératives} : regroupés, ils accèdent plus facilement au crédit, aux intrants, aux formations et aux marchés, et pèsent davantage dans les commissions mixtes de règlement des conflits.

\courssub{Les nouveaux acteurs : les éleveurs « citadins »}

Depuis les années 1990, une catégorie nouvelle s'impose : les \cle{éleveurs citadins}. Ce sont des \textbf{fonctionnaires, commerçants, cadres retraités} qui investissent leur épargne dans l'élevage, sans vivre auprès des animaux. Ils achètent de jeunes bovins maigres à la fin de la saison sèche (prix bas), les font \textbf{engraisser} quelques mois dans des parcs d'engraissement en \textbf{périphérie des grandes villes} (Yaoundé, Douala, Bafoussam), puis les revendent au prix fort aux abattoirs urbains, notamment avant les fêtes.

Ce phénomène a trois conséquences majeures : il \textbf{approvisionne} les villes en viande ; il \textbf{intensifie} la production (embouche, compléments alimentaires) ; mais il fait aussi \textbf{monter le prix des terres} en périphérie urbaine et accentue la concurrence foncière avec les petits agriculteurs et les pasteurs traditionnels.

\begin{experience}[TP : comparer un campement transhumant et un ranch moderne]
\textbf{Matériel :} deux photographies (ou deux descriptions) : un campement transhumant mbororo de saison sèche et le ranch de Koundini (Adamaoua).
\textbf{Protocole :} observer chaque document et remplir un tableau comparatif à quatre rubriques : habitat, gestion de l'eau, alimentation du troupeau, mobilité.
\textbf{Observations attendues :}
\begin{itemize}
\item \emph{Campement transhumant} : cases temporaires en matériaux végétaux, abreuvoirs naturels (mares, cours d'eau), alimentation exclusivement au pâturage, mobilité permanente du troupeau et de la famille.
\item \emph{Ranch de Koundini} : bâtiments fixes en dur (logements, hangars, couloirs de vaccination), abreuvoirs permanents alimentés par forages, parcelles cultivées de \emph{Brachiaria} pour le fourrage, troupeau contenu dans des paddocks clôturés.
\end{itemize}
\textbf{Interprétation :} le ranch matérialise la sédentarisation (fixité, intensification, contrôle sanitaire possible), tandis que le campement exprime la logique de mobilité (adaptation au climat, faible investissement, mais vulnérabilité aux conflits fonciers).
\end{experience}

\begin{exoresolu}[Exercice résolu : analyser un conflit agro-pastoral]
\textbf{Énoncé.} Dans un arrondissement de l'Adamaoua, en janvier, un troupeau de 300 zébus en transhumance détruit 12 ha de champs de maïs non encore récoltés. Les agriculteurs saisissent le sous-préfet ; les éleveurs affirment suivre un couloir officiel rétréci par de nouvelles plantations. En appliquant la grille A.E.E.T.R., analyse ce conflit et propose deux solutions.
\tcblower
\textbf{Solution détaillée.}
\begin{enumerate}
\item \textbf{Acteurs} : les agriculteurs villageois, les éleveurs transhumants (et leurs chefs de troupeau), le sous-préfet (autorité administrative), les services du MINADER et du MINEPIA, les chefs traditionnels.
\item \textbf{Enjeux} : la terre (parcours contre champs), la sécurité alimentaire des agriculteurs (récolte perdue), la survie du troupeau (accès à l'eau et au pâturage), la paix sociale.
\item \textbf{Espace} : un couloir de transhumance officiel empiété par les cultures, au contact du village et des champs.
\item \textbf{Temps} : janvier, pleine saison sèche = période de transhumance descendante ; les récoltes de maïs tardives ne sont pas rentrées, d'où la confrontation.
\item \textbf{Règlementation} : la loi n° 74/1 du 6 juillet 1974 (régime foncier, terres de parcours), les arrêtés délimitant les couloirs, la commission mixte agro-pastorale de l'arrondissement, la médiation traditionnelle.
\end{enumerate}
\textbf{Solutions proposées} : (1) réunir la commission mixte pour constater les dégâts, fixer l'indemnisation des agriculteurs et re-matérialiser le couloir par des bornes ; (2) harmoniser le calendrier : date officielle d'ouverture des parcours après la fin des récoltes, et création d'un point d'eau aménagé hors des champs pour détourner le troupeau.
\end{exoresolu}

\begin{exercice}[À faire]
1. Définis le ranching et cite deux ranchs camerounais de référence.
2. Explique pourquoi la traction animale est développée au Nord et marginale au Sud du Cameroun.
3. À l'aide du graphique de la leçon, calcule le taux d'évolution du nombre de conflits déclarés dans l'Adamaoua entre 2015 et 2023.
4. Rédige un paragraphe argumenté : « Les éleveurs citadins : une chance ou une menace pour l'élevage traditionnel ? »
\end{exercice}

\begin{corrige}[Exercice]
1. Le ranching est une exploitation d'élevage sédentaire, clôturée, gérée de façon extensive ou semi-intensive sur de grandes surfaces. Références : ranchs de Ngaoundéré et de Koundini (Adamaoua).
2. Au Nord : plaine soudanienne, ceinture cotonnière, climat sec sans tsé-tsé, donc bovins de trait possibles et utiles (charrues, charrettes). Au Sud : forêt dense, relief accidenté et mouche tsé-tsé (trypanosomiase) qui empêchent le maintien des bœufs de labour.
3. Taux d'évolution $= \dfrac{140 - 120}{120} \times 100 \approx 16,7\,\%$ : le nombre de conflits déclarés a augmenté d'environ 17 \% entre 2015 et 2023 dans l'Adamaoua.
4. Attendu : une chance (investissements, embouche, approvisionnement des villes en viande, modernisation) mais aussi une menace (hausse du prix des terres, concurrence foncière avec les petits éleveurs, spéculation). La conclusion doit être nuancée et justifiée.
\end{corrige}

\begin{aretenir}[L'essentiel de la section]
\begin{itemize}
\item La \textbf{sédentarisation} (ranchs de Ngaoundéré et Koundini, parcs améliorés, fourrages comme le Brachiaria et le Stylosanthes) fixe les troupeaux : gains sanitaires, scolaires et agronomiques (fumure), mais perte de la mobilité d'adaptation au climat.
\item Les \textbf{conflits agro-pastoraux} naissent de la concurrence pour la terre et l'eau ; ils se règlent par la médiation traditionnelle, les commissions mixtes et le cadre de la loi 74/1.
\item La \textbf{modernisation} combine traction animale (Nord seulement), amélioration génétique (IA, croisements zébu $\times$ taurin) et campagnes de vaccination.
\item Les \textbf{politiques publiques} (PNDP, PEAJ, PRODEL : plus de 500 km de couloirs et 200 points d'eau) et les groupements (GIC, coopératives) structurent la filière.
\item Les \textbf{éleveurs citadins} développent l'engraissement périurbain autour de Yaoundé, Douala et Bafoussam.
\end{itemize}
\end{aretenir}

\coursec{La pêche traditionnelle : techniques, espaces et mutations}

Après avoir analysé les mutations de l'élevage traditionnel, confronté à la sédentarisation, aux conflits fonciers et à l'introduction d'espèces améliorées, nous nous tournons maintenant vers le second pilier de l'agropastoralisme traditionnel : la pêche. Si l'élevage est l'affaire des savanes et des hautes terres, la pêche structure la vie des populations riveraines des grands bassins hydrographiques et du littoral atlantique. Elle partage avec l'élevage une même logique : une exploitation extensive des ressources naturelles, reposant sur un savoir-faire ancestral, aujourd'hui bousculée par la motorisation, la demande urbaine croissante et la pression industrielle.

\courssub{Typologie des pêches : une géographie de l'eau et des hommes}

Le Cameroun offre une mosaïque d'écosystèmes aquatiques qui conditionne une diversité de pratiques halieutiques. On distingue quatre grands espaces de pêche artisanale, chacun avec ses dynamiques propres.

\begin{definition}[Pêche artisanale (définition administrative et technique)]
La \cle{pêche artisanale} désigne l'activité de capture réalisée à pied ou à bord d'embarcations de faible tonnage (généralement < 15 m), propulsées à la voile, à la pagaie ou par moteur hors-bord de faible puissance (< 25-40 CV). Elle mobilise une main-d'œuvre majoritairement familiale, un capital limité, des engins sélectifs (mais parfois non réglementaires) et assure une commercialisation locale ou régionale (circuits courts). Elle se distingue de la pêche industrielle (chalutiers, senneurs > 25 m, équipages salariés, exportation) et de la pêche purement \emph{traditionnelle} (engins et pirogues strictement traditionnels, non motorisés).
\end{definition}

\begin{popfigure}
\begin{tikzpicture}[scale=0.85, every node/.style={transform shape}]
% Fond marin / eau
\fill[popblue!10] (-7,-5) rectangle (7,5);
% Contour très simplifié du Cameroun (polygone approximatif)
\draw[popdark, thick, fill=popcream] 
(-6.5, -4.5) -- (-6, -3) -- (-4, -2.5) -- (-2, -1.5) -- (-0.5, -0.5) -- (1, 0.5) -- (2.5, 1) -- (4, 0.5) -- (5, -0.5) -- (5.5, -2) -- (4.5, -3.5) -- (3, -4.5) -- (1, -4.5) -- (-1, -4) -- (-3.5, -4.2) -- (-5.5, -4.5) -- cycle;
% Lac Tchad (extrême nord)
\fill[popblue!40] (-5.5, 3.5) ellipse ({1.2} and {0.6});
\node[popdark, font=\tiny\bfseries] at (-5.5, 3.5) {Lac Tchad};
\draw[poporange, thick] (-5.5, 3.5) -- (-5.5, 2.8) node[below, font=\tiny, poporange] {Fluviale/Lacustre};
% Bassin Bénoué (Nord)
\draw[popblue, very thick, decorate, decoration={snake, amplitude=1pt, segment length=10pt}] (-3, 2.5) .. controls (-1, 2) and (1, 1.5) .. (3, 1);
\node[popblue, font=\tiny\bfseries, rotate=-15] at (0, 2.2) {Bénoué};
% Barrage Lagdo
\fill[poppurple] (1.5, 1.7) circle (3pt);
\node[above, font=\tiny, poppurple] at (1.5, 1.7) {Lagdo};
% Fleuves du Sud (Sanaga, Nyong, Wouri)
\draw[popblue, very thick, decorate, decoration={snake, amplitude=1pt, segment length=8pt}] (-1, 0.5) .. controls (0, 0) and (1, -0.5) .. (2.5, -1.5);
\node[popblue, font=\tiny\bfseries, rotate=-20] at (0.5, 0.2) {Sanaga};
\draw[popblue, very thick, decorate, decoration={snake, amplitude=1pt, segment length=8pt}] (-2.5, 0) .. controls (-1.5, -0.5) and (-0.5, -1) .. (0.5, -2);
\node[popblue, font=\tiny\bfseries, rotate=-25] at (-1.5, -0.5) {Nyong};
% Wouri / Estuaire
\fill[popblue!30] (2.5, -1.5) ellipse ({0.8} and {0.4});
\node[font=\tiny\bfseries, popdark] at (2.5, -1.5) {Wouri};
% Lagunes / Côte (Kribi, Limbé, Douala)
% Douala - Bonabéri
\fill[popgreen!50] (3.2, -1.8) circle (0.25);
\node[font=\tiny\bfseries, popdark] at (3.2, -1.8) {Douala};
\node[font=\tiny, popdark] at (3.2, -2.2) {Bonabéri};
% Kribi
\fill[popgreen!50] (4.5, -3) circle (0.2);
\node[font=\tiny\bfseries, popdark] at (4.5, -3) {Kribi};
% Limbé
\fill[popgreen!50] (1.5, -0.5) circle (0.2);
\node[font=\tiny\bfseries, popdark] at (1.5, -0.5) {Limbé};
% Symboles proportionnels tonnages (cercles rouges)
% Douala : gros tonnage
\draw[popred, thick] (3.2, -1.8) circle (0.6);
\node[font=\tiny, popred] at (3.2, -1.8) {\textasciitilde 40\%};
% Kribi
\draw[popred, thick] (4.5, -3) circle (0.4);
\node[font=\tiny, popred] at (4.5, -3) {\textasciitilde 25\%};
% Limbé
\draw[popred, thick] (1.5, -0.5) circle (0.35);
\node[font=\tiny, popred] at (1.5, -0.5) {\textasciitilde 15\%};
% Lac Tchad
\draw[popred, thick] (-5.5, 3.5) circle (0.35);
\node[font=\tiny, popred] at (-5.5, 3.5) {\textasciitilde 10\%};
% Intérieur (Lagdo, fluvial)
\draw[popred, thick] (1.5, 1.7) circle (0.25);
\node[font=\tiny, popred] at (1.5, 1.7) {\textasciitilde 10\%};
% Légende intégrée
\begin{scope}[xshift=7.5cm, yshift=-2cm]
\draw[popdark] (0,1.5) rectangle (4.5, -3.5);
\node[font=\small\bfseries, popdark] at (2.25, 1.2) {LÉGENDE};
% Zones
\fill[popblue!40] (0.3, 0.8) rectangle (0.7, 1.1); \node[anchor=west, font=\tiny] at (0.9, 0.95) {Eaux continentales (Lacs, Fleuves)};
\fill[popblue!30] (0.3, 0.4) rectangle (0.7, 0.7); \node[anchor=west, font=\tiny] at (0.9, 0.55) {Estuaires / Lagunes (Wouri)};
\fill[popblue!10] (0.3, 0.0) rectangle (0.7, 0.3); \node[anchor=west, font=\tiny] at (0.9, 0.15) {Domaine maritime (3 milles)};
% Ports
\fill[popgreen!50] (0.3, -0.4) circle (0.15); \node[anchor=west, font=\tiny] at (0.9, -0.4) {Ports de débarquement / Criées};
% Tonnages
\draw[popred, thick] (0.5, -0.8) circle (0.15); \node[anchor=west, font=\tiny] at (0.9, -0.8) {Tonnages débarqués (proportionnel)};
% Barrages
\fill[poppurple] (0.5, -1.2) circle (0.1); \node[anchor=west, font=\tiny] at (0.9, -1.2) {Barrages hydroélectriques (Lagdo, Mapé, Mbakaou)};
% Conflits
\draw[poporange, thick, decorate, decoration={zigzag, amplitude=1.5pt}] (0.4, -1.6) -- (0.8, -1.6); \node[anchor=west, font=\tiny] at (0.9, -1.6) {Zones de conflit (chalutiers)};
% Flèche Nord
\draw[->, popdark, thick] (3.5, -2.5) -- (3.5, -1.8);
\node[font=\tiny, popdark] at (3.5, -1.5) {N};
\end{scope}
\end{tikzpicture}
\legende{Carte schématique des zones de pêche artisanale au Cameroun. Les cercles rouges sont proportionnels aux tonnages annuels débarqués par port principal (estimations MINEPIA). Les zones hachurées bleues distinguent les domaines : maritime (côte), lagunaire/estuarien (Wouri), fluvial (Bénoué, Sanaga, Nyong, Logone/Chari) et lacustre (Tchad, Maga, barrages).}
\end{popfigure}

\textbf{1. La pêche maritime artisanale (Côte atlantique, \textasciitilde 400 km).}\par
Elle se concentre sur trois pôles majeurs : \textbf{Kribi} (Sud, pêche démersale et pélagique, crevettes), \textbf{Limbé} (Sud-Ouest, pêche pélagique intense, sardinelles, bonga) et \textbf{Douala-Bonabéri} (Estuaire Wouri, premier port de débarquement en volume, crevettes, poissons-chats, ethmalose). Les pêcheurs opèrent dans la bande des \textbf{3 milles nautiques} (zone réservée par la loi), mais y côtoient illégalement les chalutiers industriels.

\textbf{2. La pêche lagunaire et estuarienne (Wouri, Sanaga, Nyong).}\par
L'estuaire du Wouri (Douala) est une zone de nurserie critique. La pêche y est pratiquée en pirogues légères, ciblant l'\emph{Ethmalosa fimbriata} (bonga), les crevettes juvéniles (\emph{Penaeus notialis}) et le tilapia. La pollution industrielle (huiles, métaux lourds, plastiques) y est critique.

\textbf{3. La pêche fluviale (Bénoué, Logone, Chari, Sanaga, Nyong).}\par
Sur le fleuve Bénoué (navigable 6 mois/an de Garoua au Nigeria), la pêche est saisonnière, liée à la crue (\emph{flood pulse}). Les barrages de \textbf{Lagdo} (Bénoué), \textbf{Mapé} (Mbam), \textbf{Mbakaou} (Noun) et \textbf{Maga} (Logone) ont créé de vastes lacs de retenue (\textbf{pêche lacustre de barrage}) qui concentrent désormais l'essentiel de la production nordique (Tilapia, Capitaine \emph{Lates niloticus}, Silure \emph{Clarias gariepinus}).

\textbf{4. La pêche lacustre naturelle : Lac Tchad et Lac Maga.}\par
Le Lac Tchad (partie camerounaise : archipels de Blarigui, Darak) voit sa production fluctuer avec le niveau du lac (variabilité climatique). Le Lac Maga (retenue sur le Logone) est un réservoir artificiel productif pour le tilapia et le capitaine.

\begin{exemplebox}[Répartition spatiale de la production halieutique artisanale (Estimations MINEPIA/FAO récentes)]
\begin{center}
\begin{tabularx}{\linewidth}{|l|X|c|X|}\hline
\rowcolor{popblueL}
\textbf{Espace} & \textbf{Ports / Zones clés} & \textbf{Part estimée} & \textbf{Espèces dominantes} \\ \hline
Maritime & Kribi, Limbé, Douala-Bonabéri, Idenau & \textasciitilde 45-50\% & Bonga, Sardinelles, Crevettes, Carangues, Barracudas \\ \hline
Estuarien / Lagunaire & Douala (Wouri), Mouanko, Kribi (embouchures) & \textasciitilde 15\% & Ethmalosa (Bonga), Crevettes juvéniles, Tilapia, Poisson-chat \\ \hline
Fluvial (Bénoué) & Garoua, Lagdo, Rey Bouba, Garoua-Boulaï & \textasciitilde 20\% & Capitaine (Nilotique), Tilapia, Silure, Distichodus \\ \hline
Lacustre (Tchad, Maga, Barrages) & Blarigui, Darak, Maga, Lagdo, Mapé, Mbakaou & \textasciitilde 15-20\% & Tilapia, Capitaine, Silure, Carpes \\ \hline
\end{tabularx}
\end{center}
\legende{La pêche maritime domine les volumes, mais la pêche continentale (fluviale + lacustre) assure la sécurité alimentaire protéique du Nord et de l'intérieur du pays.}
\end{exemplebox}

\courssub{Techniques et engins traditionnels : l'ingéniosité adaptée au milieu}

La pêche traditionnelle repose sur une connaissance intime du comportement des espèces (migrations, reproduction, nycthéméralisme) et de l'hydrodynamique (courants, marées, stratification thermique).

\begin{methode}[Analyse technique d'un engin de pêche traditionnel]
Pour comprendre l'efficacité et l'impact d'un engin, suivez ces 4 étapes :
\begin{enumerate}
\item \textbf{Identifier le principe de capture :} Passif (le poisson se prend seul : nasse, filet maillant) vs Actif (le pêcheur poursuit : harpon, ligne, barrageage).
\item \textbf{Analyser la sélectivité (taille / espèce) :} Déterminée par la \cle{maille} (ouverture du filet) ou la taille de l'hameçon. Une maille carrée de 25 mm capture des juvéniles ; 45-50 mm cible les adultes.
\item \textbf{Évaluer l'empreinte spatiale et temporelle :} Zone de pêche (fond, surface, colonne d'eau), durée de calée (filet), effort humain (nombre de pêcheurs/pirogues).
\item \textbf{Mesurer le capital technique :} Matériaux (fibres végétales \emph{vs} nylon), motorisation, coût de maintenance, durée de vie.
\end{enumerate}
\end{methode}

\begin{center}
\begin{tabularx}{\linewidth}{|l|X|X|X|}\hline
\rowcolor{popblueL}
\textbf{Engin / Technique} & \textbf{Principe \& Milieu} & \textbf{Matériaux traditionnels} & \textbf{Espèces cibles / Sélectivité} \\ \hline
\textbf{Nasse (Trappe)} & Passif, fond (fleuves, lacs, barrages). Appâtée ou non. & Rotin, raphia, lianes (non retournables). & Silure, Tilapia, Capitaine (gros spécimens). Très sélective par l'entrée. \\ \hline
\textbf{Filet maillant (Gillnet)} & Passif, dérivant ou calé. Le poisson s'engouffre par les ouïes. & Fibres végétales (écorce de \emph{Raphia}, \emph{Hibiscus}) $\rightarrow$ Nylon monofilament. & \textbf{Tout pélagique/démersal.} Sélectivité \textbf{strictement liée à la maille}. \\ \hline
\textbf{Senne tournante / Senne de plage} & Actif, encerclement. Deux bras rabattus vers la plage ou la pirogue. & Cordages végétaux, flotteurs bois (bambou), plombs pierre. & Bancs pélagiques (Bonga, Sardinelles). \textbf{Peu sélective (tout le banc)}. \\ \hline
\textbf{Ligne (Palangre, main)} & Actif/Passif. Hameçons appâtés. Fond ou surface. & Fil végétal, hameçons forgés (recyclés). & Capitaine, Silure, Carangues (gros). Sélective par taille d'hameçon. \\ \hline
\textbf{Harpon / Flèche} & Actif, visée directe. Eau claire, faible profondeur (lac, lagune). & Bois dur (Iroko, Dibétou), pointe métal. & Grosses proies (Capitaine, Silure, Tortues - interdit). \\ \hline
\textbf{Barrageage (Akadja / Brush park)} & Passif, attraction. Fascines plantées créant un habitat artificiel. & Branches, bambous, herbes (renouvelé 2-3 fois/an). & Tilapia, Silure, Crevettes (nurserie). \textbf{Impact écosystémique fort}. \\ \hline
\textbf{Pêche à main / Épuisette} & Actif, marée basse, mangroves, rizicoles. & Mains, paniers tressés. & Crabes, crevettes, alevins (juvéniles). \\ \hline
\end{tabularx}
\end{center}

\textbf{La pirogue traditionnelle : l'outil universel.}\par
Creusée dans un tronc unique d'\textbf{Iroko} (\emph{Milicia excelsa}) ou de \textbf{Dibétou} (\emph{Lovoa trichilioides}), essences imputrescibles et légères. Longueur 6-12 m, largeur 0,8-1,2 m. Propulsion : \textbf{pagaie} (fleuves, lagunes), \textbf{voile carrée} en natte de raphia (lac Tchad, maritime côtier), \textbf{perche} (fond sableux, faible profondeur). Capacité : 2 à 5 pêcheurs + 200-500 kg de poisson. Durée de vie : 15-25 ans avec entretien (huilage, feuillardage des fissures).

\begin{popfigure}
\begin{tikzpicture}[scale=1.1, every node/.style={transform shape}]
% Coupe transversale : Filet maillant dérivant (Gillnet)
\begin{scope}[xshift=-5cm]
\draw[popblue, very thick] (-2, -1.5) -- (2, -1.5) node[midway, below, font=\tiny] {Fond marin / vase};
% Eau
\fill[popblue!10, opacity=0.5] (-2.5, -1.5) rectangle (2.5, 2);
% Flotteurs (haut)
\foreach \x in {-1.8, -1.2, -0.6, 0, 0.6, 1.2, 1.8} {
    \fill[poporange!80] (\x, 1.8) circle (0.1);
    \draw[popdark, thin] (\x, 1.8) -- (\x, 1.5);
}
\node[poporange, font=\tiny\bfseries, anchor=south] at (0, 1.9) {Ligne de flottaison (Flotteurs liège/plastique)};
% Plombs (bas)
\foreach \x in {-1.8, -1.2, -0.6, 0, 0.6, 1.2, 1.8} {
    \fill[popdark] (\x, -1.2) circle (0.07);
    \draw[popdark, thin] (\x, -1.2) -- (\x, -1.5);
}
\node[popdark, font=\tiny\bfseries, anchor=north] at (0, -1.3) {Ligne de plombage (Plombs cylindriques)};
% Filet (mailles losanges)
\draw[popred, very thin] (-2, 1.5) -- (-2, -1.2);
\draw[popred, very thin] (2, 1.5) -- (2, -1.2);
% Mailles
\foreach \y in {1.3, 0.9, 0.5, 0.1, -0.3, -0.7, -1.1} {
    \foreach \x in {-1.8, -1.0, -0.2, 0.6, 1.4} {
        \draw[popred, thin] (\x, \y) -- (\x+0.4, \y-0.2) -- (\x+0.8, \y) -- (\x+0.4, \y+0.2) -- cycle;
    }
}
% Poisson piégé (par les ouïes)
\draw[popgreen!60!black, thick] (0.2, 0.1) -- (0.8, 0.3) -- (0.6, 0.5) -- (0.2, 0.4) -- (0.0, 0.2) -- cycle; % Corps
\draw[popgreen!60!black, thick] (0.6, 0.4) -- (1.0, 0.5); % Queue
\draw[popred, thick] (0.3, 0.25) circle (0.08); % Œil
\node[popred, font=\tiny, anchor=west] at (1.1, 0.3) {Poisson maillé (ouïes bloquées)};
% Maille réglementaire vs pratiquée
\node[font=\small\bfseries, popdark] at (0, 2.3) {FILET MAILLANT DÉRIVANT (Gillnet)};
\node[font=\tiny, popred] at (0, -1.8) {Maille carrée réglementaire : 45-50 mm (adultes)};
\node[font=\tiny, poporange] at (0, -2.0) {Maille souvent pratiquée : 25-30 mm (juvéniles)};
\end{scope}
% Coupe transversale : Senne tournante (Purse Seine) - Simplifié
\begin{scope}[xshift=5cm]
\draw[popblue, very thick] (-2, -1.5) -- (2, -1.5) node[midway, below, font=\tiny] {Fond};
\fill[popblue!10, opacity=0.5] (-2.5, -1.5) rectangle (2.5, 2);
% Cercle de la senne
\draw[poppurple, very thick] (0,0) circle (1.5);
% Flotteurs haut
\foreach \a in {0, 45, 90, 135, 180, 225, 270, 315} {
    \fill[poporange!80] ({1.5*cos(\a)}, {1.5*sin(\a)}) circle (0.08);
}
% Plombs bas (anneau de plombage)
\foreach \a in {0, 45, 90, 135, 180, 225, 270, 315} {
    \fill[popdark] ({1.3*cos(\a)}, {1.3*sin(\a)-0.2}) circle (0.06);
}
% Coulisse (bourrelet) bas
\draw[poppurple, thick, decorate, decoration={snake, amplitude=0.5pt}] (-1.3, -1.4) -- (1.3, -1.4);
\node[poppurple, font=\tiny\bfseries] at (0, -1.6) {Coulisse / Bourrelet (fermeture fond)};
% Banc de poissons à l'intérieur
\foreach \i in {1..15} {
    \pgfmathsetmacro{\rx}{rand*1.0}
    \pgfmathsetmacro{\ry}{rand*1.0}
    \draw[popgreen!60!black, thin] (\rx, \ry) -- ++(0.15, 0.05) -- ++(-0.1, 0.1) -- ++(-0.15, -0.05) -- cycle;
}
\node[popgreen!60!black, font=\tiny\bfseries] at (0, 0) {BANC ENCEINTÉ};
% Pirogues
\draw[poporange!70!popdark, thick] (-2.2, 1.8) .. controls (-2.5, 1.5) .. (-2.2, 1.2);
\node[poporange!70!popdark, font=\tiny, anchor=east] at (-2.2, 1.8) {Pirogue 1 (filet)};
\draw[poporange!70!popdark, thick] (2.2, 1.8) .. controls (2.5, 1.5) .. (2.2, 1.2);
\node[poporange!70!popdark, font=\tiny, anchor=west] at (2.2, 1.8) {Pirogue 2 (fermeture)};
\node[font=\small\bfseries, popdark] at (0, 2.3) {SENNE TOURNANTE (Purse Seine)};
\node[font=\tiny, poppurple] at (0, -1.9) {Capture de masse : tout le banc (juvéniles + adultes)};
\node[font=\tiny, poppurple] at (0, -2.1) {Interdite en zone artisanale (souvent tolérée)};
\end{scope}
% Zoom Maille : Réglementaire vs Pratiquée (Impact juvéniles)
\begin{scope}[yshift=-5.5cm, xshift=-2cm]
\draw[popdark, thick] (-3, 0) rectangle (3, 1.5);
\node[font=\small\bfseries, popdark] at (0, 1.3) {COMPARAISON DES MAILLES (Vue de face - Échelle approx.)};
% Maille Réglementaire (45mm)
\begin{scope}[xshift=-3.5cm]
\node[font=\tiny\bfseries, popgreen] at (0, 1.0) {Réglementaire (45-50 mm)};
\draw[popgreen, thick] (0,0) -- (0.9, 0) -- (1.35, 0.45) -- (0.45, 0.9) -- (-0.45, 0.45) -- cycle;
\draw[popgreen, thick] (0.9,0) -- (1.8, 0) -- (2.25, 0.45) -- (1.35, 0.9) -- (0.45, 0.45) -- cycle;
\node[popgreen, font=\tiny] at (0.9, -0.3) {Diagonale \textasciitilde 45 mm};
% Poisson adulte
\draw[popgreen!60!black, thick] (0.45, 0.2) -- (1.2, 0.4) -- (1.0, 0.6) -- (0.3, 0.5) -- cycle;
\node[popgreen!60!black, font=\tiny] at (0.9, 0.65) {Adulte \checkmark};
\end{scope}
% Maille Pratiquée (25mm)
\begin{scope}[xshift=3.5cm]
\node[font=\tiny\bfseries, popred] at (0, 1.0) {Pratiquée illégalement (25-30 mm)};
\draw[popred, thick] (0,0) -- (0.5, 0) -- (0.75, 0.25) -- (0.25, 0.5) -- (-0.25, 0.25) -- cycle;
\draw[popred, thick] (0.5,0) -- (1.0, 0) -- (1.25, 0.25) -- (0.75, 0.5) -- (0.25, 0.25) -- cycle;
\node[popred, font=\tiny] at (0.5, -0.3) {Diagonale \textasciitilde 25 mm};
% Poisson juvénile piégé
\draw[popred, thick] (0.2, 0.1) -- (0.5, 0.2) -- (0.4, 0.3) -- (0.1, 0.25) -- cycle;
\node[popred, font=\tiny] at (0.5, 0.35) {Juvénile \textbf{PIÉGÉ}};
% Poisson adulte traverse
\draw[popgreen!60!black, dashed, thick] (0.6, 0.1) -- (1.2, 0.3);
\node[popgreen!60!black, font=\tiny] at (1.1, 0.35) {Adulte \textbf{ÉCHAPPE}};
\end{scope}
\end{scope}
\end{tikzpicture}
\legende{Haut : Coupe schématique d'un filet maillant dérivant (gauche) et d'une senne tournante (droite). Bas : Impact de la taille de maille. La maille réglementaire (45-50 mm) laisse passer les juvéniles et capture les adultes reproducteurs. La maille illégale (25-30 mm) capture massivement les juvéniles (surpêche de croissance) et laisse échapper les gros reproducteurs.}
\end{popfigure}

\courssub{Espèces ciblées : la ressource halieutique camerounaise}

La diversité des milieux se reflète dans la composition des captures. La connaissance des cycles biologiques (reproduction, croissance, migrations) est indispensable pour comprendre les enjeux de gestion.

\begin{center}
\begin{tabularx}{\linewidth}{|l|l|X|X|}\hline
\rowcolor{popblueL}
\textbf{Espèce (Nom scientifique / Local)} & \textbf{Milieu de prédilection} & \textbf{Biologie clé (Taille 1ère maturité, Reproduction)} & \textbf{Valeur économique / Usage} \\ \hline
\textbf{Ethmalosa fimbriata} (Bonga, \emph{Kpéssi}) & Maritime, Estuarien (Wouri), Lagunes & Pélagique, planctonivore. Maturité \textasciitilde 18-20 cm. Frai toute l'année (pic saison des pluies). & \textbf{Reine de la pêche artisanale.} Fumée (bonga sec), séchée, salée. Base alimentaire urbaine/rurale. \\ \hline
\textbf{Sardinella maderensis / aurita} (Sardinelle) & Maritime (côte), Upwellings mineurs & Pélagique côtière. Maturité \textasciitilde 15-17 cm. Frai en saison fraîche (upwelling). & Fraîche, fumée, farine de poisson. Très prisée industrie conserves (usines Kribi/Douala). \\ \hline
\textbf{Oreochromis niloticus} (Tilapia du Nil) & Lacustre (Barrages, Maga, Tchad), Fluvial, Lagunaire & Démersal, herbivore/détritivore. Maturité \textasciitilde 15-20 cm. Incubation buccale. Prolifique. & \textbf{Espèce n°1 aquaculture \& pêche continentale.} Chair blanche, peu d'arêtes. Frais, fumé. \\ \hline
\textbf{Lates niloticus} (Capitaine du Nil, \emph{Giwan ruwa}) & Lacustre (Barrages, Tchad), Fluvial (Bénoué) & Prédateur apex. Croissance rapide. Maturité \textasciitilde 50-60 cm (M), 70 cm (F). Frai en saison haute. & \textbf{Poisson noble, haute valeur marchande.} Export régional (Nigeria, Tchad). Frais, fumé (filets). \\ \hline
\textbf{Clarias gariepinus} (Silure africain, \emph{Karfi}) & Ubiquiste (Lacs, Fleuves, Marigots, Rizicoles) & Respiration aérienne (organe suprabranchial). Supporte eau anoxique. Maturité \textasciitilde 30-40 cm. & Très rustique, transport vivant facile. Fumé, frais. Résilience forte. \\ \hline
\textbf{Penaeus notialis} (Crevette rose / blanche) & Maritime (côte), Estuarien (Wouri - nurserie) & Cycle de vie estuarien : reproduction en mer, post-larves remontent en estuaire (nurserie), juvéniles retournent en mer. & \textbf{Forte valeur export (congelée).} Pêche artisanale (crevettiers piroguiers) + Industrielle. Conflit majeur nurserie. \\ \hline
\end{tabularx}
\end{center}

\begin{aretenir}[Clé de lecture : Stratégies de vie et vulnérabilité]
Les espèces \textbf{r-sélectionnées} (Bonga, Sardinelle, Tilapia, Crevette) : cycle court, forte fécondité, maturation précoce. Résistent \emph{a priori} à une pêche modérée, mais s'effondrent si mailles trop petites (capture juvéniles) + destruction nurseries (mangroves, herbiers).
Les espèces \textbf{K-sélectionnées} (Capitaine, Silure gros) : cycle long, maturation tardive, faible fécondité relative. \textbf{Très vulnérables} à la surpêche (effort de pêche ciblé sur les gros = élimination des reproducteurs). Le Capitaine du Nil a disparu de nombreux lacs sahéliens par surpêche.
\end{aretenir}

\courssub{Transformations récentes : la modernisation par le bas (Motorisation \& Matériaux)}

Depuis les années 1980-90, la pêche artisanale subit une \cle{modernisation endogène}, tirée par la rentabilité marchande et l'accès aux intrants importés (moteurs, nylon, glace), sans encadrement technique suffisant.

\begin{propriete}[Effets systémiques de la motorisation (Hors-bord 8-40 CV)]
La motorisation des pirogues traditionnelles entraine quatre transformations majeures :
\begin{enumerate}
\item \textbf{Extension du rayon d'action :} De 2-5 km (pagaie/voile) à 15-30 km (hors-bord). Accès aux zones pélagiques lointaines, contournement des zones surpêchées côtières, mais \textbf{intrusion accrue dans la zone des 3 milles} (conflit chalutiers).
\item \textbf{Augmentation de l'effort de pêche (f) :} Plus de sorties/semaine, filets plus longs (200-500 m $\rightarrow$ 1000-2000 m), temps de calée réduit. Relation $C = q \times f \times B$ (Capture = Capturabilité $\times$ Effort $\times$ Biomasse) : $f$ augmente, $B$ baisse $\rightarrow$ risque d'effondrement.
\item \textbf{Transformation économique :} Coûts fixes (moteur \textasciitilde 300-800 kFCFA) + coûts variables (carburant \textasciitilde 50-70\% recette, huile, glace). \textbf{Endettement} fréquent auprès des mareyeuses/fournisseurs (crédit intrants). Passage d'une économie de subsistance à une micro-entreprise capitalistique.
\item \textbf{Externalités négatives :} Pollution hydrocarbures (vidanges, déversements), bruit (perturbation frai), sécurité en mer précaire (pas de gilets, pas de VHF, surcharge).
\end{enumerate}
\textit{Démonstration exigible au Bac : savoir expliquer le mécanisme "Motorisation $\rightarrow$ Effort $\uparrow$ $\rightarrow$ CPUE (Capture Par Unité d'Effort) $\downarrow$ $\rightarrow$ Surpêche".}
\end{propriete}

\begin{itemize}
\item \textbf{Filets en nylon (monofilament / multifilament) :} Remplacent fibres végétales. Avantages : solidité, invisibilité (monofilament), légèreté, prix bas (importation Asie). Inconvénients majeurs : \textbf{mailles plus petites faciles à tricoter}, \textbf{fantômes (filets perdus = pêche continue pendant mois/années)}, non biodégradables.
\item \textbf{Électronique embarquée (GPS, Sondeurs) :} Diffusion rapide depuis 2010 (téléphones Android + applis GPS, sondeurs portables Lowrance/Garmin d'occasion). Permet : repérage bancs (pélagiques), localisation fonds (démersaux), retour de nuit sécurisé, \textbf{marquage zones de conflit / chalutiers}. Change la donne : pêche de recherche $\rightarrow$ pêche de capture ciblée.
\item \textbf{Chaîne du froid artisanale :} Glacières (polystyrène expansé), glace en blocs (usines Douala, Kribi, Garoua), camions frigorifiques (location). Permet l'export vers villes intérieures (Yaoundé, Bertoua, Maroua) et pays limitrophes (Tchad, RCA, Gabon). \textbf{Rôle clé des femmes mareyeuses} (voir encadré).
\end{itemize}

\begin{savaistu}[Les femmes mareyeuses : le "pouvoir invisible" de la filière]
Au Cameroun, comme en Afrique de l'Ouest, les \textbf{femmes (mareyeuses / \emph{bayam-sellam}) contrôlent > 80\% de la transformation (fumage \emph{banda}, séchage, salage) et de la distribution locale}. Elles financent les sorties (carburant, glace, filets) via le système de \textbf{crédit intrants / engagement de vente} (le pêcheur vend obligatoirement à sa "mère" à prix convenu). Elles gèrent les fours améliorés (Chorkor, Altona - moins de bois, meilleure qualité), les entrepôts frigorifiques et les camions. Elles sont le \textbf{maillon clé de la sécurité alimentaire} (disponibilité poisson sec/fumé toute l'année, prix stable) et de la \textbf{transmission du capital social} (solidarité, tontines). Toute politique halieutique ignorant leur pouvoir de négociation et leur charge de travail (tri, transformation, vente, ménage) est vouée à l'échec.
\end{savaistu}

\courssub{Enjeux critiques : Surpêche, Pollution, Conflits d'usage}

La modernisation non régulée converge vers une crise de la ressource et de l'espace.

\begin{attention}[Piège conceptuel : "Pêche traditionnelle" vs "Pêche artisanale"]
\textbf{Ne pas confondre les deux termes.}
\begin{itemize}
\item \textbf{Pêche traditionnelle} = Pratiques historiques, engins/pirogues végétaux, non motorisés, subsistance/échange local. Quasi inexistante aujourd'hui à l'état pur.
\item \textbf{Pêche artisanale} = Catégorie \textbf{administrative et juridique} (Loi 94/01, Décret 95/413) définissant un régime d'accès (permis, zone 3 milles) et des caractéristiques techniques (puissance, taille bateau). Elle \textbf{inclut} la motorisation, le nylon, le GPS, la glace.
\item \textbf{Le programme vise l'évolution DU traditionnel VERS l'artisanal moderne.} L'analyse doit porter sur la \textbf{dynamique de transition} : gains (revenus, sécurité alimentaire) vs coûts (ressource, écosystème, inégalités).
\end{itemize}
\end{attention}

\textbf{1. Surpêche et pêche illicite (INDNR).}\par
\begin{itemize}
\item \textbf{Mailles non réglementaires :} Filets "mosquito" (maille < 25 mm) pour bonga/sardinelles ; filets "capitaine" (maille 80-100 mm) mais posés en zone de nurserie. \textbf{Effet de croissance (Growth overfishing) :} capture des juvéniles avant reproduction $\rightarrow$ biomasse future amputée.
\item \textbf{Pêche juvénile ciblée :} Crevettes juvéniles en estuaire (Wouri) pour l'export (queue de crevette) ; alevins de tilapia/capitaine pour l'aquaculture (repeuplement) ou consommation locale (poisson "bébé" fumé).
\item \textbf{Effort de pêche excessif :} Nombre de pirogues motorisées $\times$ longueur filets $\times$ sorties/jour. CPUE (Capture Par Unité d'Effort) en baisse constante depuis 15 ans sur le Bonga (maritime) et le Capitaine (Lagdo/Tchad).
\end{itemize}

\textbf{2. Pollution et dégradation des habitats essentiels.}\par
\begin{itemize}
\item \textbf{Estuaire Wouri (Douala-Bonabéri) :} Récepteur des effluents industriels (raffinerie SONARA - hydrocarbures, métaux ; brasseries, savonneries - DBO$_5$, nutriments ; tanneries - chrome), urbains (eaux usées non traitées, \textbf{déchets plastiques} macro/micro), portuaires (dragage, hydrocarbures). \textbf{Destruction mangroves} (bois de chauffe fumage, extension urbaine, riziculture). $\rightarrow$ Perte nurseries (crevettes, bonga, crabes), bioaccumulation métaux lourds (risque sanitaire consommateur).
\item \textbf{Lacs de barrage (Lagdo, Mapé) :} Eutrophisation (apports azotés/phosphatés agriculture coton/maïs amont), envasement rapide (durée de vie réduite), prolifération jacinthes d'eau (\emph{Eichhornia crassipes}) bloquant accès, désoxygenant fonds.
\item \textbf{Pollution par les engins :} Filets fantômes (nylon), batteries usagées (sondeurs/GPS), bidons huile moteur jetés en mer/rivière.
\end{itemize}

\textbf{3. Conflits pêcheurs artisanaux / Chalutiers industriels (Zone des 3 milles nautiques).}\par
C'est le conflit emblématique du littoral camerounais.
\begin{itemize}
\item \textbf{Cadre légal :} Loi 94/01 (Forêt, Faune, Pêche) + Décret 95/413 : Zone des \textbf{3 milles nautiques} (5,5 km) réservée \textbf{exclusivement} à la pêche artisanale. Au-delà : pêche industrielle (chalutiers crevettiers, senneurs pélagiques).
\item \textbf{Réalité :} Chalutiers (souvent pavillons étrangers, affrétés par sociétés locales) opèrent \textbf{de nuit} ou par mer mauvaise dans la zone des 3 milles.
\item \textbf{Mécanismes de destruction :}
    \begin{itemize}
    \item \textbf{Destruction physique des engins :} Chaluts (portes lourdes, bourrelets) arrachent/coupent les filets maillants calés (perte 500k-2M FCFA/pirogue/an).
    \item \textbf{Prédation de la ressource :} Chalutiers ciblent crevettes (taux de prise accessoire \emph{bycatch} 80-90\% : juvéniles poissons, raies, tortues) et poissons démersaux (rejetés morts ou vendus bas prix "poisson chalut" concurrencent pêcheurs).
    \item \textbf{Destruction habitat :} Ratissage fond (benthos, herbiers, nurseries).
    \end{itemize}
\item \textbf{Réponses artisanales :} Manifestations (blocus port Kribi 2019, 2021), saisies citoyennes de chalutiers, pétitions MINEPIA, demande surveillance satellitaire (VMS - Vessel Monitoring System) obligatoire et fonctionnelle, patrouilles mixtes (Marine + Pêche).
\end{itemize}

\begin{exoresolu}[Analyse d'un conflit d'usage : Le chalutier dans les 3 milles]
\textbf{Énoncé :} Un chalutier crevettier de 25 m est intercepté par la Marine à 2,5 milles nautiques des côtes de Kribi (zone réservée artisanale). Sa cale contient 2 tonnes de crevettes et 8 tonnes de prise accessoire (juvéniles de capitaines, carangues, raies, tortues). Les pêcheurs artisanaux de Kribi estiment leur perte de revenu à 15 millions FCFA/mois du fait de la destruction de leurs filets et de la raréfaction du poisson. Proposez une analyse géographie du conflit et deux mesures de gestion durables.
\tcblower
\textbf{Solution détaillée :}
\begin{enumerate}
\item \textbf{Analyse spatiale :} Violation flagrante de la zonation légale (3 milles). Le chalutier opère sur le \textbf{stock partagé} (crevettes adultes migrant vers le large, juvéniles poissons résidents). C'est un conflit \textbf{asymétrique} : capital intensif (chalutier, équipage salarié, export) vs capital extensif (pirogues, main-d'œuvre familiale, marché local).
\item \textbf{Analyse écologique :} Le \emph{bycatch} (80\% de la capture) détruit le \textbf{recrutement futur} (juvéniles capitaines, carangues) et la biodiversité (raies, tortues protégées). Le ratissage détruit l'habitat benthique (nurserie crevettes).
\item \textbf{Analyse socio-économique :} Pêcheurs artisans : perte d'engins (capital), baisse CPUE (revenu), insécurité alimentaire locale. Chalutier : valeur ajoutée export (devises), mais faible retombée locale (équipage souvent étranger, carburant détaxé ?).
\item \textbf{Mesure 1 (Contrôle/Surveillance) :} \textbf{Généralisation VMS (Vessel Monitoring System) temps réel + Patrouilles nocturnes dissuasives (Marine + MINEPIA) + Sanctions lourdes (saisie bateau, retrait licence, amendes proportionnelles valeur capture).} Coût partagé par redevances pêche industrielle.
\item \textbf{Mesure 2 (Gestion partagée / Co-gestion) :} \textbf{Mise en place de Comités Locaux de Gestion Pêche (CLGP) à Kribi/Limbé/Douala} intégrant délégués pêcheurs, mareyeuses, MINEPIA, Marine, Mairie. Rôle : surveillance participative (téléphones/GPS pêcheurs), définition zones tampons, négociation calendriers (repos biologique crevettes), gestion aires marines protégées (AMP) côtières.
\end{enumerate}
\end{exoresolu}

\courssub{Synthèse de la section : Une transition à encadrer pour la durabilité}

La pêche "traditionnelle" camerounaise a muté en une \textbf{pêche artisanale motorisée, connectée, marchande}. Elle assure \textbf{> 70\% de la production nationale} (\textasciitilde 300 000 t/an) et emploie directement \textasciitilde 200 000 pêcheurs + \textasciitilde 500 000 personnes en transformation/commerce (majoritairement femmes). Pourtant, le \textbf{déficit national} est comblé par \textbf{> 200 000 t/an d'importations} (poissons surgelés : maquereau, chinchard, sardinelle importée - "poisson glace"), signe que la productivité halieutique nationale stagne ou régresse par habitant.

\begin{exemplebox}[Bilan chiffré clé (Ordres de grandeur MINEPIA / FAO / Banque Mondiale)]
\begin{itemize}
\item \textbf{Production totale capture :} \textasciitilde 300 000 tonnes/an (dont > 70\% artisanale).
\item \textbf{Consommation apparente :} \textasciitilde 500 000 t/an $\rightarrow$ \textbf{Déficit \textasciitilde 200 000 t/an} (importations surgelées > 150 milliards FCFA/an).
\item \textbf{Part artisanale dans emplois directs :} \textasciitilde 200 000 pêcheurs (95\% hommes) + \textasciitilde 500 000 transformatrices/commerçantes (95\% femmes).
\item \textbf{Valeur à la première vente :} \textasciitilde 300-400 milliards FCFA/an.
\item \textbf{Taux d'exploitation stocks clés :} Bonga (pleinement exploité / surexploité localement), Crevettes (surexploité - juvéniles), Capitaine Lagdo/Tchad (surexploité), Tilapia (modérément exploité, résilient).
\item \textbf{Investissement public :} < 1\% budget national $\rightarrow$ Faible encadrement, recherche halieutique (IRAD) sous-dotée, surveillance maritime quasi inexistante.
\end{itemize}
\end{exemplebox}

\begin{aretenir}[Enjeux transversaux pour le développement durable (Objectif DD n°14 : Vie aquatique)]
La durabilité de la pêche artisanale camerounaise repose sur 4 piliers indissociables :
\begin{enumerate}
\item \textbf{Gouvernance spatiale efficace :} Respect strict zone 3 milles (VMS, patrouilles), création AMP (Aires Marines Protégées) côtières avec zones de non-prélèvement (réserves de reproduction), zonage estuarien (sanctuaires nurseries crevettes/bonga).
\item \textbf{Sélectivité des engins \& Réduction effort :} \textbf{Contrôle mailles à l'embarquement/au débarquement} (jauges mailles obligatoires), rachat filets illégaux, limitation longueur filets / nombre pirogues par zone (permis d'effort), promotion engins sélectifs (nasses, lignes, filets mailles carrées $\ge$ 45 mm).
\item \textbf{Valorisation de la chaîne de valeur (Femmes \& Froid) :} Appui fours améliorés (économie bois, qualité sanitaire), accès crédit/formation mareyeuses, développement infrastructures glace/frigo solaires (réduction pertes post-capture \textasciitilde 30-40\%), certification qualité (label "Pêche Artisanale Durable Cameroun").
\item \textbf{Gestion adaptative par la science participative :} Suivi débarquements (échantillonnage tailles, espèces, effort) par pêcheurs formés + IRAD/Universités $\rightarrow$ Avis scientifiques pour plans de gestion (quotas, repos biologiques, tailles minimales de capture). Intégration savoirs locaux (pêcheurs connaissent frayères, migrations).
\end{enumerate}
\end{aretenir}

\begin{exercice}[Application : Croquis commenté d'un espace de pêche artisanale]
\textbf{Consigne :} À partir de la carte de la section 3.1 et des connaissances du cours, réalisez un croquis légendé de l'espace de pêche artisanale de l'estuaire du Wouri (Douala-Bonabéri / Kribi / Limbé). Votre croquis doit faire apparaître :
\begin{enumerate}
\item Les unités spatiales (Domaine maritime 3 milles, Estuaire/Lagune, Mangroves, Zone portuaire industrielle).
\item Les infrastructures (Criées, Usines glace, Fours fumage, Usines conserves/congélation).
\item Les flux (Pirogues artisanales, Chalutiers industriels, Camions frigo vers Yaoundé/Nord, Export crevettes).
\item Les conflits / pressions (Zone chevauchement 3 milles, Rejets industriels, Déforestation mangroves, Évacuation filets fantômes).
\item Une légende organisée, hiérarchisée (surfaces, points, lignes, flèches) et un titre précis.
\end{enumerate}
\textit{Indication méthode :} Utilisez la méthode du croquis géographique (fond de carte simplifié $\rightarrow$ localisation précise $\rightarrow$ symboles proportionnels/qualitatifs $\rightarrow$ légende structurée $\rightarrow$ titre + orientation + échelle approximative).
\end{exercice}

\begin{corrige}[Exercice : Croquis estuaire Wouri / Côte Kribi-Limbé]
\textbf{Éléments attendus sur le croquis :}
\begin{itemize}
\item \textbf{Fond :} Contour côte (Kribi au Sud, Limbé au Nord-Ouest, Douala/Wouri au centre). Flèche Nord. Échelle graphique (0-50 km).
\item \textbf{Surfaces (Hachures / Couleurs claires) :}
    \item \textcolor{popblue!30}{Bleu clair} : Domaine maritime (limite 3 milles = trait pointillé bleu offshore).
    \item \textcolor{popgreen!30}{Vert clair} : Mangroves (Wouri, Mouanko, Rio del Rey) - \textbf{hachures croisées vert foncé} si dégradées.
    \item \textcolor{popyellow!30}{Jaune} : Zone portuaire/industrielle (Douala-Bonabéri, Kribi Port Autonome).
\item \textbf{Points (Symboles proportionnels) :}
    \item \textcolor{popred}{Carré rouge} : Criées / Postes de débarquement principaux (Douala, Kribi, Limbé, Idenau, Mouanko). Taille $\propto$ tonnage.
    \item \textcolor{poporange}{Triangle orange} : Usines transformation (Congélation Kribi/Douala, Conserves Limbé, Farine poisson).
    \item \textcolor{poppurple}{Losange violet} : Sites fours fumage (groupements femmes) - nombreux, dispersés rive Wouri, Kribi, Limbé.
    \item \textcolor{popblue}{Cercle bleu} : Usines glace.
\item \textbf{Lignes / Flèches :}
    \item \textcolor{popgreen}{Flèche verte épaisse} : Flux pirogues artisanales (côte $\leftrightarrow$ 3 milles, entrée estuaire marée montante).
    \item \textcolor{popred}{Flèche rouge pointillée} : Intrusion chalutiers industriels (zone 3 milles, nuit).
    \item \textcolor{popdark}{Flèche noire double} : Camions frigo (Douala/Kribi $\rightarrow$ Yaoundé, Bafoussam, Nord, Tchad/RCA).
    \item \textcolor{popblue}{Flèche bleue courbe} : Export crevettes/poissons nobles (Port Kribi/Douala $\rightarrow$ Europe/Asie).
\item \textbf{Signes de pression (Éclairs / Croix) :}
    \item \textcolor{popred}{Éclair rouge} : Zones conflit chalutiers/artisanaux (au large Kribi, Limbé, entrée Wouri).
    \item \textcolor{poporange}{Croix orange} : Points rejets industriels (SONARA Limbé, Usines Douala).
    \item \textcolor{popdark}{Trait discontinu noir} : Front déforestation mangroves (avancée urbaine, bois fumage).
\item \textbf{Légende :} Organisée par thèmes (Milieux, Acteurs/Infrastructures, Flux, Conflits/Pressions). Titre : \emph{"L'espace de la pêche artisanale dans l'estuaire du Wouri et la côte Sud-Ouest : activités, flux et tensions (2024)"}.
\end{itemize}
\end{corrige}

\coursec{L'artisanat traditionnel : diversité, organisation et évolution}

Après la pêche, activité de prélèvement direct sur les ressources aquatiques, nous abordons une autre pratique traditionnelle fondamentale des sociétés camerounaises : l'\cle{artisanat}. Comme la pêche, il fournit des biens de consommation et des revenus à des millions de ruraux et de citadins ; comme elle, il connaît aujourd'hui de profondes mutations (outillage moderne, marché touristique, exportation). Mais à la différence de la pêche, l'artisanat est une activité de \cle{transformation} : il part de matières premières locales (argile, bois, cuir, fibres, métaux) pour produire des objets utilitaires, rituels ou décoratifs. C'est aussi un formidable marqueur identitaire : chaque aire culturelle du Cameroun possède ses spécialités, ses gestes et ses styles.

\courssub{Qu'est-ce que l'artisanat traditionnel ?}

\begin{definition}[Artisanat et artisanat d'art]
L'\cle{artisanat} est une activité de production manuelle d'objets, réalisée par un ouvrier qualifié (l'artisan) maîtrisant un savoir-faire technique, généralement à petite échelle et avec un outillage simple. On distingue :
\begin{itemize}
\item l'\cle{artisanat d'usage} (ou utilitaire) : poteries de cuisson, paniers, calebasses, houes, meubles, destinés à la vie quotidienne ;
\item l'\cle{artisanat d'art} : production esthétique, unique ou en petite série, exigeant la maîtrise d'un geste technique complexe et porteur d'une forte valeur symbolique et identitaire (exemples : pipes en terre cuite Bamoun, bronzes de Foumban, masques des Grassfields).
\end{itemize}
\end{definition}

L'intuition est simple : là où l'industrie produit en série des objets identiques par des machines, l'artisan produit des objets \emph{uniques} par ses mains et son outillage. L'objet artisanal porte la trace d'un geste, d'un terroir et d'une culture : c'est ce qui fait sa valeur, notamment sur le marché touristique et de l'exportation.

\begin{exemplebox}[Un secteur informel massif]
Au Cameroun, le secteur artisanal est dominé par l'\cle{informel} : l'essentiel des artisans (plus de 90\,\% selon les estimations courantes) travaille sans immatriculation, sans comptabilité formelle, souvent dans des ateliers familiaux ou en bordure de route. La contribution de l'artisanat au PIB est donc difficile à quantifier, mais son rôle dans l'\cle{emploi} est massif : il occupe des jeunes, des femmes et des ruraux, et constitue souvent une activité complémentaire à l'agriculture. L'État soutient le secteur à travers le ministère en charge des Petites et Moyennes Entreprises, de l'Économie sociale et de l'Artisanat (MINPMEESA), qui mène des actions de formation, d'équipement et de structuration des artisans (programmes d'appui à l'artisanat, villages artisanaux, foires).
\end{exemplebox}

\courssub{Les grandes filières artisanales par aire culturelle}

Le Cameroun, « Afrique en miniature », présente une remarquable diversité artisanale, liée à la fois aux ressources disponibles et aux traditions culturelles de chaque région. On distingue schématiquement trois grandes aires de spécialisation.

\courssub{Le Grand-Nord (Adamaoua, Nord, Extrême-Nord)}

L'aire soudano-sahélienne du Nord est dominée par :
\begin{itemize}
\item le \cle{cuir} : tannage des peaux de bovins et d'ovins, fabrication de babouches, selles, sacs, poufs (Maroua est la capitale du cuir) ;
\item la \cle{poterie} : jarres, canaris, poteries décorées, souvent produites par les femmes ;
\item le \cle{tissage du coton} : pagnes tissés sur métiers à bandes (Garoua, Maroua), boubous brodés ;
\item la \cle{vannerie} : paniers, nattes, greniers en paille tressée ;
\item la \cle{ferronnerie} : forgerons produisant houes, machettes, lances, bijoux.
\end{itemize}

\courssub{L'Ouest et les Hauts-Plateaux (Grassfields)}

L'aire des Grassfields (Ouest, Nord-Ouest) est le cœur de l'artisanat d'art camerounais :
\begin{itemize}
\item le \cle{bois} : sculpture de statues, trônes, masques, pipes en terre cuite et en bois (chefferies Bamoun, Bamiléké) ;
\item les \cle{perles} : parures et objets perlés réservés autrefois aux notables et aux chefs ;
\item le \cle{bronze} et la \cle{fonte à la cire perdue} : statuettes, plaques, pipes, objets royaux (Foumban, Baham, Bafoussam) ;
\item le \cle{tissage} : le fameux pagne \cle{ndop} (toile de coton teintée à motifs réservés) et les sacs en rafia tressé, très prisés lors des cérémonies traditionnelles.
\end{itemize}

\courssub{Le Sud, le Centre et l'Est (aire forestière)}

Dans la zone forestière équatoriale, les matériaux végétaux dominent :
\begin{itemize}
\item la \cle{vannerie} : paniers, pièges, meubles en rotin et en bambou ;
\item la \cle{poterie} : poteries culinaires et jarres ;
\item la \cle{sculpture sur bois} : masques, statuettes, tam-tams, sièges ;
\item la \cle{pyrogravure} : décor gravé au fer rouge sur calebasses et bois ;
\item les \cle{écorces battues} : tissus végétaux (tapa) autrefois utilisés comme vêtements, notamment chez les populations forestières.
\end{itemize}

\begin{popfigure}
\begin{center}
\begin{tikzpicture}[scale=1.0]
% Contour très simplifié du Cameroun
\draw[thick, popdark, fill=popcream] (0,0) -- (2.2,0.2) -- (3.4,1.2) -- (3.6,2.6) -- (2.8,3.4) -- (3.2,4.6) -- (2.2,6.2) -- (1.2,7.2) -- (0.4,6.4) -- (0.8,5.0) -- (0.2,3.6) -- (-0.6,2.4) -- (-0.4,1.0) -- cycle;
% Aires culturelles
\fill[popgoldL, opacity=0.55] (0.2,4.2) -- (2.9,4.2) -- (3.2,4.6) -- (2.2,6.2) -- (1.2,7.2) -- (0.4,6.4) -- (0.8,5.0) -- cycle;
\fill[poppurpleL, opacity=0.55] (0.2,4.2) -- (2.9,4.2) -- (2.8,3.4) -- (1.2,3.2) -- (0.2,3.6) -- cycle;
\fill[popgreenL, opacity=0.55] (0.2,3.6) -- (1.2,3.2) -- (2.8,3.4) -- (3.6,2.6) -- (3.4,1.2) -- (2.2,0.2) -- (0,0) -- (-0.4,1.0) -- (-0.6,2.4) -- cycle;
% Nord
\node[font=\small\bfseries, popdark] at (1.6,5.9) {NORD};
\node[font=\scriptsize, popdark, align=center] at (1.6,5.3) {cuir, poterie,\\ tissage coton,\\ vannerie, ferronnerie};
% Ouest
\node[font=\small\bfseries, popdark] at (1.5,3.85) {OUEST};
\node[font=\scriptsize, popdark, align=center] at (1.5,3.45) {bois, bronze,\\ perles, ndop};
% Sud
\node[font=\small\bfseries, popdark] at (1.4,1.9) {SUD / CENTRE / EST};
\node[font=\scriptsize, popdark, align=center] at (1.4,1.2) {vannerie, poterie,\\ sculpture, pyrogravure,\\ écorces battues};
% Symboles ponctuels
% Marteau (fer) - Nord
\draw[popink, line width=1pt] (0.9,6.6) -- (1.15,6.85);
\draw[popink, line width=2pt] (1.05,6.9) -- (1.3,6.65);
% Pot (poterie) - Nord
\draw[popink, thick] (2.4,5.7) ellipse (0.12 and 0.15);
\draw[popink, thick] (2.32,5.82) -- (2.48,5.82);
% Ciseau (sculpteur) - Ouest
\draw[popink, line width=1pt] (0.7,4.0) -- (0.95,3.75);
\draw[popink, line width=1pt] (0.95,4.0) -- (0.7,3.75);
\draw[popink, fill=popink] (0.82,3.87) circle (0.04);
% Métier à tisser - Ouest
\draw[popink, thick] (2.3,3.7) rectangle (2.7,4.0);
\draw[popink] (2.4,3.7) -- (2.4,4.0) (2.5,3.7) -- (2.5,4.0) (2.6,3.7) -- (2.6,4.0);
% Ciseau Sud
\draw[popink, line width=1pt] (0.3,1.6) -- (0.55,1.35);
\draw[popink, line width=1pt] (0.55,1.6) -- (0.3,1.35);
\draw[popink, fill=popink] (0.42,1.47) circle (0.04);
% Pot Sud
\draw[popink, thick] (2.6,0.9) ellipse (0.12 and 0.15);
\draw[popink, thick] (2.52,1.02) -- (2.68,1.02);
% Villes
\fill[popdark] (1.7,6.9) circle (0.06); \node[font=\scriptsize, popdark, right] at (1.8,6.9) {Maroua};
\fill[popdark] (0.9,4.3) circle (0.06); \node[font=\scriptsize, popdark, left] at (0.85,4.3) {Foumban};
\fill[popdark] (1.0,2.6) circle (0.06); \node[font=\scriptsize, popdark, below] at (1.0,2.5) {Yaoundé};
\fill[popdark] (0.1,0.7) circle (0.06); \node[font=\scriptsize, popdark, below left] at (0.1,0.6) {Douala};
% Nord flèche
\draw[->, thick, popdark] (3.9,6.6) -- (3.9,7.2) node[above, font=\scriptsize\bfseries]{N};
\end{tikzpicture}
\end{center}
\legende{Carte schématique des aires de spécialisation artisanale du Cameroun. Symboles : marteau = ferronnerie, pot = poterie, ciseaux = sculpture sur bois, rectangle rayé = tissage. Les contours des aires culturelles sont très simplifiés.}
\end{popfigure}

\courssub{Les matières premières : un artisanat adossé aux ressources locales}

L'artisanat traditionnel tire sa richesse de l'utilisation des \cle{matières premières locales}, disponibles dans l'environnement immédiat :

\begin{center}
\begin{tabularx}{\linewidth}{|l|X|X|}\hline
\rowcolor{popblueL}
\textbf{Matériau} & \textbf{Exemples} & \textbf{Filières concernées} \\ \hline
Argiles & latéritiques (rouges), kaoliniques (blanches) & poterie, pipes, terres cuites \\ \hline
Bois & iroko, ébène, dibétou & sculpture, masques, meubles, tam-tams \\ \hline
Cuirs et peaux & peaux de bovins et d'ovins & babouches, selles, sacs, tambours \\ \hline
Fibres végétales & rafia, rotin, bambou, paille & vannerie, nattes, meubles, tissage \\ \hline
Minerais et métaux & fer latéritique, cuivre, laiton recyclé & ferronnerie, bronzes, bijoux \\ \hline
\end{tabularx}
\end{center}

Cette dépendance aux ressources locales explique la répartition géographique des filières : le bois et le rotin dominent en zone forestière (Sud), le cuir et la paille en zone d'élevage soudanienne (Nord), l'argile et le minerai de fer sur les plateaux de l'Ouest. On parle d'un artisanat \cle{adossé au terroir} : la ressource disponible oriente la spécialisation, et la spécialisation finit par identifier la région.

\courssub{Le statut de l'artisan : transmission et organisation sociale}

\courssub{Transmission héréditaire ou apprentissage libre}

Dans les sociétés traditionnelles camerounaises, le savoir-faire artisanal se transmet selon deux modalités :
\begin{itemize}
\item la \cle{transmission héréditaire} : dans certaines sociétés, les métiers sont réservés à des groupes sociaux spécialisés, parfois organisés en \cle{castes} endogames (on se marie à l'intérieur du groupe). C'est le cas des \cle{forgerons} (souvent à la fois respectés et marginalisés), des \cle{griots} (dépositaires de la parole et de la musique) et de certaines familles de \cle{potiers}. Le métier passe alors de père en fils (ou de mère en fille pour la poterie) ;
\item l'\cle{apprentissage libre} : de plus en plus répandu aujourd'hui, il consiste en un apprentissage auprès d'un maître-artisan, moyennant une durée de formation (3 à 7 ans) et parfois une compensation financière ou en nature. Le jeune apprenti devient compagnon, puis ouvre son propre atelier.
\end{itemize}

\courssub{L'organisation en corporations et quartiers spécialisés}

Les artisans se regroupent souvent par métier, sous le contrôle des \cle{chefferies traditionnelles} qui réglementent la production et la vente. Ce regroupement donne naissance à des \cle{quartiers spécialisés} :
\begin{itemize}
\item le quartier des artisans à \cle{Foumban} (royaume Bamoun), où sculpteurs, bronziers et brodeurs sont organisés autour du palais royal et du musée ;
\item le village des artisans de \cle{Baham} (Ouest), célèbre pour ses sculpteurs et ses fondeurs ;
\item les villages de potières du Nord et de l'Adamaoua.
\end{itemize}
Ce regroupement facilite l'apprentissage, le contrôle de qualité, l'accès aux matières premières et la visibilité commerciale (le client sait où trouver le produit).

\courssub{L'évolution de l'artisanat : de l'objet rituel à l'objet marchand}

\courssub{Le changement de fonction des objets}

Traditionnellement, l'objet artisanal avait une fonction \cle{utilitaire} (poterie de cuisson, panier, houe) ou \cle{rituelle} (masque de société secrète, statue royale, pipe de notable). Aujourd'hui, une part croissante de la production devient \cle{marchande} et \cle{touristique} : les masques et statuettes sont produits pour être vendus aux touristes et aux collectionneurs, les poteries deviennent des objets de décoration, le pagne ndop s'exporte comme tissu de mode. Ce passage de la fonction rituelle à la fonction commerciale transforme les formes, les tailles et les motifs (objets plus petits, plus transportables, motifs adaptés au goût des clients).

\courssub{La modernisation des outils et des designs}

L'outillage traditionnel (marteau de forge, ciseau à bois, tour à potier à main) est progressivement complété ou remplacé par des \cle{outils modernes} : disqueuse et perceuse électriques, chalumeau pour la soudure et la fonte, tour électrique pour la poterie et la menuiserie, machines à coudre pour la broderie. Parallèlement, de \cle{nouveaux designs} apparaissent sous l'effet des commandes : hôtels et restaurants commandent du mobilier et de la décoration « ethnique chic », des designers étrangers passent commande pour l'exportation (bijoux, luminaires, mobilier).

\begin{methode}[Analyser une filière artisanale : la chaîne opératoire]
Pour analyser rigoureusement une filière artisanale (poterie, bronze, cuir...), on suit la \cle{chaîne opératoire}, c'est-à-dire la succession des étapes de production :
\begin{enumerate}
\item \textbf{Approvisionnement en matières premières} : d'où vient l'argile, le bois, la peau, le métal ? Qui l'extrait et à quel coût ?
\item \textbf{Transformation} : quels gestes techniques, quels outils, quel temps de travail ?
\item \textbf{Finition} : décor, polissage, teinture, vernissage, cuisson.
\item \textbf{Commercialisation} : qui vend, où, à qui, à quel prix ?
\item \textbf{Recyclage et déchets} : que deviennent les chutes, les objets cassés, les résidus (ex. métal refondu, copeaux brûlés) ?
\end{enumerate}
Cette démarche permet de repérer à quelle étape se situent les difficultés (rareté du bois, faiblesse des prix, concurrence industrielle) et les marges de progrès.
\end{methode}

\begin{savaistu}[La fonte à la cire perdue, une technique millénaire]
La \cle{fonte à la cire perdue}, pratiquée par les bronziers Bamoun et Tikar (Foumban, Baham), permet de réaliser des pièces creuses et complexes (statuettes, pipes, cloches) impossibles à obtenir par martelage. Le sculpteur modèle l'objet en cire d'abeille, l'enrobe d'argile, puis chauffe le moule : la cire fond et s'écoule (« perdue »), laissant un vide dans lequel on coule le métal en fusion. Après refroidissement, on brise le moule : chaque pièce est donc unique. Cette technique, vieille de plusieurs millénaires, n'a pratiquement pas changé ; seul le métal a évolué : le cuivre et le bronze d'autrefois sont aujourd'hui souvent remplacés par du laiton et des métaux recyclés (robinetterie, cartouches, fil de cuivre).
\end{savaistu}

\courssub{La commercialisation : des marchés locaux au e-commerce}

Les débouchés de l'artisanat camerounais se sont considérablement diversifiés :
\begin{itemize}
\item les \cle{marchés locaux} : l'artisanat d'usage (poteries, paniers, meubles) s'écoule sur les marchés hebdomadaires ruraux et urbains ;
\item les \cle{foires et salons} : le salon PROMOTE à Yaoundé, le SIAO (Salon International de l'Artisanat de Ouagadougou, au Burkina Faso) offrent une vitrine nationale et internationale ;
\item les \cle{boutiques d'hôtels} et les villages artisanaux (Foumban, marché central de Yaoundé) ciblent la clientèle touristique ;
\item l'\cle{exportation} : commandes de designers européens et américains, vente aux galeries d'art africain ;
\item le \cle{e-commerce} naissant : vente en ligne via les réseaux sociaux et les plateformes spécialisées, encore freinée par les difficultés de paiement et de logistique.
\end{itemize}

\begin{popfigure}
\begin{center}
\begin{tikzpicture}
\begin{axis}[
  width=0.72\linewidth,
  height=6cm,
  ybar,
  bar width=0.9cm,
  ymin=0, ymax=70,
  ylabel={Part estimée (\%)},
  symbolic x coords={Usage local, Tourisme, Export},
  xtick=data,
  nodes near coords,
  nodes near coords style={font=\small\bfseries, popdark},
  axis lines=left,
  tick label style={font=\small},
  ylabel style={font=\small},
]
\addplot[fill=popblue, draw=popdark] coordinates {(Usage local,55) (Tourisme,30) (Export,15)};
\end{axis}
\end{tikzpicture}
\end{center}
\legende{Répartition estimée de la clientèle de l'artisanat camerounais (ordres de grandeur indicatifs, sans valeur statistique officielle) : l'usage local domine (environ 55\,\%), devant le tourisme (environ 30\,\%) et l'exportation (environ 15\,\%). Total : 55 + 30 + 15 = 100\,\%.}
\end{popfigure}

\begin{propriete}[Labellisation et Indications Géographiques]
La labellisation « \cle{Made in Cameroon} » et les \cle{Indications Géographiques} (IG) en cours de mise en place (exemples évoqués : cuir de Maroua, pipes de Foumban) visent à valoriser le terroir et à protéger le savoir-faire artisanal. Une IG garantit au consommateur qu'un produit tire ses qualités de son origine géographique et d'un savoir-faire local codifié ; elle protège les artisans contre les contrefaçons et leur permet de vendre plus cher. C'est un outil juridique de développement, déjà utilisé avec succès pour d'autres produits africains (café, cacao, tissus).
\end{propriete}

\begin{attention}[Modernisation et perte d'authenticité]
L'usage croissant de \cle{produits chimiques} (teintures synthétiques, vernis industriels, acides pour le traitement du cuir) remplace progressivement les techniques naturelles traditionnelles (teintures aux écorces et aux feuilles, tannage végétal, boues et huiles naturelles). Cette substitution pose trois problèmes : une \textbf{perte d'authenticité} des produits (le client recherche justement le « fait main naturel »), des \textbf{risques sanitaires} pour les artisans (manipulation d'acides sans protection) et une \textbf{pollution} des sols et des eaux autour des ateliers. Dans une copie, ne présente pas la modernisation comme uniquement positive : elle a un coût culturel et environnemental.
\end{attention}

\begin{exoresolu}[Analyser la filière du bronze de Foumban]
Énoncé : À l'aide de vos connaissances, analysez la filière du bronze de Foumban en suivant la chaîne opératoire, puis montrez comment elle illustre les mutations de l'artisanat camerounais.
\tcblower
Solution détaillée :
\begin{enumerate}
\item \textbf{Approvisionnement} : le métal (laiton, cuivre) provient aujourd'hui largement du recyclage (vieux robinets, cartouches, fil électrique) ; la cire d'abeille est collectée localement.
\item \textbf{Transformation} : modelage de la cire, moulage en argile, fonte à la cire perdue dans des fours traditionnels, de plus en plus aidés par le chalumeau.
\item \textbf{Finition} : ébarbage, polissage, patine ; certains artisans utilisent désormais la disqueuse électrique.
\item \textbf{Commercialisation} : vente dans les ateliers du quartier des artisans, au musée du palais royal, aux touristes, et par commandes pour l'export (galeries, designers).
\item \textbf{Recyclage} : les scories et chutes de métal sont refondues ; les moules brisés retournent à la terre.
\end{enumerate}
Cette filière illustre les mutations de l'artisanat : matière première recyclée (et non plus extraite), outillage partiellement modernisé, production passée des objets royaux et rituels aux objets touristiques et d'exportation, et organisation maintenue autour de la chefferie et du quartier spécialisé. Elle montre aussi la fragilité du secteur : dépendance au tourisme, concurrence des copies industrielles, protection encore insuffisante du savoir-faire (IG non pleinement opérationnelles).
\end{exoresolu}

\begin{exercice}[À faire]
Choisis une filière artisanale de ta région (poterie, vannerie, couture, sculpture...). Décris sa chaîne opératoire en cinq étapes, puis rédige un court paragraphe justifié répondant à la question : cette filière est-elle en voie de modernisation ou de disparition ? Argumente avec au moins deux éléments concrets observés.
\end{exercice}

\begin{corrige}[Exercice]
Exemple de réponse attendue (filière vannerie en zone forestière). Chaîne opératoire : 1) collecte du rotin et du bambou en forêt ; 2) fendage et séchage des fibres ; 3) tressage à la main (fond, puis montée des parois) ; 4) finition (coupe des brins, parfois teinture) ; 5) vente au marché hebdomadaire ou sur commande. Argumentation possible : la filière est en difficulté car les paniers en plastique importés, moins chers et plus durables, concurrencent la vannerie ; de plus, la ressource en rotin s'éloigne des villages (déforestation). Mais des signes de modernisation existent : nouveaux modèles (corbeilles décoratives, meubles en rotin pour hôtels) et vente en ligne par des coopératives. La conclusion doit être justifiée par ces observations concrètes.
\end{corrige}

\begin{aretenir}[L'essentiel de la section]
\begin{itemize}
\item L'artisanat camerounais s'organise en trois grandes aires : le \cle{Nord} (cuir, poterie, tissage coton, vannerie, ferronnerie), l'\cle{Ouest}/Hauts-Plateaux (sculpture sur bois, perles, bronze à la cire perdue, ndop), le \cle{Sud}/Centre/Est (vannerie, poterie, sculpture, pyrogravure, écorces battues).
\item Il repose sur des matières premières locales (argile, bois, cuir, fibres, minerais et métaux) et sur une transmission héréditaire (castes) ou par apprentissage libre, avec une organisation en corporations et quartiers spécialisés (Foumban, Baham).
\item Il évolue : passage de l'objet utilitaire/rituel à l'objet marchand et touristique, modernisation des outils, nouveaux designs, diversification des débouchés (salons, hôtels, export, e-commerce).
\item Ses défis : informalité (plus de 90\,\% des artisans), concurrence industrielle, perte d'authenticité liée aux produits chimiques, protection insuffisante du savoir-faire (enjeu des IG et du label « Made in Cameroon »).
\end{itemize}
\end{aretenir}

\coursec{Synthèse : Enjeux transversaux et développement durable des pratiques traditionnelles}

\courssub{Transition : L'artisanat au cœur du système productif traditionnel}

Après avoir analysé successivement l'élevage, la pêche et l'artisanat, il apparaît clairement que ces trois activités ne forment pas des mondes étanches. Elles constituent les piliers d'un même \textbf{système agropastoral et artisanal intégré}, caractéristique des sociétés rurales camerounaises. L'artisanat, que nous venons d'étudier, ne se limite pas à la production d'objets : il fournit les outils de l'élevage (cannes, couteaux, cordes en cuir), les engins de la pêche (nasses, filets, pirogues) et transforme les produits bruts (lait, cuir, poisson fumé, bois sculpté). Cette interdépendance technique et économique crée une communauté de destin face aux défis contemporains. Comprendre les enjeux transversaux, c'est saisir comment la pression sur l'espace, le climat, la transmission des savoirs et les cadres juridiques redessinent ensemble l'avenir de ces pratiques.

\courssub{Pression foncière : la compétition pour l'espace et les ressources naturelles}

\noindent\textbf{Intuition :} Imaginez une carte du Nord-Cameroun ou de l'Adamaoua. Au même endroit, l'agriculteur veut labourer pour le coton ou le maïs, l'éleveur a besoin d'un parcours pour son troupeau, le pêcheur accède à la rivière ou au plan d'eau, et l'artisan prélève l'argile, le bois ou les teintures. L'espace est fini, les besoins croissent : c'est l'équation de la \cle{pression foncière}.

\begin{definition}[Pression foncière et conflit d'usage]
La \cle{pression foncière} désigne l'intensification de la demande sur un espace limité, entraînant une compétition entre usages rivaux (agricole, pastoral, halieutique, forestier, artisanal, urbain, minier). Au Cameroun, elle est exacerbée par la croissance démographique (taux > 2,6\%/an), l'expansion des cultures de rente (palmier à huile, coton, hévéa) et l'insécurité foncière (droit coutumier vs droit écrit, Loi n°74-1 du 6 juillet 1974).
\end{definition}

\noindent\textbf{Mécanismes et observations de terrain :}
\begin{itemize}
    \item \textbf{Agriculture extensive vs Parcours :} L'avancée du front pionnier agricole (ex. : plaine de la Bénoué, Diamaré) grignote les couloirs de transhumance et les zones de pâturage sec. Résultat : raccourcissement de la transhumance, surpâturage résiduel, conflits agriculteurs-éleveurs (ex. : Logone-et-Chari, Mayo-Rey).
    \item \textbf{Zones humides : agriculture de décrue vs pêche vs pâturage :} Les plaines d'inondation (Yaérés, Logone, Mbam, Wouri) sont des zones de convergence. En saison sèche, elles sont pâturages de choix ; en saison des pluies, zones de frai et de pêche ; en décrue, terres agricoles fertiles. L'aménagement hydro-agricole (barrages, riziculture irriguée type SEMRY, SODERIM) rigidifie ces cycles, bloquant la migration des poissons et l'accès à l'eau pour le bétail.
    \item \textbf{Espaces matières premières (Artisanat) :} Les carrières d'argile (Foumban, Baham), les forêts de teintures (indigo, garance) ou de bois d'œuvre (ébène, iroko) sont souvent situées sur des terres convoitées pour l'agriculture ou l'habitat. La déforestation (taux ~0,6\%/an au Cameroun) raréfie le bois de sculpture et le combustible pour la cuisson potière.
\end{itemize}

\begin{exemplebox}[Cas concret : Le bassin du Lac Tchad et la plaine du Logone]
Dans l'Extrême-Nord, la régression du Lac Tchad (-90\% de sa surface depuis les années 60) concentre les activités sur les terres émergées. Les \emph{karal} (champs de décrue), les pâturages de \emph{bourgou} (\textit{Echinochloa stagnina}) et les zones de pêche se superposent. Les aménagements de la SEMRY (riz) et de la SODERIM (riz/maïs) ont fermé des couloirs de transhumance et modifié l'hydrologie, réduisant les zones de frai. Les artisans potiers (Kotoko) voient leurs gisements d'argile inondés ou asséchés artificiellement.
\end{exemplebox}

\begin{popfigure}
\begin{tikzpicture}[scale=0.9, every node/.style={font=\small}]
% Styles
\tikzset{
    zone/.style={draw=popdark, thick, fill=popblue!10, rounded corners=5pt, minimum height=1.8cm, minimum width=3.5cm, align=center},
    arrow/.style={->, >=Stealth, thick, poporange},
    conflict/.style={draw=poppink, very thick, dashed, fill=poppink!10, rounded corners=5pt, minimum height=1.2cm, minimum width=4cm, align=center},
    title/.style={font=\bfseries\large, text=popblue},
    subtitle/.style={font=\bfseries\small, text=popdark},
}
% Title
\node[title] at (0, 5.5) {Schéma synthèse : Compétition pour l'espace et les ressources (Extrême-Nord / Nord)};
% Zones
\node[zone, align=center] (agri) at (-5, 3){AGRICULTURE EXTENSIVE\\\textit{(Coton, Maïs, Riz irrigué)}\\\textbf{Besoin :} Terres planes, eau, engrais};
\node[zone, align=center] (elev) at (5, 3){ÉLEVAGE TRADITIONNEL\\\textit{(Transhumance, Pâturage)}\\\textbf{Besoin :} Parcours, couloirs, points d'eau, bourgou};
\node[zone, align=center] (peche) at (0, 0.5){PÊCHE TRADITIONNELLE\\\textit{(Lac, Rivière, Yaérés)}\\\textbf{Besoin :} Eau permanente, frayères, accès rives};
\node[zone, align=center] (art) at (0, -2.5){ARTISANAT\\\textit{(Poterie, Teinture, Vannerie)}\\\textbf{Besoin :} Argile, bois, teintures, espace atelier};
% Arrows competition
\draw[arrow, bend left=20] (agri.east) to node[midway, above, font=\tiny] {Empiètement sur parcours} (elev.west);
\draw[arrow, bend right=20] (elev.west) to node[midway, below, font=\tiny] {Piétinement cultures} (agri.east);
\draw[arrow] (agri.south) -- node[right, font=\tiny] {Assèchement / Barrages} (peche.north);
\draw[arrow] (elev.south) -- node[left, font=\tiny] {Pollution / Piétinement rives} (peche.north);
\draw[arrow] (art.north west) -- node[above left, font=\tiny] {Déforestation / Accès argile} (agri.south west);
\draw[arrow] (art.north east) -- node[above right, font=\tiny] {Bois énergie / Teintures} (elev.south east);
% Conflict Box
\node[conflict, align=center] at (0, 5){CONFLITS D'USAGE \& INSÉCURITÉ FONCIÈRE\\ (Droit coutumier vs Loi 74-1 vs Cadre pastoral)};
% Drivers
\node[font=\bfseries, align=center, text width=4cm] at (-5, 1) {MOTEURS EXTERNES\\ Démographie + Climatique + Politiques publiques (SEMRY, ZICGC)};
\node[font=\bfseries, align=center, text width=4cm] at (5, 1) {CONSÉQUENCES\\ Conflits violents, Appauvrissement sols, Perte biodiversité, Exode rural};
% North Arrow
\node at (-7, -3.5) {\textbf{N}};
\draw[->, thick] (-7, -3.8) -- (-7, -3.2);
% Scale
\node at (7, -3.5) {Échelle indicative : 1 cm $\approx$ 50 km};
\end{tikzpicture}
\legende{Schéma synthèse des interactions et conflits d'usage de l'espace dans les plaines du Nord-Cameroun. Les flèches oranges indiquent les pressions directes d'un secteur sur l'autre. La zone rose centrale résume la conflictualité foncière.}
\end{popfigure}

\courssub{Changement climatique : multiplicateur de vulnérabilité}

\noindent\textbf{Intuition :} Le climat n'est plus un décor stable, c'est un acteur instable. La variabilité pluviométrique (sécheresses récurrentes au Nord, inondations soudaines dans les plaines) frappe simultanément les quatre piliers : l'herbe ne pousse pas (élevage), l'eau manque ou détruit (pêche/agriculture), le bois meurt (artisanat).

\begin{propriete}[Impacts différenciés mais systémiques du changement climatique]
\begin{itemize}
    \item \textbf{Ressources fourragères (Élevage) :} Au Sahel (Maroua, Kousseri), la pluviométrie a baissé de ~15-20\% depuis les années 1970. La biomasse herbacée chute, la période de soudure fourragère s'allonge (4 à 6 mois). Mortalité accrue, baisse production lait/viande.
    \item \textbf{Ressources halieutiques (Pêche) :} Variabilité des crues (hauteur, date, durée). Crues faibles = pas de mise en eau des yaérés = effondrement de la reproduction (poissons "de crue" comme \textit{Alestes}, \textit{Clarias}). Crues violentes = destruction des engins, érosion rives.
    \item \textbf{Ressources forestières (Artisanat/Énergie) :} Mortalité des essences de teinture (\textit{Indigofera}, \textit{Lonchocarpus}) et de sculpture. Pénurie de bois de chauffe pour la cuisson potière (fours traditionnels à 800-1000°C).
\end{itemize}
\end{propriete}

\begin{experience}[Analyse de données climatiques : Maroua (1980-2020)]
\textbf{Protocole :} Exploiter les séries pluviométriques annuelles (source : ASECNA / OMM).
\textbf{Observations attendues :}
\begin{enumerate}
    \item Tendance à la baisse de la pluviométrie totale (mm).
    \item Augmentation de l'écart-type (variabilité interannuelle forte).
    \item Décalage du début de saison (instabilité dates de semis/transhumance).
    \item Augmentation fréquence événements extrêmes (> 50 mm/jour = inondations ; < 400 mm/an = sécheresse sévère).
\end{enumerate}
\textbf{Interprétation :} Cette instabilité rend la planification pastorale (date de départ en transhumance) et halieutique (date de mise à l'eau des nasses) aléatoire. Elle favorise les stratégies de survie à court terme (surpêche, coupe bois vert) au détriment de la gestion durable.
\end{experience}

\begin{popfigure}
\begin{tikzpicture}
\begin{axis}[
    width=\linewidth, height=8cm,
    title={Évolution simulée pluviométrie annuelle Maroua (mm) \& Moyenne mobile 5 ans},
    xlabel={Année}, ylabel={Pluviométrie (mm)},
    xmin=1980, xmax=2020, ymin=400, ymax=1200,
    grid=major, legend pos=south west,
]
% Données simulées réalistes (variabilité sahélienne)
\addplot[popmuted, mark=*, mark size=1.5, only marks] coordinates {
(1980, 780) (1981, 620) (1982, 550) (1983, 480) (1984, 420) (1985, 650)
(1986, 710) (1987, 590) (1988, 850) (1989, 760) (1990, 680) (1991, 920)
(1992, 540) (1993, 730) (1994, 810) (1995, 600) (1996, 770) (1997, 880)
(1998, 650) (1999, 720) (2000, 580) (2001, 840) (2002, 700) (2003, 610)
(2004, 790) (2005, 550) (2006, 820) (2007, 670) (2008, 910) (2009, 740)
(2010, 630) (2011, 560) (2012, 890) (2013, 780) (2014, 690) (2015, 520)
(2016, 860) (2017, 750) (2018, 640) (2019, 930) (2020, 710)
};
\addlegendentry{Pluviométrie annuelle}
% Moyenne mobile 5 ans (calculée approximativement pour l'exemple)
\addplot[popcurve, thick, domain=1982:2018, samples=37] {720 + 30*sin(deg(x-1980)/5) - 0.5*(x-1980)}; % Tendance légère baisse + cycle
\addlegendentry{Tendance (Moy. mobile 5 ans / Ajustement)}
\draw[dashed, poppink, thick] (axis cs:1980, 750) -- (axis cs:2020, 700) node[right, font=\small] {Tendance baissière $\approx$ -5 mm/an};
\end{axis}
\end{tikzpicture}
\legende{Variabilité pluviométrique à Maroua (Extrême-Nord). La ligne pointillée rose montre la tendance structurelle à la baisse, les points la forte variabilité interannuelle. Source : données types ASECNA/OMM.}
\end{popfigure}

\courssub{Transmission des savoirs : la rupture intergénérationnelle}

\noindent\textbf{Intuition :} Le "savoir-faire" ne s'écrit pas dans les manuels, il se transmet par l'observation, l'imitation, l'initiation (maître/apprenti). Si le jeune part à l'école ou en ville, la chaîne se brise. La mémoire technique (gestes du potier, nœuds du filet, chants de transhumance, recettes de teinture) s'efface en une génération.

\begin{definition}[Patrimoine culturel immatériel (PCI) et transmission]
Au sens de l'UNESCO (Convention 2003), le PCI inclut les \cle{savoirs traditionnels} : techniques de production, connaissances de la nature (phénologie, éthnozoologie), rites associés. La transmission est \textbf{verticale} (parents $\to$ enfants) et \textbf{horizontale} (pairs, sociétés d'initiation). Sa rupture entraîne une \cle{perte irréversible de biodiversité culturelle}.
\end{definition}

\noindent\textbf{Facteurs de rupture (Enquête type) :}
\begin{enumerate}
    \item \textbf{Scolarisation formelle :} Temps d'école vs temps d'apprentissage "sur le tas". Dévalorisation sociale des métiers manuels ("métiers d'échec scolaire").
    \item \textbf{Exode rural / Migration :} Jeunes hommes partis (élevage, pêche), jeunes filles mariées hors du village (poterie, teinture). Vieillissement des détenteurs du savoir (âge moyen > 55 ans dans beaucoup de chefferies).
    \item \textbf{Substitution technique :} Pièces détachées industrielles (pompes, filets nylon, bassines plastique) remplacent objets artisanaux. Perte de la demande = perte de l'incitation à apprendre.
    \item \textbf{Perte rituelle :} Certaines techniques sont liées à des rites (protection du four, bénédiction des filets, rites de transhumance). La conversion religieuse (islam, christianisme) ou la modernité peuvent interdire ces rites, vidant la technique de son sens symbolique et de sa transmission codifiée.
\end{enumerate}

\begin{attention}[Piège conceptuel : "Traditionnel" $\neq$ "Archaïque/Figé"]
Ne pas opposer \textbf{Traditionnel} (figé, archaïque, inefficace) et \textbf{Moderne} (performant, innovant, rationnel). La réalité est un \textbf{CONTINUUM d'hybridation} :
\begin{itemize}
    \item Pirogue en bois + moteur hors-bord 15-40 CV (pêche).
    \item Four à poterie traditionnel + isolation en brique réfractaire / thermomètre (artisanat).
    \item Tour à bois manuel $\to$ tour à bois électrique (menuiserie Baham/Foumban).
    \item Téléphone mobile pour cours coton/bétail/poisson + GPS pour transhumance (élevage).
\end{itemize}
L'innovation endogène existe : l'artisan/éleveur/pêcheur adapte \textbf{en permanence}. Le défi n'est pas de "moderniser" mais d'\textbf{accompagner l'hybridation} sans perdre le cœur du savoir (gestuelle, connaissance milieu).
\end{attention}

\begin{savaistu}[Le Ndop et le Toghu : vers une protection par Indication Géographique (OAPI)]
Le \textbf{Ndop} (tissu de coton bleu indigo à motifs blancs, réserve, tissé par les Bamileke, teint à l'indigo \textit{Indigofera arrecta}, motifs à la réserve cousue/nouée) et le \textbf{Toghu} (velours noir brodé de fils colorés, Nord-Ouest, porté par les notables) sont des symboles identitaires forts.
Actuellement, des démarches sont en cours auprès de l'\textbf{OAPI} (Organisation Africaine de la Propriété Intellectuelle, Yaoundé) pour les enregistrer comme \textbf{Indications Géographiques (IG)}.
\textbf{Enjeu :} Protéger le nom contre l'usurpation (contrefaçons chinoises imprimées), valoriser le prix à la production (revenus artisans), obliger le respect du cahier des charges (coton local, teinture naturelle, tissage main). C'est un levier juridique pour la transmission : "Si ça rapporte, les jeunes restent."
\end{savaistu}

\courssub{Valorisation économique : du survie au marché durable}

\noindent\textbf{Intuition :} Ces activités ne sont pas des "musées vivants". Elles nourrissent des millions de Camerounais (sécurité alimentaire ODD 2) et génèrent des revenus (ODD 1, 8). Le défi : passer de l'économie de subsistance / marché local informel à des \cle{circuits courts}, \cle{commerce équitable}, \cle{tourisme culturel responsable}, sans aliéner les producteurs.

\begin{methode}[Analyse de chaîne de valeur (Filière artisanale / produits d'élevage / pêche)]
Pour proposer une valorisation durable, suivre ces étapes :
\begin{enumerate}
    \item \textbf{Cartographier les acteurs :} Producteur $\to$ Collecteur/Intermédiaire $\to$ Grossiste/Transformateur $\to$ Détaillant/Exportateur $\to$ Consommateur final.
    \item \textbf{Quantifier la valeur ajoutée :} Prix payé au producteur vs Prix final. Part du producteur souvent < 20-30\%.
    \item \textbf{Identifier les goulets :} Qualité irrégulière, manque de conditionnement, transport, absence de marque/label, financement.
    \item \textbf{Proposer des leviers :} Groupement (GIC, Coopérative), Certification (Label "Produit du Terroir", Bio, IG/OAPI), E-commerce / Marketplace (Jumia, Afrimarket, sites spécialisés), Contrats de filière avec hôtels/restaurants (circuit court).
\end{enumerate}
\end{methode}

\begin{exemplebox}[Chiffres clés : Tourisme culturel et artisanat (Estimation MINTOUL/ONMT)]
\begin{itemize}
    \item \textbf{Festivals majeurs :} Ngondo (Douala, Sawa), Nyem-Nyem (Ngaoundéré, Adamaoua), Mpo'o (Edéa, Bakoko), Medumba (Bafang), Ngoun (Foumban). Fréquentation : 50 000 à 200 000 visiteurs/édition.
    \item \textbf{Revenus directs artisans/guides :} Estimés 50-200 millions FCFA / festival pour les exposants officiels.
    \item \textbf{Risque "Folklorisation" :} Spectacle déconnecté du sens (danses sacrées devenues "animation"), objets produits pour le touriste (qualité basse, motifs déformés), captation des revenus par intermédiaires extérieurs (agences, hôtels), non-reversement aux communautés détentrices du savoir.
    \item \textbf{Bonne pratique :} Festival Ngondo : gestion par l'association des peuples Sawa, cahier des charges exposition, pourcentage sur ventes reversé à un fonds de développement communautaire.
\end{itemize}
\end{exemplebox}

\begin{exoresolu}[Cas pratique : Coopérative de potières de Baham (Ouest)]
\textbf{Énoncé :} Un groupement de 40 potières produit des marmites traditionnelles (cuisson bois, four ouvert). Revenus faibles, pénurie bois, casse élevée (30\%), vente au bord de route prix dérisoire. Proposer un plan de valorisation sur 3 ans.
\tcblower
\textbf{Solution structurée :}
\textbf{An 1 - Diagnostic \& Organisation :} Création GIC "Terre de Baham". Inventaire savoirs (argiles, formes, décors). Formation gestion/comptabilité (ONG d'appui).
\textbf{An 2 - Innovation technique \& Qualité :} Construction 2 fours améliorés (type "four à voûte" ou "four à gazéification" : rendement bois -50\%, cuisson homogène, casse < 5\%). Standardisation tailles (S, M, L). Création marque collective "Baham Pottery" + logo. Tests labo (aptitude alimentaire, plomb/cadmium).
\textbf{An 3 - Marché \& Circuit court :} Showroom à Bafoussam + site web / Instagram. Convention avec 10 restaurants "Terroir" (Yaoundé, Douala) pour marmites "signature". Participation Salon PROMOTE / Foire de Yaoundé. Demande IG "Poterie de Baham" à l'OAPI.
\textbf{Indicateurs de réussite :} Revenu x3, taux casse < 5\%, bois -50\%, 10 emplois jeunes (préparation argile, livraison, digital), transmission assurée (3 apprenties formées).
\end{exoresolu}

\courssub{Cadre institutionnel : les leviers juridiques pour un développement durable}

\noindent\textbf{Intuition :} Le droit n'est pas seulement contraignant, il peut être \textbf{protecteur} et \textbf{habilitant}. Connaître les textes permet aux acteurs de revendiquer leurs droits (accès aux ressources, propriété intellectuelle) et aux pouvoirs publics de cibler les politiques d'appui.

\begin{center}
\begin{tabularx}{\linewidth}{|l|X|X|}\hline
\rowcolor{popblueL}
\textbf{Texte / Cadre} & \textbf{Objet principal} & \textbf{Levier pour pratiques traditionnelles} \\ \hline
\textbf{Loi n°2011/022 (14 juil. 2011)} & Organisation de l'artisanat & Carte d'artisan (statut social, accès formation/financement), Chambre des Métiers, Label "Artisanat du Cameroun", protection dessins/modèles. \\ \hline
\textbf{Cadre pastoral (Loi 74-1 foncier, Décret 78/263 transhumance, Loi 2005/017 forêt/faune)} & Gestion parcours, transhumance, points d'eau & Sécurisation couloirs/pâturages, Commissions de gestion foncière pastorale, prévention conflits, reconnaissance droits d'usage coutumiers. \\ \hline
\textbf{Loi n°94/01 (Forestière) \& Loi n°98/005 (Régime faune/pêche)} & Pêche continentale & Plans de gestion participative (PGP), zones de reproduction protégées, droits pêcheurs artisanaux vs industriels, co-gestion. \\ \hline
\textbf{OAPI (Accord de Bangui 1999, révisé 2015)} & Propriété Intellectuelle & \textbf{Indications Géographiques (IG)}, Marques collectives, Dessins/Modèles industriels. Protège le "savoir-faire" comme actif économique (Ndop, Toghu, Poterie Baham, Couteaux Bafia...). \\ \hline
\textbf{ODD (Agenda 2030 ONU)} & Cadre global & \textbf{ODD 1} (Pas de pauvreté), \textbf{ODD 2} (Faim zéro), \textbf{ODD 8} (Travail décent/Croissance), \textbf{ODD 12} (Conso/Prod responsables), \textbf{ODD 14} (Vie aquatique), \textbf{ODD 15} (Vie terrestre). Alignement politiques nationales (SND30). \\ \hline
\end{tabularx}
\end{center}

\begin{aretenir}[Synthèse des leviers institutionnels]
\begin{itemize}
    \item \textbf{Sécurisation foncière :} Application effective du cadre pastoral (délimitation couloirs, bornage points d'eau) + Réforme foncière (reconnaissance droits coutumiers collectifs).
    \item \textbf{Reconnaissance statutaire :} Carte d'artisan (Loi 2011/022) = accès microcrédit (MC2, ANPME), formation continue (CFP), couverture sociale (CNPS secteur informel).
    \item \textbf{Valorisation par la PI :} Dépôt IG / Marques collectives à l'OAPI = monopole d'usage du nom, cahier des charges = garantie qualité + transmission.
    \item \textbf{Gouvernance participative :} Comités de gestion pêche (PGP), Commissions pastorales, GIC/Coopératives = acteurs légitimes pour négocier avec l'État/privés.
\end{itemize}
\end{aretenir}

\courssub{Croquis de synthèse guidé (TP) : Espaces productifs traditionnels du Cameroun}

\noindent\textbf{Consigne :} Réaliser un croquis de synthèse à l'échelle du Cameroun intégrant les 3 secteurs (Élevage, Pêche, Artisanat) avec une légende hiérarchisée. Utiliser la base cartographique fournie (contour, villes, fleuves, axes).

\begin{methode}[Réalisation du croquis de synthèse (Épreuve type Bac)]
\begin{enumerate}
    \item \textbf{Analyse du sujet :} "Espaces productifs traditionnels" $\to$ Montrer la localisation, la diversité, les interactions, les enjeux.
    \item \textbf{Choix des symboles (Sémiologie graphique) :}
        \begin{itemize}
            \item \textbf{Surfaces (Zones) :} Hachures / Couleurs pastel pour les \textbf{aires de dominance} (ex. : Zone d'élevage extensif Adamaoua/Nord, Zone de pêche lacustre/fluviale, Bassins artisanaux majeurs).
            \item \textbf{Points (Localisation précise) :} Symboles spécifiques : $\blacktriangle$ (Marché à bétail), $\bullet$ (Centre artisanat majeur / Chefferie), $\ast$ (Port de pêche / Foyer halieutique), $\blacksquare$ (Festival culturel).
            \item \textbf{Lignes (Flux / Réseaux) :} Flèches proportionnelles (Transhumance), Traits discontinus (Routes commerciales artisanat), Flèches bleues (Migration poisson).
        \end{itemize}
    \item \textbf{Construction de la légende (Hiérarchisée en 3 titres) :}
        \begin{enumerate}
            \item \textbf{Les grands ensembles de production} (Surfaces : Élevage, Pêche, Artisanat).
            \item \textbf{Les nœuds et pôles de valorisation} (Points : Marchés, Centres artisanaux, Ports, Festivals).
            \item \textbf{Les dynamiques et vulnérabilités} (Flèches : Transhumance, Flux commerciaux ; Symboles ponctuels : Zones de conflit, Zones menacées climatique, Aires protégées/ZICGC).
        \end{enumerate}
    \item \textbf{Esthétique :} Propreté, netteté, coloriage crayon de couleur (pas feutre), écriture lisible, flèche Nord, échelle, titre complet, sources.
\end{enumerate}
\end{methode}

\begin{popfigure}
\begin{tikzpicture}[scale=0.65, every node/.style={font=\footnotesize}]
% --- BASE CARTOGRAPHIQUE SIMPLIFIÉE CAMEROUN ---
% Contour approximatif (polygone)
\draw[popdark, thick, fill=popcream!30] 
(0,0) -- (1,1.5) -- (0.5,3) -- (2,4.5) -- (4,5) -- (6,4.5) -- (7.5,3.5) -- (8.5,2) -- (9,0.5) -- (7.5,-0.5) -- (6,-1) -- (4,-1.2) -- (2,-0.5) -- (0.5,0.5) -- cycle;
% Fleuves principaux
\draw[popblue, thick] (2, -0.5) .. controls (3, 1) and (4, 2.5) .. (4.5, 4) node[above right, popblue] {Bénoué};
\draw[popblue, thick] (1, 1.5) .. controls (2.5, 2.5) and (3.5, 3.5) .. (4, 4.5) node[above, popblue] {Sanaga};
\draw[popblue, thick] (0.5, 0.5) .. controls (1.5, 1.5) and (2.5, 2) .. (3, 3) node[above left, popblue] {Nyong};
\draw[popblue, thick] (7, 0) .. controls (7.5, 1) .. (7.5, 2) node[right, popblue] {Wouri};
\draw[popblue, thick] (8, 3) .. controls (7.5, 3.5) .. (6.5, 4) node[left, popblue] {Logone};
\draw[popblue, thick] (5.5, 4.5) -- (4.5, 4.5) node[left, popblue] {Lac Tchad};
% Villes repères
\foreach \x/\y/\name in {
    1.5/0.5/Douala, 3.5/2.5/Yaoundé, 5.5/2.5/Bafoussam, 4.5/4.2/Ngaoundéré, 
    7.5/1.5/Maroua, 6.5/3.5/Garoua, 2.5/4.5/Bertoua, 1.5/3.5/Ebolowa, 
    8.5/3.5/Kousseri, 4.5/4.8/LacTchad} {
    \fill[popdark] (\x,\y) circle (2.5pt);
    \node[anchor=south west, font=\tiny] at (\x+0.1,\y+0.1) {\name};
}
% --- COUCHES THÉMATIQUES (CROQUIS ÉLÈVE) ---
% 1. SURFACES : AIRES DE DOMINANCE (Hachures / Patterns)
% Élevage : Adamaoua, Nord, Extrême-Nord (zone sahélienne)
\pattern[pattern=north east lines, pattern color=popgreen!60!popdark] 
(3.5,3.5) -- (5,4.5) -- (7,4.5) -- (8,3.5) -- (7.5,2.5) -- (5.5,3) -- cycle;
\node[popgreen!60!popdark, rotate=30, font=\bfseries\small] at (6, 4) {ZONE ÉLEVAGE EXTENSIF};
% Pêche : Lac Tchad, Logone, Bénoué, Sanaga maritime, Wouri, Côte
\pattern[pattern=north west lines, pattern color=popblue!60!popdark] 
(7.5, 3.5) -- (8.5, 3.5) -- (9, 2) -- (8, 1) -- (7.5, 1.5) -- cycle; % Lac Tchad / Logone
\pattern[pattern=north west lines, pattern color=popblue!60!popdark] 
(3.5, 4) -- (5, 4.5) -- (5.5, 3.5) -- (4, 3) -- cycle; % Bénoué moyenne
\pattern[pattern=north west lines, pattern color=popblue!60!popdark] 
(1.5, 0.5) -- (2.5, 1) -- (3, 0.5) -- (2, 0) -- cycle; % Wouri / Estuaire
\node[popblue!60!popdark, rotate=-20, font=\bfseries\small] at (8, 2.5) {ZONE PÊCHE LACUSTRE/FLUVIALE};
% Artisanat : Bassins majeurs (Ouest, Nord-Ouest, Foumban, Grand Sud)
\pattern[pattern=crosshatch dots, pattern color=poporange!70!popdark] 
(4.5, 2) -- (6, 2.5) -- (6.5, 1.5) -- (5, 1) -- cycle; % Ouest (Bafoussam/Baham/Foumban)
\pattern[pattern=crosshatch dots, pattern color=poporange!70!popdark] 
(2.5, 1.5) -- (4, 2) -- (4.5, 1) -- (3, 0.5) -- cycle; % Centre/Sud (Ebolowa/Sangmelima - sculpture)
\pattern[pattern=crosshatch dots, pattern color=poporange!70!popdark] 
(3.5, 4) -- (4.5, 4.5) -- (5, 3.5) -- (4, 3) -- cycle; % Nord-Ouest (Bamenda - Toghu) / Adamaoua (Ngaoundéré - cuir)
\node[poporange!70!popdark, font=\bfseries\small] at (5.5, 1.8) {BASSINS ARTISANAUX MAJEURS};
% 2. POINTS : NŒUDS DE VALORISATION
% Marchés à bétail (Triangle noir)
\foreach \x/\y in {6.5/3.5, 7.5/1.5, 4.5/4.2, 8.5/3.5} {
    \node[draw=popdark, fill=popgreen!20, regular polygon, regular polygon sides=3, inner sep=2pt] at (\x,\y) {};
}
% Centres artisanaux / Chefferies (Carré orange)
\foreach \x/\y in {5.5/2.5, 4.8/2.2, 3.5/4, 2.5/3.5} {
    \node[draw=popdark, fill=poporange!20, minimum size=5pt] at (\x,\y) {};
}
% Ports / Foyers halieutiques (Étoile bleue)
\foreach \x/\y in {1.5/0.5, 7.5/3.5, 5/4} {
    \node[draw=popblue, fill=popblue!10, star, star points=5, star point ratio=2.25, inner sep=1.5pt] at (\x,\y) {};
}
% Festivals (Losange violet)
\foreach \x/\y in {1.5/0.5, 4.5/4.2, 3.5/4, 6.5/3.5} {
    \node[draw=poppurple, fill=poppurple!10, diamond, inner sep=2pt] at (\x,\y) {};
}
% 3. FLUX / DYNAMIQUES
% Transhumance (Flèches vertes épaisses)
\draw[popgreen!60!popdark, very thick, ->, >=Stealth, decorate, decoration={snake, amplitude=1mm, segment length=4mm}] (7, 4) -- (5, 2.5) node[midway, above left, font=\tiny] {Transhumance Nord $\to$ Sud};
\draw[popgreen!60!popdark, very thick, ->, >=Stealth, decorate, decoration={snake, amplitude=1mm, segment length=4mm}] (5, 2.5) -- (7, 4) node[midway, below right, font=\tiny] {Retour};
% Flux artisanat/commercial (Traits pointillés orange)
\draw[poporange, dashed, thick] (5.5, 2.5) -- (3.5, 2.5) node[midway, above, font=\tiny] {Flux produits artisanaux};
\draw[poporange, dashed, thick] (5.5, 2.5) -- (1.5, 0.5) node[midway, below, font=\tiny] {Export portuaire};
% Zones de conflit / Vulnérabilité (Cercles pointillés rouges)
\draw[poppink, thick, dashed] (7.5, 2.5) circle (0.8) node[above, font=\tiny, poppink] {Conflits Agri/Elev};
\draw[poppink, thick, dashed] (8.5, 3) circle (0.6) node[below, font=\tiny, poppink] {Sécheresse / Lac Tchad};
\draw[poppink, thick, dashed] (4.5, 3.5) circle (0.7) node[left, font=\tiny, poppink] {Pression foncière};
% --- LÉGENDE INTÉGRÉE (Simulée pour le rendu final) ---
\begin{scope}[shift={(10, 0)}, font=\tiny]
\node[font=\bfseries, align=left] at (0, 5.5) {LÉGENDE HIÉRARCHISÉE};
% Titre 1
\node[font=\bfseries, anchor=west] at (0, 5) {1. GRANDS ENSEMBLES DE PRODUCTION};
\fill[pattern=north east lines, pattern color=popgreen!60!popdark] (0, 4.6) rectangle (0.8, 4.9); \node[anchor=west] at (1, 4.75) {Élevage extensif (Adamaoua, Nord, Extrême-Nord)};
\fill[pattern=north west lines, pattern color=popblue!60!popdark] (0, 4.2) rectangle (0.8, 4.5); \node[anchor=west] at (1, 4.35) {Pêche lacustre \& fluviale (Lac Tchad, Bénoué, Sanaga, Wouri, Côte)};
\fill[pattern=crosshatch dots, pattern color=poporange!70!popdark] (0, 3.8) rectangle (0.8, 4.1); \node[anchor=west] at (1, 3.95) {Bassins artisanaux majeurs (Ouest, N-O, Sud, Adamaoua)};
% Titre 2
\node[font=\bfseries, anchor=west] at (0, 3.4) {2. NŒUDS ET PÔLES DE VALORISATION};
\node[draw=popdark, fill=popgreen!20, regular polygon, regular polygon sides=3, inner sep=2pt] at (0.4, 3.0) {}; \node[anchor=west] at (1, 3.0) {Marché à bétail majeur};
\node[draw=popdark, fill=poporange!20, minimum size=5pt] at (0.4, 2.6) {}; \node[anchor=west] at (1, 2.6) {Centre artisanat / Chefferie (Ndop, Toghu, Poterie, Sculpture)};
\node[draw=popblue, fill=popblue!10, star, star points=5, star point ratio=2.25, inner sep=1.5pt] at (0.4, 2.2) {}; \node[anchor=west] at (1, 2.2) {Port / Foyer halieutique};
\node[draw=poppurple, fill=poppurple!10, diamond, inner sep=2pt] at (0.4, 1.8) {}; \node[anchor=west] at (1, 1.8) {Festival culturel (Ngondo, Nyem-Nyem, Ngoun, Mpo'o)};
% Titre 3
\node[font=\bfseries, anchor=west] at (0, 1.4) {3. DYNAMIQUES ET VULNÉRABILITÉS};
\draw[popgreen!60!popdark, very thick, ->, >=Stealth, decorate, decoration={snake, amplitude=1mm, segment length=4mm}] (0, 1.0) -- (1.5, 1.0) node[anchor=west, midway] {Couloirs de transhumance (aller/retour)};
\draw[poporange, dashed, thick] (0, 0.6) -- (1.5, 0.6) node[anchor=west, midway] {Flux commerciaux artisanat / produits halieutiques};
\draw[poppink, thick, dashed] (0.4, 0.2) circle (0.2) node[anchor=west] at (1, 0.2) {Zones de conflits d'usage / Vulnérabilité climatique forte};
% Nord / Échelle
\node at (0, -0.5) {\textbf{N}}; \draw[->, thick] (0, -0.8) -- (0, -0.2);
\node at (0, -1.2) {Échelle : 1 cm $\approx$ 150 km};
\node at (0, -1.5) {Source : MINESEC, MINADER, MINEPIA, MINPMEESA, ONMT, OAPI};
\end{scope}
\end{tikzpicture}
\legende{Croquis de synthèse guidé : « Espaces productifs traditionnels du Cameroun ». Intégration des trois secteurs (surfaces), des nœuds de valorisation (points) et des dynamiques/vulnérabilités (flux, zones de tension). Légende hiérarchisée en 3 niveaux. Base cartographique simplifiée.}
\end{popfigure}

\courssub{Synthèse transversale : Vers un développement durable intégré}

\noindent\textbf{Intuition finale :} Élevage, Pêche, Artisanat partagent le même territoire, les mêmes acteurs (les ruraux), le même patrimoine culturel et la même vulnérabilité climatique. Les traiter séparément conduit à des politiques sectorielles inefficaces (ex. : barrage pour riziculture qui tue la pêche et bloque la transhumance). Le \cle{développement durable intégré} impose une gouvernance transversale.

\begin{definition}[Développement durable des pratiques traditionnelles]
C'est la \textbf{satisfaction des besoins présents} (revenus décents, identité culturelle, sécurité alimentaire, cohésion sociale) \textbf{sans compromettre la capacité des générations futures} à disposer de \textbf{ressources naturelles renouvelées} (pâturages, stocks de poisson, forêts, argile, eau), de \textbf{savoirs techniques vivants} (transmis, adaptés, protégés) et d'\textbf{écosystèmes fonctionnels} (cycles hydrologiques, fertilité sols, biodiversité).
\end{definition}

\begin{aretenir}[Les 4 piliers de l'action publique et privée pour l'intégration]
\begin{enumerate}
    \item \textbf{Gouvernance foncière unifiée :} Cadastre pastoral + Zonage pêche + Réserves matières premières artisanat $\to$ Plan d'occupation des sols (POS) communal / régional concerté.
    \item \textbf{Adaptation climatique "basée sur les écosystèmes" (EbA) :} Restauration berges (frayères + pâturage + bois énergie), Agroforesterie parc à karité/nére (fourrage + fruit + teinture + sol), Gestion intégrée des ressources en eau (GIRE) bassins versants.
    \item \textbf{Éducation \& Transmission rénovée :} Intégration modules "Savoirs locaux" programmes scolaires (primaire/secondaire technique), Centres de formation artisanale (CFA) dual (entreprise/école), Valorisation statut maître-artisan / éleveur-modèle.
    \item \textbf{Économie de la connaissance \& PI :} Labellisation (IG OAPI, Label "Artisanat Cameroun"), Commerce équitable, Plateformes numériques gérées par coopératives, Tourisme communautaire équitable (revenus réinvestis dans transmission/environnement).
\end{enumerate}
\end{aretenir}

\begin{methode}[Pour la dissertation / argumentation au Bac : Sujet type "Les pratiques agropastorales et artisanales traditionnelles face aux défis du développement durable au Cameroun"]
Suivez cette structure en 5 temps (garantie de rigueur) :
\begin{enumerate}
    \item \textbf{Définir \& Délimiter :} Définir "pratiques traditionnelles" (élevage, pêche, artisanat), "développement durable" (3 piliers + culture), espace camerounais. Annoncer le paradoxe : vitalité démographique/économique vs vulnérabilité structurelle.
    \item \textbf{Constater la vitalité ET la mutation :} Chiffres clés (emplois, PIB informel, sécurité alimentaire, exportations artisanat). Montrer l'hybridation (moteurs, pirogues, fours améliorés, téléphonie). \emph{Ne pas figer dans le passé.}
    \item \textbf{Analyser les moteurs (Internes / Externes) :} 
        \begin{itemize}
            \item \textit{Internes :} Démographie, érosion savoirs, conflits générationnels, gestion coutumière.
            \item \textit{Externes :} Changement climatique (variabilité), Pression foncière (agro-industrie, urbain, conservation), Marchés mondiaux (prix matières premières, contrefaçons), Politiques publiques (sectorielles vs intégrées).
        \end{itemize}
    \item \textbf{Évaluer les impacts (Grille 3 piliers + Culture) :}
        \begin{itemize}
            \item \textit{Environnement :} Dégradation sols/eau/biodiversité OU gestion durable (ex. pâturage tournant, zones de non-prise).
            \item \textit{Social :} Inégalités genre/âge, conflits, exode OU cohésion, identité, inclusion.
            \item \textit{Économie :} Précarité, faible VA, dépendance OU revenus diversifiés, valeur ajoutée locale, insertion marchés niche.
            \item \textit{Culturel :} Perte mémoire/rites OU patrimonialisation vivante (IG, festivals).
        \end{itemize}
    \item \textbf{Proposer des pistes de valorisation durable (Opérationnelles) :} 
        \begin{itemize}
            \item Institutionnel : Application Codes (Pastoral, Pêche, Artisanat), Réforme foncière, OAPI/IG.
            \item Économique : GIC/Coopératives, Circuits courts, Label, Financement vert/climat (Fonds Vert, Adaptation).
            \item Technique : Hybridation accompagnée (R\&D participative : IRAD, Universités, Artisans), Énergies renouvelables (solaire séchage poisson, biogaz cuisson poterie).
            \item Humain : Formation duale, Éducation au patrimoine, Tourisme responsable charte communautaire.
        \end{itemize}
    \item \textbf{Conclusion ouverte :} Le "traditionnel" n'est pas un héritage à museifier, mais un \textbf{capital d'innovation frugale} pour l'avenir du Cameroun (SND30, ODD). L'enjeu : passer de la \textbf{subissance} à la \textbf{résilience choisie}.
\end{enumerate}
\end{methode}

\begin{exercice}[Entraînement Bac - Sujet de synthèse]
\textbf{Sujet :} « Les pratiques agropastorales et artisanales traditionnelles au Cameroun : entre résilience et vulnérabilité. Discutez. »
\textbf{Consignes :} Mobiliser les connaissances des sections 1 à 5. Structurer la réponse selon la méthode en 5 temps ci-dessus. Temps conseillé : 2h. Attendus : Croquis de localisation (obligatoire), données chiffrées précises, référence aux textes de loi (2011/022, Code Pastoral, OAPI), articulation changement climatique / foncier / transmission.
\end{exercice}

\begin{corrige}[Exercice n°1 - Éléments de corrigé type]
\textbf{Introduction :} Accroche (ex. citation ODD 12 / Proverbe). Définitions. Paradoxe. Annonce du plan (3 parties : I. Vitalité d'un patrimoine productif vivant ; II. Vulnérabilités structurelles croisées ; III. Leviers pour une résilience durable choisie).
\textbf{I. Vitalité :} Démographie (60\% pop active), PIB informel (~30\% PIB non pétrolier), Sécurité alimentaire (lait, poisson, céréales), Identité (Ndop, Toghu, chefferies). Hybridation : pirogue motorisée, four amélioré, mobile money transhumance.
\textbf{II. Vulnérabilités :} 1. Foncière (Conflits agri/élev, accaparement, barrage SEMRY). 2. Climatique (Sécheresse Nord -20\% pluie, Inondations plaines, Perte biodiversité). 3. Savoirs (Rupture transmission, déscolarisation métiers, perte rites). 4. Économique (Faible VA, intermédiaires, contrefaçons, précarité sociale).
\textbf{III. Leviers :} 1. Juridique (Loi 2011/022, Cadre pastoral, OAPI/IG). 2. Organisation (GIC, Coopératives, Comités gestion pêche, Commissions pastorales). 3. Technique/Éco (EbA, Agroforesterie, Énergies propres, R\&D participative). 4. Marché (Circuits courts, Labels, Tourisme responsable charte, E-commerce coopératif). 5. Humain (Formation duale, Éducation patrimoine, Statut maître-artisan).
\textbf{Conclusion :} Bilan. Ouverture : Nécessité d'une "Politique nationale intégrée du patrimoine productif traditionnel" alignée SND30/ODD.
\end{corrige}

\newpage

\coursec{Activité d'intégration}

\coursec{Les pratiques agropastorales et artisanales traditionnelles}

\courssub{Objectifs de l'activité}
\begin{itemize}
    \item Mobiliser les connaissances sur les systèmes agropastoraux et artisanaux traditionnels (élevage, pêche, artisanat) pour analyser une situation territoriale complexe.
    \item Exploiter un dossier documentaire hétérogène (cartes diachroniques, statistiques, photographies, textes juridiques, témoignages) pour identifier les dynamiques spatiales et les conflits d'usage.
    \item Réaliser un croquis de synthèse cartographiant les conflits et les potentialités de la vallée de la Bénoué.
    \item Rédiger une note d'argumentation structurée (environ 300 mots) proposant des mesures d'aménagement concerté et durable, justifiées par les documents.
\end{itemize}

\courssub{Situation : Aménager la vallée de la Bénoué, un territoire sous tension}

La vallée de la Bénoué (région du Nord, Cameroun) constitue un espace de convergence majeur entre le fleuve Bénoué, ses affluents (Mayo Rey, Faro, Déo) et la plaine alluviale. Ce territoire, large de 20 à 40 km sur plus de 300 km de long, concentre des enjeux forts : il abrite le parc national de la Bénoué (180 000 ha, créé en 1968, classé réserve de biosphère UNESCO en 1981), des zones de chasse villageoises (ZIC GC 1, 2, 3, 4), des périmètres irrigués (riziculture de Lagdo, diversification SODECOTON), et des zones d'habitat dense (Garoua, \num{436 000} hab. en 2020 ; Lagdo, Guider).

\textbf{Contexte biophysique et démographique.}\par
Le climat est de type soudano-sahélien (pluviométrie 900--1 100 mm/an, saison sèche 6--7 mois). Les sols sont des vertisols (argiles gonflantes) en plaine et des sols ferrugineux lessivés sur les glacis. La population du département de la Bénoué est passée de \num{1 200 000} hab. (RGPH 2005) à \num{1 850 000} hab. (estimations 2020), soit une croissance annuelle moyenne de 3,2 \%. Cette pression démographique entraîne une extension des surfaces cultivées (mil, sorgho, maïs, coton, riz pluvial) au détriment des parcours et des zones humides.

\textbf{Données documentaires fournies (extrait du dossier élève).}

\textbf{Document 1 : Carte d'occupation du sol -- Diachronie 2000 vs 2020 (simplifiée).}
\begin{center}
\begin{tabularx}{\linewidth}{|l|X|X|}\hline
\textbf{Unité paysagère} & \textbf{2000 (\% surface)} & \textbf{2020 (\% surface)} \\ \hline
Savane arborée / parcours extensifs & 48 \% & 28 \% \\
Zones humides / bancs de sable (pêche) & 12 \% & 7 \% \\
Agriculture pluviale (mil/sorgho/coton) & 25 \% & 42 \% \\
Agriculture irriguée (riz Lagdo) & 5 \% & 8 \% \\
Habitat / infrastructures & 3 \% & 6 \% \\
Forêts galeries / zones protégées (PN Bénoué) & 7 \% & 9 \% \\ \hline
\end{tabularx}
\end{center}
\emph{Source : Traitement d'images Landsat / Sentinel, adaptation MINEPAT 2021.}

\textbf{Document 2 : Évolution des effectifs d'élevage et captures de pêche (Département de la Bénoué).}
\begin{center}
\begin{tabular}{|l|cccc|}\hline
\textbf{Indicateur} & \textbf{2000} & \textbf{2010} & \textbf{2020} & \textbf{Variation 2000-2020} \\ \hline
Cheptel bovin (têtes) & 420 000 & 580 000 & 750 000 & + 78,6 \% \\
Cheptel ovin/caprin (têtes) & 310 000 & 450 000 & 620 000 & + 100 \% \\
Transhumants entrants (têtes/an) & 150 000 & 220 000 & 300 000 & + 100 \% \\
Captures pêche artisanale (tonnes/an) & 12 500 & 9 800 & 6 200 & - 50,4 \% \\
Nombre de pirogues motorisées & 120 & 450 & 1 200 & $\times$ 10 \\ \hline
\end{tabular}
\end{center}
\emph{Source : DREPIA Nord, Rapports annuels ; Enquête cadre MINEPIA 2020.}

\textbf{Document 3 : Photographies légendées (observations de terrain 2022).}
\begin{itemize}
    \item \textit{Photo A :} Berges du Bénoué à l'aval du barrage de Lagdo. Érosion sévère, cultures de bord de rive (riz, patate douce) jusqu'au fil de l'eau. Absence de bande riveraine.
    \item \textit{Photo B :} Point d'eau pastoral (mare aménagée) envahi par des champs de coton. Troupeaux contraints de traverser les champs pour s'abreuver (janvier 2022, saison sèche).
    \item \textit{Photo C :} Atelier de poterie traditionnelle à Figuil (argile de la plaine alluviale). Four à bois alimenté par bois mort prélevé en zone protégée périphérique. Stocks d'argile menacés par l'extension des rizières irriguées (périmètre de Lagdo).
    \item \textit{Photo D :} Marché hebdomadaire de Garoua. Vente de cuir tanné (artisanat), poisson fumé (pêche), bétail sur pied. Présence de collecteurs pour export vers le Nigeria (bétail) et Douala (poisson).
\end{itemize}

\textbf{Document 4 : Extraits de textes réglementaires et témoignages.}
\begin{itemize}
    \item \textit{Loi n° 98/005 du 14 avril 1998 (régime de l'élevage) :} Art. 12 -- « Les couloirs de transhumance et les aires de repos sont inaliénables. » Art. 15 -- « Tout cultivateur qui détruit des cultures par son bétail en dehors des couloirs est passible d'amende. »
    \item \textit{Arrêté préfectoral 2015/045 :} Interdiction de la pêche au filet maillant < 25 mm et de la motorisation des pirogues > 15 CV sur le Bénoué amont. \textbf{Non appliqué} (manque de moyens de surveillance).
    \item \textit{Témoignage Ardo (chef pastoral) de Rey Bouba :} « Avant, nous négocions le "kola" (droit de passage) avec les lamibé. Aujourd'hui, les champs bouchent les couloirs officiels. Nos bêtes meurent de soif ou mangent le mil. Les jeunes éleveurs achètent des motos pour surveiller les troupeaux, ça change la donne. »
    \item \textit{Témoignage Présidente Groupement de pêcheuses de Garoua :} « Le poisson devient rare et petit. Les filets "moustiquaires" (mailles 10 mm) prennent tout, même les alevins. L'eau baisse vite depuis le barrage. Nous devons aller plus loin, mais l'essence coûte cher. »
    \item \textit{Témoignage Artisan potier Figuil :} « L'argile de qualité est sous les rizières de la SEMRY. On nous refuse l'accès. Le bois pour les fours se fait rare, les gardes du parc nous pourchassent. »
\end{itemize}

\courssub{Consignes}
\begin{enumerate}
    \item \textbf{Analyse cartographique et statistique (Croquis).} À partir des Documents 1 et 2, réaliser un croquis légendé et titré (format A4 paysage) de la vallée de la Bénoué mettant en évidence :
    \begin{itemize}
        \item Les principales unités spatiales (fleuve, plaine alluviale, piémont, zone protégée, ville).
        \item Les dynamiques d'occupation du sol 2000 $\rightarrow$ 2020 (flèches d'expansion agricole, régression parcours/zones humides).
        \item Les zones de conflits d'usage (minimum 4 zones distinctes : agriculture/élevage, agriculture/pêche, élevage/pêche, artisanat/ressources).
        \item Les infrastructures structurantes (barrage Lagdo, route nationale N1, périmètres rizicoles de Lagdo, couloirs de transhumance théoriques).
    \end{itemize}
    \item \textbf{Analyse des transformations des pratiques (Rédaction courte).} En vous appuyant sur l'ensemble du dossier, expliquez en 150 mots comment la motorisation (motos éleveurs, pirogues pêcheurs), la sédentarisation des éleveurs et la commercialisation de l'artisanat ont modifié les rapports aux ressources naturelles et aux espaces sur 20 ans.
    \item \textbf{Note d'argumentation : Proposition d'aménagement concerté (300 mots $\pm$ 10\%).} En votre qualité de consultant géographe mandaté par le Préfet de la Bénoué, rédigez une note proposant \textbf{trois mesures concrètes et articulées} d'aménagement durable (ex: zonage participatif, sécurisation couloirs/aires de repos, gestion communautaire ressources halieutiques/argile/bois, label artisanat, concertation plateforme multi-acteurs). Chaque mesure doit être justifiée par au moins une référence explicite aux documents du dossier (chiffres, photos, textes, témoignages).
\end{enumerate}

\courssub{Ressources à mobiliser}
\begin{aretenir}[Notions clés de la leçon pour l'activité]
\begin{enumerate}
    \item \textbf{Systèmes d'élevage :} Élevage nomade (transhumance Nord-Sud), semi-sédentaire (éleveurs Peuls sédentarisés), ranch (SODECOTON). Notions de \cle{couloir de transhumance}, \cle{aire de repos}, \cle{point d'eau pastoral}, \cle{charge pastorale}.
    \item \textbf{Pêche traditionnelle :} Pêche de subsistance vs commerciale. Engins : nasses, filets maillants, lignes. Notions de \cle{surpêche}, \cle{juvéniles}, \cle{période de repos biologique}, \cle{motorisation artisanale}.
    \item \textbf{Artisanat traditionnel :} Poterie (Figuil, argile alluviale), tannerie (cuir, peaux), vannerie, forge. Matières premières : argile, bois (énergie), peaux (sous-produit élevage). Enjeux : \cle{gestion durable ressources}, \cle{valorisation circuits courts/export}, \cle{concurrence ressources (bois/argile)}.
    \item \textbf{Conflits d'usage :} Agriculture/Élevage (divagation, couloirs bouchés), Élevage/Pêche (piétinement frayères, pollution), Agriculture/Pêche (assèchement zones humides, pesticides), Artisanat/Environnement (déforestation énergie, extraction argile).
    \item \textbf{Outils d'aménagement :} \cle{Plan d'occupation des sols (POS)}, \cle{Schémas directeurs}, \cle{Gestion communautaire des ressources naturelles (GCRN)}, \cle{Approche paysage}, \cle{Concertation multi-acteurs (plateformes)}.
    \item \textbf{Méthode croquis :} Fond de carte simplifié $\rightarrow$ Zonage (couleurs/hachures) $\rightarrow$ Figé ponctuels (villes, points d'eau) $\rightarrow$ Figé linéaires (fleuve, routes, couloirs, flèches flux) $\rightarrow$ Légende organisée (thèmes) $\rightarrow$ Titre, échelle, orientation, source.
\end{enumerate}
\end{aretenir}

\courssub{Démarche et corrigé commenté}
\begin{exoresolu}[Activité d'intégration : Aménagement de la vallée de la Bénoué]
\textbf{Énoncé :} À partir d'un dossier documentaire sur la vallée de la Bénoué (Nord-Cameroun), l'élève doit : 1) Réaliser un croquis de synthèse des dynamiques spatiales et conflits d'usage (2000-2020). 2) Rédiger une analyse des transformations des pratiques (150 mots). 3) Produire une note d'argumentation préconisant trois mesures d'aménagement concerté (300 mots).

\tcblower

\textbf{1. Croquis de synthèse : « La vallée de la Bénoué : dynamiques et conflits d'usage (2000-2020) »}

\begin{popfigure}
\begin{tikzpicture}[scale=0.85, every node/.style={font=\small}]
    % Fond : Fleuve Bénoué (sinueux vertical)
    \draw[popblue, line width=3pt, decorate, decoration={snake, amplitude=1.5pt, segment length=15pt}] (5,0.5) .. controls (5.2,2) and (4.8,4) .. (5,5.5) .. controls (5.3,7) and (4.7,9) .. (5,11.5);
    % Affluents
    \draw[popblue, line width=1.5pt] (2,3) .. controls (3.5,3.5) .. (5,4); % Mayo Rey
    \draw[popblue, line width=1.5pt] (8,6) .. controls (6.5,6.5) .. (5,7); % Faro
    \draw[popblue, line width=1.5pt] (2,9) .. controls (3.5,9.5) .. (5,10); % Déo
    % Barrage Lagdo
    \draw[popdark, line width=2.5pt] (5,8.5) -- (5,8.8) node[midway, right=2pt] {\tiny Barrage Lagdo};
    % Zones : Plaine alluviale (Ouest fleuve) - Agriculture expansive
    \fill[poporange!30, pattern=north east lines, pattern color=poporange] (1,1) rectangle (5,11);
    \node[align=center, poporange] at (3, 5.5) {Plaine alluviale\\ \textbf{Expansion agricole}\\ (Mil, Coton, Riz Lagdo)};
    % Zones : Piémont/Plateau (Est fleuve) - Parcours / Parc
    \fill[popgreen!20, pattern=dots, pattern color=popgreen] (5,1) rectangle (11,11);
    \node[align=center, popgreen] at (8, 5.5) {Piémont / Plateau\\ \textbf{Parcours / Parc Bénoué}};
    % Parc National Bénoué (zone protégée Est)
    \draw[popdark, line width=1.5pt, dashed] (6,7) rectangle (10.5,11);
    \node[popdark] at (8.25, 9) {PN Bénoué};
    % Villes
    \node[circle, fill=popdark, inner sep=2pt, label=below:\textbf{Garoua}] at (3, 2) {};
    \node[circle, fill=popdark, inner sep=2pt, label=below:\textbf{Lagdo}] at (5, 8.5) {};
    \node[circle, fill=popdark, inner sep=2pt, label=right:\textbf{Rey Bouba}] at (9, 9) {};
    \node[circle, fill=popdark, inner sep=2pt, label=right:\textbf{Figuil}] at (7, 3.5) {};
    % Route N1
    \draw[popline, line width=2pt, double=popdark, double distance=1pt] (3,0.5) -- (3,11.5) node[above] {N1};
    % Couloirs de transhumance (théoriques vs bouchés)
    \draw[->, poppurple, line width=2pt, dashed] (9, 10.5) .. controls (7, 8) and (6, 5) .. (4.5, 2.5) node[below left, poppurple] {Couloir Nord-Sud (bouché)};
    \draw[->, poppurple, line width=1.5pt] (9, 10.5) -- (9, 11.5) node[above] {Entrée transhumants};
    % Zones de conflits (symboles éclairs)
    \node[poppink, scale=1.5] at (4.5, 4) {$\otimes$}; % Agri/Elevage
    \node[poppink, scale=1.5] at (4.8, 6) {$\otimes$}; % Peche/Agri (berges)
    \node[poppink, scale=1.5] at (7, 8) {$\otimes$}; % Elevage/Parc (braconnage/paturage)
    \node[poppink, scale=1.5] at (6.5, 3.5) {$\otimes$}; % Artisanat/Argile (Figuil)
    % Flèches flux commerciaux
    \draw[->, popgold, line width=1.5pt] (3, 2) -- (1, 2) node[left] {Bétail $\rightarrow$ Nigeria};
    \draw[->, popgold, line width=1.5pt] (3, 2) -- (3, 0) node[below] {Poisson/Artisanat $\rightarrow$ Douala};
    % Légende (simplifiée dans le code, l'élève la fait à part)
    \node[draw, fill=popcream, inner sep=5pt, anchor=south west, align=left] at (0.2, 0.2) {
        \textbf{LÉGENDE SIMPLIFIÉE}\\
        \textcolor{poporange}{\rule{4pt}{4pt}} Expansion agricole (2000-2020)\\
        \textcolor{popgreen}{\rule{4pt}{4pt}} Parcours / Zone pastorale\\
        \textcolor{popdark}{\rule{4pt}{4pt}} Zone protégée (PN Bénoué)\\
        \textcolor{popblue}{\rule{4pt}{4pt}} Fleuve Bénoué \& affluents\\
        \textcolor{poppurple}{\rule{4pt}{4pt}} Couloir transhumance (théorique)\\
        $\otimes$ \textcolor{poppink}{Conflits d'usage identifiés}\\
        $\bullet$ \textcolor{popdark}{Villes / Infrastructures}\\
        \textcolor{popgold}{\rule{4pt}{4pt}} Flux commerciaux sortants
    };
\end{tikzpicture}
\legende{Croquis de synthèse : Dynamiques spatiales et conflits d'usage dans la vallée de la Bénoué (2000-2020).}
\end{popfigure}

\textbf{Justification du croquis :}
\begin{itemize}
    \item Le fond de carte oppose la plaine alluviale Ouest (agricole) au piémont Est (pastoral/protégé), structuré par le fleuve Bénoué (axe Nord-Sud).
    \item L'expansion agricole (hachures orange) grignote les parcours (verts) et les zones humides (bleu), confirmant Doc 1 (Agri pluviale +17\%, Parcours -20\%, Zones humides -5\%).
    \item Le couloir de transhumance (flèche violette pointillée) traverse la zone agricole : c'est le point de friction principal (Doc 3 Photo B, Témoignage Ardo).
    \item Les 4 conflits ($\otimes$) sont localisés : 1) Agri/Élevage (couloir), 2) Agri/Pêche (berges cultivées Doc 3 Photo A), 3) Élevage/Parc (pression sur ressources protégées), 4) Artisanat/Agri (argile Figuil sous rizières Doc 3 Photo C).
    \item Les flux commerciaux (or) montrent l'insertion dans les réseaux longue distance (Nigeria, Douala), moteur de la commercialisation (Doc 3 Photo D).
\end{itemize}

\textbf{2. Transformations des pratiques sur 20 ans (Rédaction type ~150 mots)}

Entre 2000 et 2020, la motorisation a profondément reconfiguré l'exploitation des ressources. Chez les éleveurs, l'usage de la moto (Doc 2 : cheptel +78\%, transhumants $\times$2 ; Doc 4 Témoignage Ardo) permet une surveillance accrue des troupeaux sur de vastes parcours, mais facilite aussi l'intrusion rapide dans les champs et les zones protégées, exacerbant les conflits fonciers. La sédentarisation partielle des Peuls (création de campements fixes près des points d'eau) transforme l'élevage nomade en système agro-sylvo-pastoral, augmentant la pression locale sur le fourrage et l'eau. Chez les pêcheurs, l'explosion du nombre de pirogues motorisées ($\times$10 en 20 ans, Doc 2) a décuplé l'effort de pêche et le rayon d'action, expliquant l'effondrement des captures (-50\%) malgré l'augmentation des moyens, et la capture de juvéniles (filets moustiquaires, Doc 4). Pour l'artisanat, la commercialisation (Doc 3 Photo D : export cuir, poisson fumé) monétise les productions traditionnelles (poterie Figuil, tannerie), créant une demande accrue en matières premières (argile, bois, peaux) qui entre en concurrence directe avec l'agriculture (argile sous rizières) et la conservation (bois dans le parc). Ces trois dynamiques convergent vers une \cle{privatisation de l'accès aux ressources communes} au détriment de la régulation coutumière.

\textbf{3. Note d'argumentation : Proposition d'aménagement concerté de la vallée de la Bénoué (300 mots)}

\textit{Objet : Trois mesures pour une gouvernance partagée de la vallée de la Bénoué.}

Monsieur le Préfet,

Le diagnostic croisé des documents révèle un territoire en rupture : l'expansion agricole (+17\% de surfaces en 20 ans, Doc 1) a fragmenté les couloirs de transhumance (Doc 3, Photo B) et grignoté les zones de frayères et d'argile (Doc 3, Photos A \& C), tandis que la motorisation non régulée (pirogues $\times$10, Doc 2) a divisé par deux la ressource halieutique. Pour concilier sécurité alimentaire, paix sociale et préservation du capital naturel, je propose trois mesures articulées.

\textbf{Mesure 1 : Instituer un « Plan de Gestion Intégrée des Terroirs » (PGIT) par arrondissement, matérialisé par un zonage participatif contraignant.} Ce plan, élaboré en concertation avec les Lamibé, Ardo, groupements de pêcheurs, artisans et services techniques (MINEPAT, MINEPIA, MINADER), délimitera sur le terrain (bornage GPS) : des \textbf{couloirs de transhumance élargis à 500 m} (Loi 98/005 Art. 12) avec \textbf{aires de repos tous les 15 km} équipées en points d'eau (réponse au témoignage Ardo : « bêtes meurent de soif ») ; des \textbf{bandes riveraines inconstructibles de 100 m} le long du Bénoué et affluents pour protéger les frayères et limiter l'érosion (Doc 3 Photo A) ; des \textbf{zones d'extraction artisanale (argile Figuil, bois mort)} négociées avec les gestionnaires des périmètres rizicoles (SODECOTON/État) et la conservation du Parc. Ce zonage, annexé au POS communal, rendra opposable le droit de passage et l'accès aux ressources.

\textbf{Mesure 2 : Mettre en place une « Gestion Communautaire de la Pêche et des Zones Humides » (GCPZH) avec quota et repos biologique.} Sur le modèle des ZIC villageoises, des comités de gestion par groupement de pêcheurs (Doc 4 Témoignage Pêcheuses) fixeront des \textbf{quotas annuels par espèce}, imposeront un \textbf{repos biologique de 3 mois} (juin-août, montaison) et \textbf{interdiront les moteurs > 15 CV et mailles < 25 mm} (Arrêté 2015 appliqué). La surveillance sera assurée par des \textbf{brigades mixtes} (pêcheurs, garde-pêche, marine) financées par une redevance sur les captures commerciales (export Douala/Nigeria, Doc 3 Photo D). Cela inverse la logique de « course au poisson » motorisée (Doc 2 : effort $\uparrow$, captures $\downarrow$).

\textbf{Mesure 3 : Créer un « Label « Savoir-Faire Bénoué » pour l'artisanat, adossé à un fonds de gestion durable des matières premières.} Ce label (poterie Figuil, cuir Garoua, vannerie) garantira la traçabilité et la qualité pour accéder aux marchés urbains et export (Doc 3 Photo D). En contrepartie, une \textbf{redevance « matière première »} (1\% du CA) alimentera un fonds géré par les groupements d'artisans pour : (a) sécuriser l'accès aux gisements d'argile (convention avec les aménagements hydro-agricoles pour zones tampons), (b) financer des \textbf{fours à gaz/biomasse améliorés} remplaçant le bois mort du Parc (Doc 3 Photo C), (c) former aux nouvelles techniques (design, finition). Cela transforme l'artisanat en levier de conservation et non plus en prédateur de ressources.

Ces trois mesures, indissociables, reposent sur la \cle{concertation multi-acteurs} (plateforme préfectorale trimestrielle) et la \cle{responsabilisation des usagers} par la délégation de gestion. Elles transforment les conflits d'usage en contrats territoriaux, condition \textit{sine qua non} de la durabilité de la vallée.

\textit{Fait à Garoua, le [Date]. Le Consultant Géographe.}
\end{exoresolu}

\courssub{Bilan de l'activité}
\begin{aretenir}[Ce que l'activité a permis de mettre en évidence]
\begin{enumerate}
    \item \textbf{Interdépendance des systèmes :} L'activité montre que l'élevage, la pêche, l'agriculture et l'artisanat ne sont pas des secteurs étanches mais partagent les mêmes ressources finies (eau, sol, bois, espace, argile) dans la vallée. Une action sur l'un (ex: extension rizicole Lagdo) impacte immédiatement les autres (perte argile potiers, bouchage couloirs, assèchement frayères).
    \item \textbf{Rôle accélérateur de la motorisation/commercialisation :} La modernisation technique (motos, moteurs hors-bord) et l'insertion marchande (export bétail/poisson, label artisanat) ne sont pas neutres : elles augmentent la pression d'exploitation (\cle{surpêche}, \cle{surpâturage}, \cle{déforestation énergie}) et affaiblissent les modes de régulation coutumiers (kola, repos biologique respecté).
    \item \textbf{Limites de l'approche répressive / sectorielle :} Les textes existent (Loi 98, Arrêté 2015) mais sont inappliqués faute de moyens, de légitimité locale et de concertation. L'aménagement durable ne peut être que \textbf{contractuel} (PGIT, GCPZH, Label) et \textbf{territorialisé} (zonage au GPS, bornage physique).
    \item \textbf{Compétences géographiques mobilisées :} Analyse diachronique (cartes 2000/2020), lecture de tableaux statistiques (taux de variation, parts), interprétation de photographies (indices visuels de conflits), analyse de textes juridiques et sociologiques (témoignages), synthèse cartographique (croquis), rédaction argumentée à visée opérationnelle (note préfectorale).
    \item \textbf{Transférabilité :} La démarche (diagnostic partagé $\rightarrow$ zonage contractuel $\rightarrow$ gestion déléguée $\rightarrow$ valorisation économique) est réplicable dans d'autres bassins versants camerounais (Logone, Sanaga, Lac Tchad) confrontés aux mêmes tensions agropastorales.
\end{enumerate}
\end{aretenir}

\newpage

\coursec{Méthodes et exercices résolus}

\coursec{Les pratiques agropastorales et artisanales traditionnelles}

\courssub{L'élevage traditionnel : diversité des systèmes et organisation de l'espace}

\begin{methode}[Analyser un système d'élevage traditionnel et son organisation de l'espace]
\textbf{Objectif :} Caractériser un système (nomadisme, transhumance, élevage sédentaire) et expliquer ses logiques spatiales (parcours, points d'eau, cultures).\\
\textbf{Démarche :}
\begin{enumerate}
    \item \textbf{Identifier le type d'élevage :} Nomadisme (déplacement constant), Transhumance (aller-retour saisonnier Nord/Sud ou altitude), Élevage sédentaire/ranching (fixe).
    \item \textbf{Repérer les ressources clés :} Eau (mares, fleuves, forages), Pâturages (herbacées, ligneuses), Sel/minéraux.
    \item \textbf{Dessiner le croquis d'organisation spatiale :} Tracer les axes de transhumance (flèches), localiser les zones de séjour sec/humide, les points d'eau, les marchés, les zones de cultures (conflits potentiels).
    \item \textbf{Expliquer l'adaptation au milieu :} Lien climat (saisonnalité des pluies) $\leftrightarrow$ mobilité du troupeau $\leftrightarrow$ gestion de la ressource herbagère.
    \item \textbf{Conclure sur la viabilité :} Pression démographique, extension des cultures, changement climatique $\rightarrow$ réduction des parcours.
\end{enumerate}
\textbf{Réflexe Bac :} Sur un croquis, la légende doit être hiérarchisée (Flèches = flux/transhumance ; Surfaces = zones de pâture ; Points = points d'eau/marchés). Citer des exemples camerounais (Bororo/Fulbe au Nord, Mbororo à l'Adamaoua, Bakweri au Sud-Ouest).
\end{methode}

\begin{exoresolu}[Analyse de la transhumance des Mbororo à l'Adamaoua (Type Bac)]
\textbf{Énoncé :} À partir de vos connaissances et du document ci-après (extrait fictif de rapport), analysez l'organisation spatiale de la transhumance des éleveurs Mbororo dans la région de l'Adamaoua et expliquez les conflits d'usage qui en découlent.
\begin{quote}
\textit{« En saison sèche (nov.-mars), les troupeaux quittent les hauts plateaux de l'Adamaoua (zones de séjour des pluies) pour descendre vers les plaines de la Bénoué et le bassin du Lac Tchad (zones de séjour de la sèche) à la recherche d'eau et de pâturages verts. Le retour s'effectue avec les premières pluies. Les couloirs de passage traversent des zones de cultures cotonnières et vivrières en expansion. »}
\end{quote}
\textbf{Consignes :}
\begin{enumerate}
    \item Réalisez un croquis schématique légendé de cette transhumance (échelle non exigée).
    \item Expliquez les deux causes principales des conflits agriculteurs-éleveurs dans cette zone.
    \item Proposez deux mesures de gestion durable de l'espace pastoral.
\end{enumerate}
\tcblower
\textbf{Solution détaillée :}

\textbf{1. Croquis schématique :}
\begin{popfigure}
\begin{tikzpicture}[scale=0.9]
% Fond
\fill[popcream] (-5,-3) rectangle (5,3);
% Zones
\fill[popgreenL!50] (-5,1) rectangle (5,3); % Zone humide (Nord/Adamaoua)
\fill[popgoldL!50] (-5,-3) rectangle (5,-1); % Zone sèche (Sud/Bénoué)
\fill[poporangeL!30] (-5,-1) rectangle (5,1); % Zone de transition/cultures

% Flèches Transhumance
\draw[popblue, very thick, ->] (0,2.5) -- (0,1.2) node[midway, right] {\small Saison sèche (Descente)};
\draw[poporange, very thick, ->] (0.5,-2.5) -- (0.5,-1.2) node[midway, right] {\small Saison des pluies (Montée)};

% Points d'eau
\foreach \y in {-2, 0, 2} {
    \fill[popblue] (3,\y) circle (3pt);
    \node[right] at (3.2,\y) {\small Point d'eau / Mare};
}
% Marché
\fill[poppurple] (-3,0) circle (4pt);
\node[left] at (-3.2,0) {\small Marché à bétail};

% Légendes zones
\node[align=center, font=\bfseries\small] at (0,2) {ZONE DE SÉJOUR DES PLUIES\\(Adamaoua : pâturages abondants)};
\node[align=center, font=\bfseries\small] at (0,-2) {ZONE DE SÉJOUR DE LA SÈCHE\\(Plaines Bénoué / Lac Tchad : eau permanente)};
\node[align=center, font=\bfseries\small, text=popdark] at (0,0) {ZONE DE TRANSITION / CULTURES\\(Coton, Maïs, Mil - Conflits)};

% Nord
\node[popblue, font=\bfseries] at (-4.5, 2.5) {N $\uparrow$};
\end{tikzpicture}
\legende{Croquis schématique : Transhumance Mbororo Adamaoua / Plaines de la Bénoué. Flèches bleues = descente (sèche) ; Flèches orange = montée (pluies).}
\end{popfigure}

\textbf{2. Causes des conflits agriculteurs-éleveurs :}
\begin{itemize}
    \item \textbf{L'empiétement des cultures sur les couloirs de transhumance et les pâturages :} L'extension des périmètres cotonniers (SODECOTON) et vivriers réduit la largeur des couloirs (souvent < 50m au lieu de 100-300m légaux) et détruit les points d'eau. Les troupeaux traversent les champs $\rightarrow$ dégâts cultureaux.
    \item \textbf{La non-application de la réglementation foncière (Loi 74-1 du 6 juil. 1974) :} L'absence de délimitation physique effective des couloirs et aires de pâture, la corruption de certains chefs traditionnels vendant des terres pastorales à des agriculteurs, et le non-respect du calendrier cultural (récoltes tardives bloquant le passage) exacerbent les tensions.
\end{itemize}

\textbf{3. Mesures de gestion durable :}
\begin{enumerate}
    \item \textbf{La matérialisation et la sécurisation foncière des couloirs de transhumance :} Bornage physique, plantation d'arbres marqueurs, inscription au registre foncier (titres collectifs pour les groupements d'éleveurs).
    \item \textbf{La mise en place de Commissions Locales de Gestion des Conflits (CLGC) paritaires :} Intégrant éleveurs, agriculteurs, administration (Sous-préfet, Délégué MINEPIA/MINADER) et chefferies traditionnelles pour la médiation préventive (calendrier de passage) et curative (indemnisation rapide via fonds de garantie).
\end{enumerate}

\textbf{Conclusion :} L'organisation spatiale de la transhumance Mbororo est une adaptation rationnelle à la saisonnalité climatique (ressource herbe/eau). Sa durabilité est menacée par l'anthropisation croissante de l'espace (agriculture, démographie), nécessitant une gouvernance foncière inclusive et contraignante.
\end{exoresolu}

\courssub{Transformations de l'élevage traditionnel : modernisation, conflits et politiques publiques}

\begin{methode}[Expliquer les transformations de l'élevage traditionnel : modernisation et politiques publiques]
\textbf{Objectif :} Analyser le passage de l'extensif pur à des formes intensifiées (ranching, amélioration génétique, alimentation) et le rôle de l'État.\\
\textbf{Démarche :}
\begin{enumerate}
    \item \textbf{Caractériser l'élevage « traditionnel » de référence :} Extensif, itinérant, race locale (Zébu Goudali, Azawak, Kuri), productivité faible (lait/viande), capital sur pied (épargne, prestige), main-d'œuvre familiale.
    \item \textbf{Identifier les leviers de modernisation (politiques publiques) :}
        \begin{itemize}
            \item \textbf{Sanitaire :} Campagnes de vaccination (peste bovine, CBPP, fièvre aphteuse), parcs de vaccination, couloirs sanitaires.
            \item \textbf{Génétique :} Centres d'insémination artificielle (CIA), importation de taureaux améliorateurs (Gir, Sahiwal), croisement F1.
            \item \textbf{Fourragère / Alimentation :} Aménagement de pâturages (réserves fourragères, ensilage, tourteaux de coton/arachide), blocs multi-nutritionnels.
            \item \textbf{Infrastructures :} Pistes pastorales, marchés à bétail modernes, abattoirs frigorifiques, laiteries (ex: LAITERIE DU BÉNOUÉ).
            \item \textbf{Foncier / Sédentarisation :} Ranchs d'État (ex: Ranch de Ngaoundaba, Koundini), périmètres d'élevage villageois.
        \end{itemize}
    \item \textbf{Analyser les freins/limites :} Coût des intrants, perte de rusticité des races améliorées, accaparement des ranchs par élites, résistance culturelle (refus de vendre le « capital »), changement climatique (sécheresses récurrentes).
    \item \textbf{Typologier les exploitations actuelles :} Traditionnel pur, Traditionnel amélioré (sanitaire + alimentation), Semi-intensif (pâturage + appoint), Intensif/Peri-urbain (zéro pâturage).
\end{enumerate}
\textbf{Formule clé :} Taux d'accroissement du cheptel $\approx$ Taux de natalité $-$ Taux de mortalité $-$ Taux de réforme/vente. La modernisation vise à faire baisser la mortalité et augmenter la natalité (intervalle vêlage-vêlage raccourci).
\end{methode}

\begin{exoresolu}[Évaluation d'un programme de modernisation : Le cas du Ranch de Ngaoundaba]
\textbf{Énoncé :} Le Ranch de Ngaoundaba (Adamaoua, créé en 1948, $\approx$ 10 000 ha) a vocation de sélection de race Goudali et de diffusion de taureaux améliorateurs. Malgré les investissements publics, son impact sur l'élevage villageois environnant reste limité.
\textbf{Consignes :}
\begin{enumerate}
    \item Présentez les trois piliers techniques de la modernisation mise en œuvre dans ce ranch.
    \item Expliquez pourquoi la diffusion des taureaux améliorateurs vers les éleveurs traditionnels (Mbororo) échoue souvent. Donnez 3 raisons socio-économiques et techniques.
    \item Calculez le taux de réforme annuel nécessaire pour maintenir un cheptel stable de 500 têtes si le taux de mortalité est de 4\% et le taux de natalité de 55\% (taux de vêlage sur femelles reproductrices). On suppose 60\% de femelles dans le cheptel et un sex-ratio à la naissance 50/50.
\end{enumerate}
\tcblower
\textbf{Solution détaillée :}

\textbf{1. Trois piliers techniques du Ranch de Ngaoundaba :}
\begin{itemize}
    \item \textbf{Amélioration génétique :} Sélection intra-race Goudali (performance laitière/viande), station d'insémination artificielle (IA), production de taureaux reproducteurs certifiés.
    \item \textbf{Maîtrise de l'alimentation :} Aménagement de pâturages (Brachiaria, Panicum), conservation de fourrage (foin, ensilage), distribution d'aliments concentrés (tourteau coton, son) en saison sèche critique.
    \item \textbf{Prophylaxie rigoureuse :} Calendrier vaccinal strict (Peste bovine éradiquée, CBPP, Fièvre aphteuse, Pasteurellose), lutte contre les trypanosomes (tsé-tsé), traitements vermifuges réguliers.
\end{itemize}

\textbf{2. Raisons de l'échec de la diffusion vers le traditionnel :}
\begin{enumerate}
    \item \textbf{Inadaptation du génotype amélioré au milieu extensif :} Les taureaux F1 ou races exotiques (Gir, Holstein) ont des besoins nutritionnels élevés et une sensibilité accrue aux maladies (trypanosomose, tiques) que le système villageois (parcours naturel, absence d'appoint) ne peut couvrir $\rightarrow$ mortalité forte ou infertilité.
    \item \textbf{Coût d'acquisition et de maintenance prohibitif :} Prix d'un taureau améliorateur (500 000 -- 1 500 000 FCFA) inaccessible pour l'éleveur moyen. Nécessité d'enclos, d'eau permanente, de sel/minéraux, de soins vétérinaires payants.
    \item \textbf{Raison culturelle et économique (Logique du capital sur pied) :} L'éleveur traditionnel privilégie le \textbf{nombre} (sécurité, prestige, dot, assurance vie) à la \textbf{productivité individuelle}. Vendre une vache pour acheter un taureau « performant » mais fragile réduit le capital social. La gestion du troupeau est collective (famille/clan), la décision d'IA ou d'achat taureau est individuelle $\rightarrow$ blocage social.
\end{enumerate}

\textbf{3. Calcul du taux de réforme (R) pour cheptel stable ($\Delta N = 0$) :}
\textbf{Données :} $N_{total} = 500$ ; $T_{mort} = 4\%$ ; $T_{natal} = 55\%$ (sur femelles reproductrices) ; $\% Femelles = 60\%$ ; Sex-ratio naissance = 1.
\textbf{Étape 1 : Calcul des naissances annuelles (B).}
Nombre de femelles $= 500 \times 0,60 = 300$ têtes.
Nombre de femelles reproductrices (estimation standard $\approx$ 70\% des femelles) $= 300 \times 0,70 = 210$ vaches.
Naissances $B = 210 \times 0,55 = 115,5 \approx 116$ veaux/an.
\textbf{Étape 2 : Calcul des décès annuels (D).}
$D = 500 \times 0,04 = 20$ têtes/an.
\textbf{Étape 3 : Équilibre démographique.}
Variation nette $\Delta N = B - D - R$ (Réformes/Ventes).
Pour $\Delta N = 0$ : $R = B - D$.
$R = 116 - 20 = 96$ têtes/an.
\textbf{Étape 4 : Taux de réforme (\% du cheptel total).}
$T_{réforme} = \frac{R}{N_{total}} \times 100 = \frac{96}{500} \times 100 = \mathbf{19,2\%}$.
\textbf{Interprétation :} Pour stabiliser le cheptel, l'éleveur doit vendre/réformer près de 20\% de son troupeau par an. C'est élevé (souvent 10-15\% en traditionnel), ce qui explique la croissance naturelle des cheptels traditionnels si pas de vente contrainte.
\end{exoresolu}

\courssub{La pêche traditionnelle : techniques, espaces et mutations}

\begin{methode}[Caractériser la pêche traditionnelle : techniques, espaces, mutations]
\textbf{Objectif :} Distinguer les types de pêche (continentale, maritime, lagunaire), leurs engins, leurs cibles et les pressions subies.\\
\textbf{Démarche :}
\begin{enumerate}
    \item \textbf{Classifier par milieu :}
        \begin{itemize}
            \item \textbf{Continentale :} Fleuves (Sanaga, Bénoué, Logone, Chari), Lacs (Tchad, Ossa, Barombi), Barrages (Lagdo, Mbakaou, Memve'ele). Engins : Filets maillants, nasses, lignes, barrage (akwala), empoisonnement (interdit).
            \item \textbf{Lagunaire / Estuarienne :} Wouri, Moungo, Rio del Rey, Ntem. Engins : Senne tournante, filets fixes (akwala), palangre, pêche à pied (coquillages, crabes).
            \item \textbf{Maritime artisanale :} Côte (Kribi, Limbé, Idenau, Campo). Pirogues à moteur hors-bord (40-100 CV). Engins : Filets maillants dérivants, senne de plage, palangre, casiers (langoustes).
        \end{itemize}
    \item \textbf{Analyser l'organisation sociale :} Pêcheurs autochtones vs migrants (Ghanéens, Béninois, Nigérians, Tchadiens), rôle des « patrons » (propriétaires pirogues/moteurs/filets) vs équipiers (partage de la pêche), organisations professionnelles (GIC, Coopératives).
    \item \textbf{Identifier les mutations (Pressions) :}
        \begin{itemize}
            \item \textbf{Surexploitation :} Mailles trop petites (juvéniles), effort de pêche croissant (nombre pirogues, puissance moteurs).
            \item \textbf{Dégradation habitats :} Pollution (hydrocarbures, déchets urbains, pesticides coton), destruction mangroves (bois de chauffe, carciniculture), érosion côtière, barrages (coupure migration poisson).
            \item \textbf{Concurrence déloyale :} Pêche industrielle (chalutiers) dans la zone réservée aux artisans (3 milles nautiques), pêche IUU (Illicite, Non déclarée, Non réglementée).
            \item \textbf{Post-récolte :} Perte de qualité (manque glace/chaîne de froid), transformation artisanale (fumage, séchage) énergivore (bois) et polluante (HAP).
        \end{itemize}
\end{enumerate}
\textbf{Données clés Bac :} Production halieutique Cameroun $\approx$ 250 000 -- 300 000 t/an (dont $\sim$ 70-80\% artisanale). Déficit couverture besoins $\approx$ 200 000 t/an (importations). Emploi direct $\sim$ 200 000 pêcheurs + 400 000 transformation/commercialisation.
\end{methode}

\begin{exoresolu}[Diagnostic de la pêche artisanale sur le Lac Tchad (Partie camerounaise)]
\textbf{Énoncé :} La partie camerounaise du Lac Tchad (Logone-et-Chari) concentre une pêche artisanale intense. Les captures sont dominées par \textit{Alestes baremoze} (Makayab), \textit{Hydrocynus} (Captain), \textit{Clarias} (Poisson-chat). Les filets maillants de maille 25-30 mm dominent. Le nombre de pirogues a triplé en 20 ans.
\textbf{Consignes :}
\begin{enumerate}
    \item Citez deux engins de pêche traditionnels actifs et un engin passif utilisés sur le Lac.
    \item Expliquez l'impact écologique de l'utilisation généralisée de filets à mailles de 25 mm sur le stock de \textit{Alestes baremoze} (taille première maturité sexuelle $\approx$ 18-20 cm).
    \item Proposez trois mesures de gestion concertée (co-gestion) adaptées au contexte transfrontalier (Cameroun, Tchad, Niger, Nigeria).
\end{enumerate}
\tcblower
\textbf{Solution détaillée :}

\textbf{1. Engins de pêche Lac Tchad :}
\begin{itemize}
    \item \textbf{Actifs :} \textbf{La senne de plage (ou senne tournante « Boukarou »)} : encerclement de bancs, main-d'œuvre nombreuse. \textbf{La pêche au barrage / Akwala (filet fixe tendu sur pieux)} : capture passive mais installation active, très efficace en période de crue/décrue.
    \item \textbf{Passif :} \textbf{Le filet maillant (Gillnet) dérivant ou calé} : le plus utilisé aujourd'hui. \textbf{La palangre (ligne mère avec hameçons)} : ciblant Captain/Clarias. \textbf{Les nasses/casiers} : pour Clarias/écrevisses.
\end{itemize}

\textbf{2. Impact des mailles 25 mm sur \textit{Alestes baremoze} :}
\begin{itemize}
    \item \textbf{Pêche de juvéniles (Growth Overfishing) :} La maille 25 mm capture des individus de $\approx$ 12-16 cm (facteur de forme $\approx$ 0,5-0,6 $\times$ périmètre maille). Or la première maturité sexuelle ($L_{50}$) est à 18-20 cm.
    \item \textbf{Conséquence :} Les poissons sont capturés \textbf{avant d'avoir reproduit au moins une fois}. Cela réduit le stock reproducteur (Biomasse Reproductrice / Spawning Stock Biomass), diminue le recrutement futur et fait chuter les captures par unité d'effort (CPUE). C'est une \textbf{surexploitation de recrutement}.
    \item \textbf{Sélection génétique inverse :} Pression de sélection favorisant les individus à croissance lente/maturité précoce à petite taille.
\end{itemize}

\textbf{3. Mesures de co-gestion transfrontalière (CBLT - Commission Bassin Lac Tchad) :}
\begin{enumerate}
    \item \textbf{Harmonisation et application stricte de la maille minimale légale :} Passage à \textbf{40-45 mm} (maille carrée) pour les filets maillants ciblant Alestes/Captain, avec contrôle conjoint aux débarquements (brigades mixtes CM/TD/NG/NI).
    \item \textbf{Instauration de repos biologiques synchronisés :} Arrêt total de la pêche (ex: 2 mois en période de frai maximal -- juin/juillet ou décrue) sur l'ensemble du lac, avec filets de compensation ou activités alternatives (agriculture de décrue, élevage) pour les pêcheurs.
    \item \textbf{Gestion de l'effort de pêche (Licences plafonnées) :} Plafonnement du nombre de pirogues motorisées par village/campement, registre biométrique des pêcheurs, interdiction de nouveaux entrants (gel des permis), promotion de la motorisation propre (4 temps) vs 2 temps polluants.
\end{enumerate}
\textbf{Conclusion :} La survie de la pêche lacustre exige une gouvernance régionale forte (CBLT), car le stock est partagé ; une mesure unilatérale camerounaise est inefficace si les riverains tchadiens/nigérians continuent les petites mailles.
\end{exoresolu}

\courssub{L'artisanat traditionnel : diversité, organisation et évolution}

\begin{methode}[Analyser l'artisanat traditionnel : diversité, organisation, évolution]
\textbf{Objectif :} Classer les métiers, comprendre l'organisation de la production (atelier, apprentissage) et les enjeux de modernisation (marché, tourisme, technologie).\\
\textbf{Démarche :}
\begin{enumerate}
    \item \textbf{Classer par matière première / filière :}
        \begin{itemize}
            \item \textbf{Bois / Sculpture :} Masques, statuaires, sièges (Bamileke, Bamoun, Tikar, Bassa), menuiserie d'art.
            \item \textbf{Terre / Céramique / Poterie :} Foumban (Bamoun), Baham, Batcham, Mokolo (argile, kaolin). Utilitaire (canaris, marmites) et artistique.
            \item \textbf{Métallurgie / Fer / Bronze :} Forgerons traditionnels (outils agricoles, armes, bijoux), Fonte à la cire perdue (Bamoun, Tikar - pipes, bijoux, statuettes).
            \item \textbf{Textile / Vannerie / Cuir :} Tissus ndop, tentes (Nord), nattes, paniers (Raphia, rotin, bambou), maroquinerie (Nord, Maroua).
            \item \textbf{Autres :} Peinture sur toile/batik, instruments de musique (balafon, tam-tam), pharmacopée (conditionnement).
        \end{itemize}
    \item \textbf{Analyser l'organisation socio-économique :}
        \begin{itemize}
            \item \textbf{Apprentissage :} Système « maître-apprenti » (contrat oral, durée 3-7 ans, « liberté » finale), transmission familiale (castes parfois : forgerons, potiers).
            \item \textbf{Atelier :} Unité de production familiale/collective, souvent au domicile, faible capitalisation, outils manuels.
            \item \textbf{Commercialisation :} Vente directe atelier, marchés hebdomadaires, boutiques artisanat (Yaoundé, Douala, Foumban), export informel, salons (SIAO Ouaga, PROMOTE Yaoundé).
        \end{itemize}
    \item \textbf{Identifier les mutations / défis :}
        \begin{itemize}
            \item \textbf{Adaptation au marché :} Standardisation, design contemporain, nouveaux matériaux (recyclage métal/pneu, plastique), certification (label « Made in Cameroon »).
            \item \textbf{Tourisme / Patrimoine :} Musées (Foumban, Baham, Bandjoun), villages artisans (Foumban, Baham), risque de « folklorisation » / contrefaçon (importation Asie).
            \item \textbf{Innovation technique :} Tours à bois électriques, fours à gaz/électriques (céramique), machines à coudre, outillage portatif $\rightarrow$ gain productivité, perte « main » ?
            \item \textbf{Formation/Insertion :} Centres de formation professionnelle (CFP), lycées techniques (spécialités Arts Appliqués), reconnaissance diplômes (CAP, BT, BTS Design/Artisanat).
        \end{itemize}
\end{enumerate}
\textbf{Notion clé :} \cle{Économie de la débrouille} vs \cle{Entrepreneuriat artisanal}. Le passage à l'échelle nécessite : Accès financement (FNE, PIAASI), Protection propriété intellectuelle (OAPI - dessins/modèles), Regroupement en clusters/pôles d'excellence.
\end{methode}

\begin{exoresolu}[Étude de cas : La poterie de Foumban face aux défis de la modernité]
\textbf{Énoncé :} Foumban (Région Ouest, Dpt Noun) est le pôle céramique majeur du Cameroun. L'argile (kaolin) est extraite localement. La production est dominée par des femmes (modelage) et des hommes (cuisson, four). Les fours traditionnels (foyers ouverts, bois) sont énergivores et polluants. La demande touristique et urbaine (décoration) croît.
\textbf{Consignes :}
\begin{enumerate}
    \item Réalisez un organigramme simplifié de la chaîne opératoire (Chaîne Opératoire) de la poterie foumbanaise, de l'extraction à la vente.
    \item Calculez le rendement énergétique approximatif d'un four traditionnel vs un four amélioré type « four à tiroir » ou « four à gaz » si : Four traditionnel : 2 stères bois (1 stère $\approx$ 0,7 t, PCI $\approx$ 15 MJ/kg) pour cuire 500 pièces (pertes 60\%). Four amélioré : 1 stère bois pour 500 pièces (pertes 20\%).
    \item Argumentez pourquoi l'introduction du tour de potier électrique (vitesse, régularité) peut être rejetée par les potentières traditionnelles.
\end{enumerate}
\tcblower
\textbf{Solution détaillée :}

\textbf{1. Chaîne Opératoire (Organigramme) :}
\begin{popfigure}
\begin{tikzpicture}[node distance=1.2cm and 1.5cm, every node/.style={font=\small, align=center}]
\node[draw, popblueL, fill=popblue!10, rounded corners, align=center] (ext){Extraction Argile\\ (Carrière kaolin)};
\node[draw, popgreenL, fill=popgreen!10, rounded corners, right=of ext, align=center] (prep){Préparation\\ (Broyage, Trempage,\\ Décantation, Malaxage)};
\node[draw, poporangeL, fill=poporange!10, rounded corners, right=of prep, align=center] (form){Façonnage / Modelage\\ (Tournage main, Colombins,\\ Estampage moules) -- \textbf{Femmes}};
\node[draw, poppurpleL, fill=poppurple!10, rounded corners, right=of form, align=center] (sech){Séchage lent\\ (Ombre $\rightarrow$ Soleil)};
\node[draw, poppinkL, fill=poppink!10, rounded corners, below=of sech, align=center] (cuis){Cuisson\\ (Four traditionnel bois OU\\ Four amélioré) -- \textbf{Hommes}};
\node[draw, popgoldL, fill=popgold!10, rounded corners, right=of cuis, align=center] (fin){Finition / Décor\\ (Engobes, Polissage,\\ Peinture, Vernis)};
\node[draw, popblueL, fill=popblue!10, rounded corners, right=of fin, align=center] (vent){Commercialisation\\ (Atelier, Marché, Salon,\\ Export, Musée)};

\draw[->, thick, popdark] (ext) -- (prep);
\draw[->, thick, popdark] (prep) -- (form);
\draw[->, thick, popdark] (form) -- (sech);
\draw[->, thick, popdark] (sech) -- (cuis);
\draw[->, thick, popdark] (cuis) -- (fin);
\draw[->, thick, popdark] (fin) -- (vent);
\end{tikzpicture}
\legende{Chaîne opératoire simplifiée de la poterie de Foumban. Division sexuelle du travail marquée (Formage = Femmes ; Cuisson = Hommes).}
\end{popfigure}

\textbf{2. Calcul rendement énergétique (Énergie utile / Énergie consommée) :}
\textbf{Hypothèses :} 1 stère bois = 0,7 tonne = 700 kg. PCI bois $\approx$ 15 MJ/kg.
Énergie contenue dans 1 stère $= 700 \times 15 = 10\,500$ MJ.
\textbf{Four Traditionnel :} Conso = 2 stères $= 21\,000$ MJ. Pertes 60\% $\rightarrow$ Rendement $\eta_{trad} = 40\%$.
Énergie utile (pièces) $= 21\,000 \times 0,40 = 8\,400$ MJ pour 500 pièces.
\textbf{Four Amélioré :} Pour la \textbf{même cuisson} (énergie utile requise identique = 8 400 MJ), avec un rendement $\eta_{am} = 80\%$ (pertes 20\%).
Énergie consommée nécessaire $= \frac{8\,400}{0,80} = 10\,500$ MJ = \textbf{1 stère}.
\textbf{Analyse :} L'énergie consommée est \textbf{divisée par 2} (2 stères $\rightarrow$ 1 stère) pour la même production. \textbf{Économie de bois : 50\%} (1 stère économisé). Réduction massive déforestation, fumées (santé femmes), coût d'achat bois.

\textbf{3. Raisons du rejet potentiel du tour électrique par les potentières :}
\begin{enumerate}
    \item \textbf{Perte de l'identité technique et du « geste » ancestral :} Le modelage main/colombins/tournage lent est un savoir-faire corporel transmis mère-fille. Le tour électrique standardise, déshumanise, transforme l'artisane en opératrice machine.
    \item \textbf{Contrainte infrastructurelle et économique :} Nécessite électricité stable (coupures fréquentes Foumban), investissement achat/maintenance tour (coût élevé), atelier fermé (perte espace de vie/socialisation).
    \item \textbf{Adaptation du produit au marché local vs touristique :} Le tour impose des formes symétriques, lisses (style « industriel/usine »). La clientèle locale/traditionnelle et une partie touristique recherchent l'irrégularité, la trace du doigt, l'authenticité du « fait main » qui justifie le prix. Le tour produit du « standard » qui entre en concurrence avec l'importé bon marché.
\end{enumerate}
\end{exoresolu}

\courssub{Synthèse : Enjeux transversaux et développement durable des pratiques traditionnelles}

\begin{methode}[Synthèse : Évaluer la durabilité des pratiques agropastorales/artisanales (Approche DD)]
\textbf{Objectif :} Croiser les 3 piliers du Développement Durable (Économique, Social, Environnemental) pour juger de la viabilité d'une pratique et proposer des pistes d'appui.\\
\textbf{Démarche (Grille d'analyse DD) :}
\begin{enumerate}
    \item \textbf{Pilier Environnemental :} Impact sur ressources (eau, sol, biodiversité, forêt, climat). \textit{Ex: Surpâturage $\rightarrow$ érosion ; Pêche mailles petites $\rightarrow$ effondrement stock ; Four bois $\rightarrow$ déforestation ; Teinture artisanale $\rightarrow$ pollution eau.}
    \item \textbf{Pilier Social / Culturel :} Emploi (jeunes, femmes), sécurité alimentaire, cohésion sociale, transmission patrimoine, équité genre, santé (pénibilité, fumées). \textit{Ex: Élevage = assurance vie/sociale ; Poterie = identité féminine ; Pêche = migrants/intégration.}
    \item \textbf{Pilier Économique :} Revenus, productivité, valeur ajoutée, accès marchés/financement, résilience aux chocs (prix, climat). \textit{Ex: Faible productivité travail ; Dépendance intermédiaires ; Saisonnalité revenus.}
    \item \textbf{Identifier les \textbf{Trade-offs} (arbitrages) :} \textit{Ex: Modernisation élevage (productivité $\uparrow$) vs Perte biodiversité races locales / Accaparement foncier. Four amélioré (Environnement $\uparrow$) vs Coût investissement (Éco $\downarrow$ pour pauvre).}
    \item \textbf{Formuler des recommandations « Gagnant-Gagnant » (Synergies) :} Agroforesterie (arbres fourragers + bois d'œuvre + carbone), Aquaculture intégrée (bouse $\rightarrow$ engrais étang $\rightarrow$ poisson), Éco-tourisme artisanat (revenu + patrimoine), Labellisation/IGP (Prix $\uparrow$ + Qualité $\uparrow$ + Environnement $\uparrow$).
\end{enumerate}
\textbf{Structure de réponse type Bac (Question synthèse) :}
\begin{quote}
\textit{« Les pratiques traditionnelles sont-elles un frein ou un levier pour le développement durable au Cameroun ? »}
\underline{Intro} : Définition pratiques, importance poids PIB/Emploi. Thèse : Ni frein ni levier \emph{per se}, tout dépend de l'accompagnement.
\underline{Corps} : 1. Atouts (Social/Éco : emploi, résilience, savoirs). 2. Limites (Environnement : dégradation ; Éco : faible productivité). 3. Conditions de la transition (Politiques : foncier, formation, financement, innovation, gouvernance).
\underline{Conclu} : Nécessité d'une \cle{modernisation endogène} (respect savoirs + apports science/marché) pour transformer le « traditionnel » en « durable ».
\end{quote}
\end{methode}

\begin{exoresolu}[Débat : Modernisation endogène vs Rupture technologique dans l'artisanat du bois (Bamiléké)]
\textbf{Énoncé :} Dans les chefferies Bamiléké (Ouest), la sculpture sur bois (portes, trônes, masques) est une activité séculaire. L'introduction de machines à commande numérique (CNC) et de logiciels de CAO est proposée par un programme d'appui pour « booster la productivité et l'export ».
\textbf{Consignes :}
\begin{enumerate}
    \item En mobilisant la grille d'analyse DD (Environnement, Social, Économique), évaluez les impacts positifs et négatifs de l'introduction de la CNC dans les ateliers de sculpture royaux.
    \item Proposez un modèle de transition « modernisation endogène » conciliant innovation technique et préservation du patrimoine immatériel.
\end{enumerate}
\tcblower
\textbf{Solution détaillée :}

\textbf{1. Évaluation DD de l'introduction CNC :}

\begin{center}
\begin{tabularx}{\linewidth}{|l|X|X|}\hline
\textbf{Pilier} & \textbf{Impacts Positifs (Leviers)} & \textbf{Impacts Négatifs (Freins/Risques)} \\ \hline
\textbf{Environnemental} & \textbf{Optimisation matière :} Imbrication CAO $\rightarrow$ taux de chute bois $\downarrow$ (30\% $\rightarrow$ 10\%). Moins d'arbres abattus pour même volume produit. Précision $\rightarrow$ moins d'erreurs/rebuts. & \textbf{Énergie/Empreinte :} Consommation électricité (réseau instable $\rightarrow$ groupes électrogènes diesel). Usure outils (carbure) $\rightarrow$ déchets métalliques. Transport machines lourdes. \\ \hline
\textbf{Social / Culturel} & \textbf{Attractivité jeunes :} Image « high-tech », moins de pénibilité (poussière, bruit manuel), salaires potentiels $\uparrow$. Formation CAO/FAO valorisante (diplômes). & \textbf{Perte du « Savoir-Faire » (Patrimoine Immatériel) :} Déconnexion geste/outil/bois. Disparition apprentissage long (maître/élève). Risque « standardisation mondiale » : perte styles locaux (iconographie royale, symboles secrets). \textbf{Exclusion :} Artisans âgés/analphabètes marginalisés. \\ \hline
\textbf{Économique} & \textbf{Productivité $\uparrow\uparrow$ :} Séries identiques rapides (commandes hôtels, export). Qualité constante (contrôle qualité). Nouveaux marchés (design contemporain, mobilier sur mesure). Réduction coûts main-d'œuvre/unité. & \textbf{Investissement lourd :} Machine (20-50 M FCFA), logiciels, formation, maintenance, électricité stable. Endettement artisans. \textbf{Dépendance :} Fournisseurs extérieurs (pièces, MAJ logicielles). \textbf{Concurrence interne :} Artisans manuels « low-cost » vs CNC « high-cost » $\rightarrow$ dualisme, conflits. \\ \hline
\end{tabularx}
\end{center}

\textbf{2. Modèle de transition « Modernisation Endogène » :}
\begin{enumerate}
    \item \textbf{Numérisation du patrimoine (Banque de données 3D) :} Scanner (photogrammétrie/LiDAR) les pièces maîtresses anciennes (portes royales, masques sacrés) $\rightarrow$ Bibliothèque numérique de motifs/authentiques. La CNC ne « crée » pas, elle \textbf{reproduit/facilite} à partir de cette banque validée par les gardiens de la tradition (Notables, Maîtres sculpteurs).
    \item \textbf{Hybridation des tâches (Co-botique artisanale) :} \textbf{CNC = Tâches lourdes/répétitives/brutes :} Débits, dégrossissage, panneaux plats, motifs géométriques répétitifs. \textbf{Main = Tâches nobles/fines/symboliques :} Finitions, gravures symboliques secrètes, polissage, patine, assemblage, « mise en vie » de l'objet (rituels). L'artisan devient « Chef de projet / Designer / Finisseur ».
    \item \textbf{Gouvernance partagée (Cluster Artisanal) :} Création d'une \textbf{SCOP / GIE « Ateliers Royaux Modernes »} : Mutualisation machine CNC (coût partagé), Bureau d'études design intégré (jeunes formés CAO + Anciens pour iconographie), Label collectif « Sculpture Bamiléké Authentique » (Marque collective OAPI) garantissant \% main + origine bois certifié (FSC/PEFC local).
    \item \textbf{Formation duale :} Cursus Lycée Technique / CFP : « Technicien en Sculpture Numérique \& Patrimoine » : 50\% atelier traditionnel (maîtres) + 50\% Labo FAO/CAO. Diplôme reconnu (BT/BTS).
\end{enumerate}
\textbf{Conclusion :} La rupture technologique pure détruit le capital symbolique (valeur ajoutée principale). La modernisation endogène utilise la technologie pour \textbf{soulager la contrainte matière/répétition} et \textbf{valoriser le geste humain} sur la haute valeur ajoutée (symbole, finition, récit), assurant la transmission vivante du patrimoine.
\end{exoresolu}

\begin{methode}[Réaliser un croquis de synthèse géographique (Pratiques agropastorales/artisanales)]
\textbf{Objectif :} Produire un croquis lisible, légendé, structuré, répondant à une problématique spatiale (localisation, flux, dynamiques).\\
\textbf{Démarche en 5 étapes (Méthode « CROQUIS ») :}
\begin{enumerate}
    \item \textbf{C} - \textbf{C}erner le sujet : Titre précis (ex: « L'élevage et la pêche traditionnels au Cameroun : espaces, flux et enjeux »), Échelle (Pays / Région), Orientation (Nord), Fond de carte (Contour pays, grands fleuves, villes repères, axes routiers/fer).
    \item \textbf{R} - \textbf{R}epérer les informations : Lister les éléments à cartographier par \textbf{type de signe} :
        \begin{itemize}
            \item \textbf{Surfaces (Zones) :} Zones d'élevage (Nord, Adamaoua, Est), Zones de pêche (Lac Tchad, Côte, Barrages), Zones artisanales (Foumban, Baham, Maroua, Bafoussam).
            \item \textbf{Lignes (Flux/Axes) :} Couloirs transhumance (flèches doubles sens), Routes bétail, Pistes pastorales, Export poisson/bois.
            \item \textbf{Points (Localisation) :} Marchés à bétail (Garoua, Ngaoundéré, Douala), Ranchs, Laiteries, Centres artisanaux, Ports de pêche (Kribi, Limbé, Guirvidig).
        \end{itemize}
    \item \textbf{O} - \textbf{O}rganiser la légende : Hiérarchiser par \textbf{thèmes} (couleurs distinctes) et \textbf{ordre d'importance}.
        \begin{itemize}
            \item Thème 1 : Élevage (Verts/Bleus) -- Surfaces pâturages, Flèches transhumance, Points marchés/ranchs.
            \item Thème 2 : Pêche (Bleus/Cyans) -- Surfaces plans d'eau, Points ports/barrages, Flèches commercialisation.
            \item Thème 3 : Artisanat (Ocres/Rouges) -- Points pôles majeurs, Flèches diffusion produits.
            \item Thème 4 : Conflits/Enjeux (Rouge/Noir) -- Zones conflits agro-pastoraux, Zones surexploitation halieutique, Déforestation bois d'œuvre.
        \end{itemize}
    \item \textbf{Q} - \textbf{Q}ualifier le tracé : Traits nets, coloriage régulier (crayons de couleur), flèches proportionnelles (largeur = importance flux), symboles normalisés (losange = mine/ressource, carré = ville, cercle = point).
    \item \textbf{U} - \textbf{U}nifier et \textbf{S} - \textbf{S}oigner : Titre en haut (centré), Légende encadrée (bas droite ou gauche), Auteur/Source/Date. \textbf{Vérifier :} Cohérence fond/carte, lisibilité, réponse au sujet.
\end{enumerate}
\textbf{Astuce Bac :} Préparez 2-3 « fonds de carte mentaux » types (Cameroun entier, Nord-Cameroun, Côte/Forêt) avec les repères fixes (Yaoundé, Douala, Garoua, Ngaoundéré, Maroua, Bertoua, Bafoussam, Kribi, Limbé, Fleuves : Sanaga, Bénoué, Logone, Chari, Nyong, Wouri).
\end{methode}

\begin{exoresolu}[Croquis de synthèse : « Les espaces des ressources naturelles renouvelables exploitées traditionnellement au Cameroun »]
\textbf{Énoncé :} Réalisez un croquis de synthèse au 1/10 000 000ème environ, mettant en évidence les grands espaces de l'élevage, de la pêche et de l'artisanat d'extraction (bois, argile) au Cameroun, ainsi que les principaux flux de commercialisation et zones de tension.
\tcblower
\textbf{Solution détaillée : Description textuelle du croquis attendu (Code TikZ fourni pour visualisation)}

\begin{popfigure}
\begin{tikzpicture}[scale=0.65]
% --- FOND DE CARTE CAMEROUN SIMPLIFIÉ (Polygone approximatif) ---
\draw[popdark, thick] 
(0,0) -- (1,2) -- (2,4) -- (4,5) -- (6,4.5) -- (8,3) -- (9,1) -- (8.5,-1) -- (7,-2) -- (5,-2.5) -- (3,-2) -- (1.5,-1) -- cycle;
% Villes repères
\node[circle, fill=popdark, inner sep=1.5pt, label=above:\tiny Yaoundé] (Y) at (4.5, 1.5) {};
\node[circle, fill=popdark, inner sep=1.5pt, label=below:\tiny Douala] (D) at (5.5, 0.5) {};
\node[circle, fill=popdark, inner sep=1.5pt, label=above:\tiny Garoua] (G) at (2.5, 3.5) {};
\node[circle, fill=popdark, inner sep=1.5pt, label=above:\tiny Ngaoundéré] (N) at (4, 3) {};
\node[circle, fill=popdark, inner sep=1.5pt, label=below:\tiny Maroua] (M) at (1.5, 2.5) {};
\node[circle, fill=popdark, inner sep=1.5pt, label=right:\tiny Bertoua] (B) at (6.5, 2) {};
\node[circle, fill=popdark, inner sep=1.5pt, label=left:\tiny Bafoussam] (Bf) at (3.5, 1) {};
\node[circle, fill=popdark, inner sep=1.5pt, label=below:\tiny Kribi] (K) at (6.5, -0.5) {};
\node[circle, fill=popdark, inner sep=1.5pt, label=below:\tiny Limbé] (L) at (7.5, 0) {};
% Fleuves
\draw[popblue, thick] (4.5,1.5) .. controls (5,0.5) .. (5.5,0.5) node[right, pos=0.5, font=\tiny] {Sanaga};
\draw[popblue, thick] (2.5,3.5) .. controls (3,2) .. (4,1.5) node[right, pos=0.5, font=\tiny] {Bénoué};
\draw[popblue, thick] (1.5,2.5) -- (1,1.5) node[left, font=\tiny] {Logone};
\draw[popblue, thick] (0,0) -- (-0.5,0.5) node[left, font=\tiny] {Chari};
% Lac Tchad (coin)
\fill[popblue!20] (-0.5,-0.5) rectangle (1,0.5);
\node[font=\tiny, popblue] at (0.25,0) {Lac Tchad};

% --- THÈME 1 : ÉLEVAGE (Verts) ---
% Zones pâturages (Hachures)
\foreach \x/\y in {1.5/3, 2.5/3.5, 3.5/3, 4.5/2.5, 3/2.5} {
    \fill[popgreen!30, pattern=north east lines, pattern color=popgreen] (\x-0.3,\y-0.2) rectangle (\x+0.3,\y+0.2);
}
% Transhumance Flèches (Nord-Sud)
\draw[popgreen, very thick, ->] (3,3.5) -- (3,1.5) node[midway, right, font=\tiny] {Transhumance};
\draw[popgreen, very thick, ->] (4,3) -- (4.5,1) node[midway, right, font=\tiny] {Transhumance};
% Marchés bétail
\foreach \pos in {(2.5,3.5), (4,3), (5.5,0.5)} {
    \fill[popgreen!80!popdark] \pos circle (3pt);
}

% --- THÈME 2 : PÊCHE (Bleus) ---
% Lacs/Barrages
\fill[popblue!30] (5.5,2.5) ellipse (0.5 and 0.3); \node[font=\tiny] at (5.5,2.5) {Lagdo};
\fill[popblue!30] (6.5,1.5) ellipse (0.3 and 0.2); \node[font=\tiny] at (6.5,1.5) {Mbakaou};
\fill[popblue!30] (7.5,1) ellipse (0.3 and 0.2); \node[font=\tiny] at (7.5,1) {Memve'ele};
% Ports pêche maritime
\fill[popblue!80!popdark] (K) circle (3pt); \fill[popblue!80!popdark] (L) circle (3pt);
% Flux commercialisation poisson
\draw[popblue, thick, ->, dashed] (0.25,0) -- (4.5,1.5);
\draw[popblue, thick, ->, dashed] (K) -- (Y);
\draw[popblue, thick, ->, dashed] (L) -- (Y);

% --- THÈME 3 : ARTISANAT BOIS/TERRE (Ocres) ---
% Pôles majeurs
\fill[poporange!80!popdark] (Bf) circle (4pt); \node[above, font=\tiny] at (Bf) {Foumban/Baham};
\fill[poporange!80!popdark] (3.5, 0.5) circle (4pt); \node[below, font=\tiny] at (3.5, 0.5) {Bafoussam};
\fill[poporange!80!popdark] (2, 2.5) circle (4pt); \node[left, font=\tiny] at (2, 2.5) {Maroua (Cuir/Vannerie)};
\fill[poporange!80!popdark] (6.5, 2) circle (4pt); \node[right, font=\tiny] at (6.5, 2) {Bertoua (Bois)};
% Flux artisanat
\draw[poporange, thick, ->, dotted] (Bf) -- (Y);
\draw[poporange, thick, ->, dotted] (3.5,0.5) -- (D);
\draw[poporange, thick, ->, dotted] (6.5,2) -- (D);

% --- THÈME 4 : TENSIONS (Rouge) ---
% Zones conflits agro-pastoraux (Nord/Bénoué)
\draw[popred, thick, pattern=north west lines, pattern color=popred] (1.5,2) rectangle (4,3);
\node[font=\tiny, popred, rotate=45] at (2.7,2.5) {Conflits Agro-Pastoraux};
% Surexploitation pêche (Lac Tchad)
\draw[popred, thick] (-0.5,-0.5) rectangle (1,0.5);
% Déforestation bois d'œuvre (Est/Sud)
\draw[popred, thick, pattern=crosshatch dots, pattern color=popred] (6,1.5) rectangle (8.5,3);
\node[font=\tiny, popred] at (7.2,2.2) {Pression Bois d'œuvre};

% Nord
\node[popdark, font=\bfseries\small] at (9, 4.5) {N $\uparrow$};
\node[font=\small] at (9, 4) {\tiny 1 cm $\approx$ 100 km};
\end{tikzpicture}
\legende{Croquis de synthèse : Espaces des ressources renouvelables exploitées traditionnellement au Cameroun. Élevage (Vert : zones pâturages hachurées, flèches transhumance, points marchés), Pêche (Bleu : plans d'eau, ports, flèches flux), Artisanat (Orange : pôles majeurs, flèches diffusion), Tensions (Rouge : hachures zones conflits, surexploitation, déforestation). Échelle indicative 1/10 000 000.}
\end{popfigure}

\textbf{Points de vigilance pour la notation :}
\begin{itemize}
    \item Légende \textbf{hiérarchisée par thèmes} (pas mélange points/surfaces/flèches).
    \item Utilisation correcte des \textbf{signes graphiques} : Surfaces = hachures/couleurs ; Flux = flèches (plein=lourd, pointillé=léger, pointillé-traité=artisanal) ; Points = symboles géométriques distincts.
    \item \textbf{Titre complet} + Échelle indicative + Nord + Source.
    \item Localisation précise des \textbf{3 bassins d'élevage} (Nord, Adamaoua, Est), des \textbf{2 façades maritimes} (Kribi, Limbé) + Lac Tchad + Barrages, des \textbf{3 pôles artisanaux} (Ouest/Bamoun, Nord/Maroua, Est/Bertoua).
\end{itemize}
\end{exoresolu}

\begin{attention}[Les 5 erreurs qui coûtent des points au Bac]
\begin{enumerate}
    \item \textbf{Confondre « Traditionnel » et « Archaïque/Immuable » :} Ne pas écrire « l'élevage traditionnel ne change pas ». \textbf{Correction :} C'est un système dynamique en mutation (sanitaire, alimentation, motorisation, sédentarisation partielle). Utiliser le terme \cle{modernisation endogène} ou \cle{transition agroécologique}.
    \item \textbf{Oublier la dimension spatiale (Géographie !) :} Réponse purement descriptive (liste techniques) sans localisation (Nord/Adamaoua/Est), sans flux (transhumance, commercialisation), sans acteurs (échelle locale/nationale/transfrontalière). \textbf{Réflexe :} Toujours mentalement (ou sur brouillon) situer : « Où ? », « Qui ? », « Vers où ? ».
    \item \textbf{Mélanger Pêche Artisanale et Pêche Industrielle dans l'analyse :} Les enjeux, engins, acteurs, zones (3 milles nautiques), statistiques sont distincts. \textbf{Piège :} Citer des chalutiers pour la pêche artisanale. \textbf{Donnée clé :} Artisanale = Pirogues $\le$ 15-20m, moteurs hors-bord, filets maillants/sennes ; Industrielle = Chalutiers/senneurs > 20m, licences État, export direct.
    \item \textbf{Calculs démographiques/énergétiques approximatifs ou sans unités :} Dans l'exercice résolu 2 (Taux de réforme) ou 4 (Rendement four), l'oubli de l'unité (\%/an, MJ, stères) ou l'erreur sur la base de calcul (cheptel total vs femelles reproductrices) coûte la moitié des points. \textbf{Méthode :} Écrire la formule, identifier la population de référence, calculer étape par étape, conclure avec unité.
    \item \textbf{Croquis « coloriage » sans légende structurée / Signes non conventionnels :} Utiliser des ronds pour des flux, des flèches pour des zones, pas de hiérarchie dans la légende (thèmes mélangés), pas de Nord, pas d'échelle, titre absent. \textbf{Règle d'or :} 1 Thème = 1 Couleur dominante = 1 Section de légende. Surfaces = Hachures/Pois ; Flux = Flèches (épaisseur $\propto$ importance) ; Points = Symboles géométriques distincts.
\end{enumerate}
\end{attention}

\newpage

\coursec{Exercices}

\courssub{Je vérifie mes connaissances}

\begin{exercice}[QCM : notions clés]
Pour chaque question, choisis la bonne réponse.
\begin{enumerate}
\item La \cle{transhumance} est :
\begin{enumerate}
\item la sédentarisation définitive des éleveurs ;
\item le déplacement saisonnier des troupeaux entre zones de pâturage ;
\item l'élevage industriel en bâtiment fermé.
\end{enumerate}
\item Le \cle{ranching} désigne :
\begin{enumerate}
\item un élevage extensif moderne sur de grands espaces clôturés ;
\item une technique de pêche au filet ;
\item un mode de cuisson de la poterie.
\end{enumerate}
\item Une \cle{mareyeuse} est :
\begin{enumerate}
\item une pêcheuse en haute mer ;
\item une commerçante qui achète et revend le poisson frais ;
\item une artisane vannier.
\end{enumerate}
\item Une \cle{indication géographique} (IG) est :
\begin{enumerate}
\item une carte topographique ;
\item un signe qui lie la qualité d'un produit à son origine géographique ;
\item une taxe à l'exportation.
\end{enumerate}
\end{enumerate}
\end{exercice}

\begin{exercice}[Vrai ou faux ? Justifie ta réponse]
Réponds par vrai ou faux, puis justifie en deux lignes maximum.
\begin{enumerate}
\item Les conflits agriculteurs--éleveurs au Cameroun opposent souvent des cultivateurs et des éleveurs peuls (Mbororo) autour de l'accès à la terre et à l'eau.
\item La pêche artisanale camerounaise utilise exclusivement des chalutiers industriels.
\item L'artisanat traditionnel camerounais a totalement disparu face à l'industrie moderne.
\item La modernisation de l'élevage passe notamment par l'amélioration des races, la vaccination et la création de ranches.
\end{enumerate}
\end{exercice}

\begin{exercice}[Définitions]
Définis en deux ou trois lignes chacune des notions suivantes : \cle{chaîne opératoire} ; \cle{pêche artisanale} ; \cle{mareyeuse} ; \cle{artisanat}.
\end{exercice}

\begin{exercice}[Citations et calcul direct]
\begin{enumerate}
\item Cite deux exemples de \cle{conflits agro-pastoraux} au Cameroun (régions concernées).
\item Cite deux \cle{engins de pêche interdits} par la réglementation camerounaise.
\item Cite deux \cle{matières premières artisanales menacées} au Cameroun.
\item En 2000, les captures artisanales étaient de 90 000 tonnes ; en 2023, elles atteignent 135 000 tonnes. Calcule le taux de variation en pourcentage entre 2000 et 2023.
\end{enumerate}
\end{exercice}

\courssub{Je m'entraîne}

\begin{exercice}[L'élevage traditionnel et ses transformations]
\begin{enumerate}
\item Présente les deux grands systèmes d'élevage traditionnel au Cameroun (élevage extensif nomade/transhumant et élevage sédentaire) en précisant pour chacun : les régions concernées, les espèces élevées et l'organisation de l'espace.
\item Explique deux transformations récentes de l'élevage traditionnel (modernisation, sédentarisation, politiques publiques).
\item Dégage une conséquence positive et une conséquence négative de ces transformations.
\end{enumerate}
\end{exercice}

\begin{exercice}[La pêche traditionnelle en mutation]
Au Cameroun, la pêche artisanale maritime (Kribi, Limbé, Douala) et continentale (lac Tchad, Logone, Wouri) fournit l'essentiel du poisson consommé.
\begin{enumerate}
\item Décris deux techniques de pêche artisanale (engins, embarcations, organisation du travail).
\item Montre comment la pêche artisanale s'est transformée (motorisation des pirogues, glace, circuits de commercialisation par les mareyeuses).
\item Explique pourquoi la surexploitation menace les ressources halieutiques.
\end{enumerate}
\end{exercice}

\begin{exercice}[L'évolution de l'artisanat]
\begin{enumerate}
\item Cite quatre filières artisanales traditionnelles du Cameroun et leur zone de production (ex. : poterie, vannerie, sculpture, tissage).
\item Pour l'une de ces filières, décris la \cle{chaîne opératoire} de la matière première au produit fini.
\item Explique deux évolutions récentes de l'artisanat camerounais (marché touristique, exportation, concurrence industrielle, labels).
\end{enumerate}
\end{exercice}

\begin{exercice}[Étude de cas : les conflits agro-pastoraux]
Dans les régions du Nord-Ouest et de l'Adamaoua, les conflits entre agriculteurs et éleveurs Mbororo se sont multipliés : destruction des champs par les troupeaux, accaparement des points d'eau, occupation des couloirs de transhumance.
\begin{enumerate}
\item Identifie deux causes de ces conflits.
\item Propose deux solutions de gestion concertée de l'espace (délimitation des couloirs de transhumance, comités de dialogue, points d'eau séparés).
\item Rédige un court paragraphe argumenté (10 lignes maximum) montrant que la coexistence agriculture--élevage peut être une source de complémentarité.
\end{enumerate}
\end{exercice}

\courssub{Je me prépare au Bac}

\begin{exercice}[Dissertation structurée (type Bac)]
\textbf{Sujet :} À l'aide d'exemples précis pris au Cameroun, montrez comment les pratiques agropastorales et artisanales traditionnelles s'adaptent aux mutations contemporaines.

\textbf{Consignes :} rédige une dissertation structurée comprenant :
\begin{enumerate}
\item une \textbf{introduction} (accroche, définition des termes clés, problématique, annonce du plan) ;
\item deux \textbf{parties équilibrées} : (I) la persistance et la diversité des pratiques traditionnelles (élevage, pêche, artisanat) ; (II) leurs transformations face aux mutations contemporaines (modernisation, marchés, conflits, politiques publiques) ;
\item une \textbf{conclusion} (bilan, ouverture sur le développement durable).
\end{enumerate}
\textbf{Attendus :} exemples localisés (Adamaoua, Extrême-Nord, Foumban, Kribi\ldots), vocabulaire géographique précis, données chiffrées.
\end{exercice}

\begin{exercice}[Croquis légendé et organisé]
\textbf{Sujet :} Réalise un croquis légendé et organisé de la répartition des principales activités agropastorales et artisanales au Cameroun.

\textbf{Consignes :}
\begin{enumerate}
\item Trace le fond de carte (contour simplifié du Cameroun, flèche du nord, titre).
\item Représente au moins \textbf{trois jeux de symboles} : (a) les grandes zones d'élevage (Adamaoua, Nord, Extrême-Nord) ; (b) les zones de pêche artisanale (côte atlantique, lac Tchad, grands fleuves) ; (c) les pôles artisanaux (Foumban, Bafoussam, Maroua\ldots).
\item Construis une \textbf{légende hiérarchisée} (symboles ordonnés du général au particulier) et donne un \textbf{titre} précis au croquis.
\item Rédige un court commentaire (5 lignes) expliquant les principaux faits de localisation mis en évidence.
\end{enumerate}
\end{exercice}

\begin{exercice}[Analyse de documents et cas pratique (type Bac)]
\textbf{Document 1.} Évolution des captures de poisson au Cameroun (en milliers de tonnes) :

\begin{center}
\begin{tabularx}{\linewidth}{|l|X|X|X|X|}
\hline
\rowcolor{popblueL} \textbf{Année} & \textbf{2000} & \textbf{2010} & \textbf{2020} & \textbf{2023} \\
\hline Pêche artisanale & 90 & 110 & 128 & 135 \\
\hline Pêche industrielle & 60 & 75 & 95 & 105 \\
\hline
\end{tabularx}
\end{center}

\textbf{Document 2.} Le code de la pêche camerounais interdit les engins non sélectifs (filets à petites mailles, pêche à la dynamite, monofilaments) et institue des périodes de repos biologique pour permettre la reproduction des espèces.

\textbf{Questions :}
\begin{enumerate}
\item Calcule le taux de variation des captures artisanales et des captures industrielles entre 2000 et 2023. Compare les deux évolutions.
\item À l'aide des documents et de tes connaissances, explique les causes de la \cle{surexploitation} des ressources halieutiques au Cameroun.
\item Propose deux mesures de gestion durable de la pêche, en t'appuyant sur le document 2.
\item \textbf{Cas pratique :} un groupement d'artisans potiers de Foumban souhaite exporter sa production. Élabore un plan d'action en quatre étapes (diagnostic de la matière première, amélioration de la cuisson avec un four amélioré, obtention d'un label de type indication géographique, marketing digital), en précisant pour chaque étape l'objectif et les moyens nécessaires.
\end{enumerate}
\end{exercice}

\newpage

\coursec{Corrigés des exercices}

\coursec{Les pratiques agropastorales et artisanales traditionnelles}
\courssub{Corrigés détaillés des exercices d'application et de synthèse}

\begin{corrige}[Exercice 1 — QCM : notions clés]
\begin{enumerate}
\item \textbf{Bonne réponse : b.} La transhumance est le \cle{déplacement saisonnier} des troupeaux entre des zones de pâturage (par exemple entre l'Adamaoua et les zones de saison sèche du Nord). La sédentarisation (a) est le phénomène inverse ; l'élevage industriel en bâtiment (c) relève de l'élevage intensif, non traditionnel.
\item \textbf{Bonne réponse : a.} Le \cle{ranching} est un élevage extensif moderne pratiqué sur de grands espaces clôturés, avec des races améliorées et un suivi vétérinaire (ex. : ranchs de l'Adamaoua gérés par la SODEPA).
\item \textbf{Bonne réponse : b.} Une \cle{mareyeuse} est une commerçante qui achète le poisson frais au débarquement et le revend sur les marchés ; elle joue un rôle central dans la commercialisation de la pêche artisanale.
\item \textbf{Bonne réponse : b.} Une \cle{indication géographique} est un signe qui lie la qualité ou la réputation d'un produit à son origine géographique (ex. : poivre de Penja, label IG enregistré à l'OAPI).
\end{enumerate}
\textbf{Points de vigilance :} ne pas confondre transhumance (déplacement saisonnier aller-retour) et nomadisme (déplacement permanent) ; ne pas confondre mareyeuse (commerçante) et pêcheuse (productrice) ; distinguer ranching (extensif moderne) et élevage intensif hors-sol.
\end{corrige}

\begin{corrige}[Exercice 2 — Vrai ou faux ? Justifie ta réponse]
\begin{enumerate}
\item \textbf{Vrai.} Dans le Nord-Ouest et l'Adamaoua, les conflits opposent fréquemment des cultivateurs sédentaires et des éleveurs peuls (Mbororo) pour l'accès à la terre et aux points d'eau, notamment lorsque les troupeaux détruisent les champs avant la récolte.
\item \textbf{Faux.} La pêche artisanale utilise des pirogues (souvent motorisées), des filets maillants, des lignes et des nasses ; les chalutiers relèvent de la pêche \emph{industrielle} (pêche de fond, forte puissance, équipage salarié).
\item \textbf{Faux.} L'artisanat traditionnel persiste et se transforme : poterie, vannerie, sculpture (Foumban), tissage (Maroua, ndop Bamiléké) trouvent de nouveaux débouchés (tourisme, exportation, labels), malgré la concurrence des produits industriels (plastique, aluminium).
\item \textbf{Vrai.} La modernisation passe par l'amélioration génétique des races (insémination, taureaux améliorateurs), la vaccination systématique (peste bovine, PPCB), la création de ranches et la construction de points d'eau pastoraux (politiques menées notamment par le MINEPIA et la SODEPA).
\end{enumerate}
\textbf{Points de vigilance :} une justification est exigée même pour une réponse « vrai » ; éviter les affirmations absolues (« totalement disparu », « exclusivement ») qui sont presque toujours fausses en géographie ; citer les acteurs institutionnels (MINEPIA, SODEPA, OAPI).
\end{corrige}

\begin{corrige}[Exercice 3 — Définitions]
\begin{itemize}
\item \textbf{Chaîne opératoire :} suite ordonnée d'opérations techniques qui transforment une matière première en produit fini (ex. pour la poterie : extraction de l'argile, triage, malaxage, façonnage, séchage, décor, cuisson).
\item \textbf{Pêche artisanale :} pêche pratiquée à petite échelle avec des embarcations simples (pirogues) et des engins peu mécanisés (filets, lignes, nasses), destinée à l'autoconsommation et au marché local ou national.
\item \textbf{Mareyeuse :} commerçante, souvent une femme, qui achète le poisson frais au débarquement et assure sa distribution (vente au détail, transformation par fumage/séchage, transport vers les marchés intérieurs).
\item \textbf{Artisanat :} activité de production manuelle ou peu mécanisée d'objets utilitaires ou décoratifs (poterie, vannerie, sculpture, tissage), reposant sur un savoir-faire traditionnel transmis de génération en génération.
\end{itemize}
\textbf{Points de vigilance :} une définition comporte le genre proche (activité, suite, commerçante) et le caractère distinctif (peu mécanisée, savoir-faire hérité, rôle commercial) ; éviter les définitions par l'exemple seul.
\end{corrige}

\begin{corrige}[Exercice 4 — Citations, localisation et calcul direct]
\begin{enumerate}
\item Deux exemples de conflits agro-pastoraux : la région du \textbf{Nord-Ouest} (cultivateurs contre éleveurs Mbororo autour des champs et des points d'eau) et la région de \textbf{l'Adamaoua} (occupation des couloirs de transhumance par les cultures, destruction des cultures par les troupeaux).
\item Deux engins de pêche interdits : les \textbf{filets à petites mailles} (non sélectifs, capturent les juvéniles) et la \textbf{pêche à la dynamite} (destructrice d'habitats) ; on peut aussi citer les \textbf{monofilaments} (filets invisibles, très capturants).
\item Deux matières premières artisanales menacées : le \textbf{rotin} (surexploitation forestière, disparition des lianes en forêt dense) et les \textbf{essences de bois précieux} utilisées en sculpture (ébène, dibétou) ; on peut aussi citer certaines argiles de qualité pour la poterie (gisements localisés, épuisement).
\item Calcul du taux de variation des captures artisanales (2000–2023) :
\[
T = \frac{V_{2023}-V_{2000}}{V_{2000}}\times 100 = \frac{135-90}{90}\times 100 = \frac{45}{90}\times 100 = 50\,\%.
\]
Les captures artisanales ont augmenté de \textbf{50\,\%} entre 2000 et 2023.
\end{enumerate}
\textbf{Points de vigilance :} le taux de variation se calcule toujours par rapport à la valeur de départ (année 2000), jamais par rapport à la valeur d'arrivée ; ne pas oublier le signe et l'unité (\%) ; localiser précisément les exemples (régions, fleuves, essences).
\end{corrige}

\begin{corrige}[Exercice 5 — L'élevage traditionnel et ses transformations]
\begin{enumerate}
\item \textbf{Élevage extensif nomade/transhumant :} pratiqué dans l'Adamaoua, le Nord et l'Extrême-Nord ; espèces dominantes : bovins (zébu), ovins et caprins ; l'espace est organisé autour des couloirs de transhumance, des pâturages de saison sèche (bas-fonds, yaérés) et de saison humide (plateaux), et des points d'eau, les troupeaux se déplaçant selon les saisons.
\textbf{Élevage sédentaire :} pratiqué dans l'Ouest, le Centre, le Sud et le Littoral ; espèces : volailles, porcs, petits ruminants, parfois bovins en stabulation ; l'espace est organisé autour de la concession familiale, avec parcs, cases d'élevage, cultures fourragères associées et fumure des champs.
\item Deux transformations majeures : (a) la \textbf{modernisation technique} — amélioration des races (croissement zébu × taurins), vaccination systématique, création de ranches clôturés (Adamaoua : ranchs de la SODEPA), forages pastoraux ; (b) la \textbf{sédentarisation encadrée} encouragée par les politiques publiques (MINEPIA, SODEPA, projets d'appui), qui fixe les éleveurs et développe l'agropastoralisme (association agriculture-élevage sur un même espace).
\item \textbf{Conséquence positive :} augmentation de la production de viande et de lait, amélioration des revenus des éleveurs, sécurisation du cheptel. \textbf{Conséquence négative :} accentuation des conflits d'accès à la terre et à l'eau, pression sur les pâturages autour des points d'eau permanents (surpâturage, dégradation des sols), perte de biodiversité pastorale.
\end{enumerate}
\textbf{Points de vigilance :} toujours localiser les exemples (régions, espèces, infrastructures) ; distinguer clairement système nomade/transhumant et système sédentaire par l'organisation de l'espace (mobilité vs fixation), pas seulement par les espèces.
\end{corrige}

\begin{corrige}[Exercice 6 — La pêche traditionnelle en mutation]
\begin{enumerate}
\item Deux techniques : (a) la \textbf{pêche au filet maillant} depuis une pirogue : l'équipage dérive en laissant le filet capturer le poisson par les ouïes (maillage adapté à l'espèce ciblée), puis le relève à la main ; le travail est collectif (patron de pirogue, matelots, partage de la capture). (b) La \textbf{pêche à la ligne et aux nasses} : engins individuels placés dans les fleuves (Wouri, Sanaga, Logone, Chari) ou les lagunes (Ébrié), relevés périodiquement, peu coûteux en capital.
\item Transformations récentes : la \textbf{motorisation} des pirogues (moteurs hors-bord 15–40 ch) élargit le rayon d'action (jusqu'au large pour la pêche maritime) et réduit la fatigue ; l'usage de la \textbf{glace} (glacières, usines à glace à Kribi, Douala, Limbé) permet de conserver le poisson frais et d'atteindre des marchés lointains (Yaoundé, Bafoussam) ; les \textbf{mareyeuses} organisent des circuits de commercialisation efficaces (achat au débarquement, fumage/séchage, transport routier vers l'intérieur du pays), intégrant la pêche artisanale à l'économie monétaire nationale.
\item La surexploitation s'explique par la croissance démographique et la demande urbaine croissante, la multiplication des pirogues motorisées (augmentation de l'effort de pêche), l'usage d'engins non sélectifs (petites mailles, monofilaments) qui capturent les juvéniles avant reproduction, et la concurrence spatiale des chalutiers industriels (pêche de fond, crevettiers) dans les zones côtières. Résultat : baisse des captures par unité d'effort (CPUE) et raréfaction des espèces à forte valeur commerciale (crevettes, gros poissons démersaux).
\end{enumerate}
\textbf{Points de vigilance :} distinguer pêche artisanale maritime (Kribi, Limbé, Douala, côte atlantique) et continentale (lac Tchad, Logone, Chari, Wouri, Sanaga, lagunes) ; relier chaque transformation technique (moteur, glace, mareyeuses) à son effet sur l'espace (rayon d'action, marchés) et le marché (prix, revenus).
\end{corrige}

\begin{corrige}[Exercice 7 — L'évolution de l'artisanat]
\begin{enumerate}
\item Quatre filières emblématiques et leurs zones : la \textbf{poterie} (Foumban, Baham, régions du Nord — Maroua, Guider) ; la \textbf{vannerie} (région du Centre — Yaoundé, Mbalmayo, Littoral — Douala, rotin et raphia) ; la \textbf{sculpture sur bois} (Foumban, Bafoussam, Baham, pays Bamiléké) ; le \textbf{tissage} (Maroua et le Grand Nord — coton, ndop et tissus de raphia chez les Bamiléké).
\item Chaîne opératoire de la poterie (ex. Foumban) : extraction de l'argile (gisement local) $\rightarrow$ triage et malaxage/dégraissage (ajout de sable, grog) $\rightarrow$ façonnage (modelage à la main, colombins, tour manuel) $\rightarrow$ séchage progressif à l'ombre $\rightarrow$ décor (engobes, motifs géométriques, incisions) $\rightarrow$ cuisson (foyer ouvert traditionnel ou four amélioré type four à bois/charbon) $\rightarrow$ produit fini (pot, jarre, canari, tuile) destiné à la vente locale ou touristique.
\item Deux évolutions contrastées : (a) l'ouverture au \textbf{marché touristique et à l'exportation}, qui valorise les produits haut de gamme (sculptures de Foumban, masques, ndop, poteries décorées vendues aux touristes et à l'étranger via foires, galeries, plateformes numériques) ; (b) la \textbf{concurrence industrielle} (plastiques, aluminium, tissus imprimés) qui réduit la demande d'objets utilitaires quotidiens, compensée partiellement par la montée des \textbf{labels et indications géographiques} (OAPI) qui protègent et valorisent le savoir-faire (ex. : poivre de Penja, futur label « Poterie de Foumban »).
\end{enumerate}
\textbf{Points de vigilance :} associer systématiquement filière et zone de production (terroir) ; la chaîne opératoire doit être ordonnée, de la matière première au produit fini, avec le vocabulaire technique (dégraissage, engobe, cuisson).
\end{corrige}

\begin{corrige}[Exercice 8 — Étude de cas : les conflits agro-pastoraux]
\begin{enumerate}
\item Deux causes structurelles : la \textbf{rareté et la compétition pour la terre et l'eau} (pression démographique, extension des champs vivriers et de rente sur les pâturages et les couloirs de transhumance) ; la \textbf{destruction des cultures par les troupeaux} (passage hors couloirs, période de culture) et l'accaparement des points d'eau en saison sèche par l'un ou l'autre groupe.
\item Deux solutions concrètes et spatialisées : la \textbf{délimitation matérialisée des couloirs de transhumance} (bornes, haies vives) et des zones de pâturage réservées, avec des points d'eau séparés (forages pastoraux distincts des forages villageois) ; la création et l'animation de \textbf{comités de dialogue} agriculteurs-éleveurs (chefs traditionnels, autorités administratives, MINEPIA, MINADER, ONG) pour prévenir et régler les litiges à l'amiable avant l'intervention judiciaire.
\item \textbf{Paragraphe argumenté :} Loin d'être rivales, l'agriculture et l'élevage peuvent être complémentaires dans une logique d'agropastoralisme intégré. Les troupeaux fournissent du fumier qui fertilise les champs et améliore les rendements céréaliers, tandis que les résidus de récolte (pailles de mil/sorgho/maïs, fanes d'arachide, tourteaux) nourrissent le bétail après les moissons : c'est le principe du contrat de pâturage ou de l'alliance « fumier contre fourrage ». L'élevage procure en outre un revenu monétaire (vente bétail) et une force de traction animale pour le labour et le transport. Dans l'Adamaoua, la sédentarisation partielle des éleveurs Mbororo montre que l'association cultures vivrières (maïs, igname) et bovins sur un même espace est viable. À condition de délimiter strictement les couloirs de transhumance, d'aménager des points d'eau dédiés et d'instituer un dialogue permanent, la coexistence devient un facteur de développement local durable plutôt qu'une source de conflits violents.
\end{enumerate}
\textbf{Points de vigilance :} au Bac, les solutions doivent être concrètes, spatialisées (couloirs, points d'eau, bornage) et institutionnelles (comités, administration) ; le paragraphe doit contenir une thèse claire, des arguments causaux (fumier, résidus, traction, revenu) et au moins un exemple localisé (Adamaoua, Nord-Ouest).
\end{corrige}

\begin{corrige}[Exercice 9 — Dissertation structurée (type Bac)]
\textbf{Barème indicatif (sur 20) :} introduction (4 pts : accroche 1, définitions 1, problématique 1, annonce du plan 1) ; développement (12 pts : deux parties équilibrées de 6 pts, avec exemples localisés et données chiffrées) ; conclusion (2 pts) ; présentation, vocabulaire géographique et rigueur (2 pts).

\textbf{Introduction.} \emph{Accroche :} L'élevage, la pêche et l'artisanat occupent encore la majorité de la population rurale camerounaise et assurent une part essentielle de la sécurité alimentaire et des revenus. \emph{Définitions :} Les pratiques agropastorales et artisanales traditionnelles désignent des activités de production fondées sur des savoir-faire hérités, une faible mécanisation, une main-d'œuvre familiale et une insertion historique dans les circuits locaux. \emph{Problématique :} Comment ces pratiques héritées font-elles face à la modernisation technique, à l'ouverture des marchés et aux politiques publiques ? \emph{Annonce du plan :} Après avoir montré leur persistance et leur diversité (I), nous analyserons leurs transformations face aux mutations contemporaines (II).

\textbf{I. Persistance et diversité des pratiques traditionnelles.}
L'élevage extensif transhumant domine dans les savanes de l'Adamaoua, du Nord et de l'Extrême-Nord (zébus, ovins, caprins), organisé autour des couloirs de transhumance et des points d'eau saisonniers, tandis que l'élevage sédentaire (volailles, porcs, petits ruminants) structure les campagnes de l'Ouest, du Centre et du Sud. La pêche artisanale, maritime (Kribi, Limbé, Douala — 135 000 t en 2023) et continentale (lac Tchad, Logone, Chari, Wouri, Sanaga), fournit l'essentiel du poisson consommé au Cameroun. L'artisanat reste vivace : poterie et sculpture à Foumban (Noun), vannerie au Centre (rotin), tissage à Maroua (coton) et ndop chez les Bamiléké. Ces activités reposent sur des chaînes opératoires héritées, une forte organisation sociale du travail (groupements, corporations) et une adaptation aux milieux (savane, forêt, littoral).

\textbf{II. Transformations face aux mutations contemporaines.}
La modernisation touche l'élevage : amélioration génétique, prophylaxie vétérinaire, ranches de l'Adamaoua (SODEPA), sédentarisation encouragée par les politiques publiques (MINEPIA). La pêche artisanale se transforme : motorisation des pirogues, chaîne du froid (glace), circuits des mareyeuses, mais aussi surexploitation des ressources (engins interdits, juvéniles) et conflits avec la pêche industrielle (chalutiers). L'artisanat s'ouvre au marché touristique et à l'exportation, affronte la concurrence industrielle (plastiques, tissus imprimés) et cherche la protection par les labels et indications géographiques (OAPI). Enfin, les conflits agriculteurs-éleveurs (Nord-Ouest, Adamaoua) révèlent les tensions foncières et appellent une gestion concertée de l'espace (couloirs, comités de dialogue).

\textbf{Conclusion.} Les pratiques traditionnelles camerounaises ne disparaissent pas : elles s'adaptent, entre persistance des savoir-faire et intégration à l'économie de marché. \emph{Ouverture :} Leur avenir dépend d'un développement durable conciliant productivité accrue, gestion concertée des ressources (repos biologique, couloirs de transhumance, quotas de pêche) et valorisation économique des savoir-faire (labels, tourisme, numérique).

\textbf{Points de vigilance :} annoncer le plan explicitement en fin d'introduction ; équilibrer les deux parties (environ 15 lignes chacune) ; citer au moins un exemple localisé et un chiffre par paragraphe ; éviter le catalogue descriptif sans analyse ; soigner le vocabulaire (transhumance, ranching, mareyeuse, chaîne opératoire, indication géographique, agropastoralisme).
\end{corrige}

\begin{corrige}[Exercice 10 — Croquis légendé et organisé]
\textbf{Barème indicatif (sur 10) :} fond de carte correct avec titre et nord (2 pts) ; trois jeux de symboles bien localisés (4 pts) ; légende hiérarchisée et complète (2 pts) ; commentaire géographique pertinent (2 pts).

\textbf{Éléments attendus sur le croquis :}
\begin{itemize}
\item Fond de carte : contour simplifié du Cameroun, flèche du nord, titre précis : « Les grandes activités agropastorales et artisanales du Cameroun ».
\item Symboles d'aire (aplats colorés ou hachures) : grandes zones d'élevage extensif en Adamaoua, Nord et Extrême-Nord (savanes).
\item Symboles linéaires (traits bleus) et ponctuels (points bleus) : zones de pêche artisanale — côte atlantique (Douala, Limbé, Kribi) ; lac Tchad (extrême nord) ; fleuves Logone, Chari, Wouri, Sanaga.
\item Symboles ponctuels (étoiles ou carrés rouges nommés) : pôles artisanaux majeurs — Foumban (poterie, sculpture), Bafoussam (sculpture, ndop), Maroua (tissage, vannerie).
\item Légende hiérarchisée : 1. Élevage (aires), 2. Pêche (linéaires/ponctuels), 3. Artisanat (ponctuels).
\end{itemize}

\begin{popfigure}
\begin{tikzpicture}[scale=0.85]
% Fond de carte simplifié Cameroun (polygone approximatif)
\draw[popdark, thick] (0,0) -- (1,3.5) -- (3,4.5) -- (5.5,4) -- (7,2.5) -- (7.5,0.5) -- (6.5,-1) -- (4.5,-1.5) -- (2.5,-1) -- (0.5,-0.5) -- cycle;
% Nord
\fill[popgreen!30, opacity=0.6] (1.5,2.5) -- (3,4.5) -- (5.5,4) -- (5,2) -- (3.5,1.5) -- (2,2) -- cycle;
\node[popdark, font=\small\bfseries] at (3.5,3.2) {Élevage extensif};
% Lac Tchad
\fill[popblue!40] (4.5,4.2) ellipse ({0.6} and {0.3});
\node[popblue, font=\tiny] at (4.5,4.6) {Lac Tchad};
% Fleuves
\draw[popblue, thick] (3.5,0.5) .. controls (4,1.5) and (4.5,2.5) .. (4.5,4.2); % Logone/Chari
\draw[popblue, thick] (2.5,1) .. controls (3,2) and (3.5,3) .. (4,3.5); % Bénoué approx
\draw[popblue, thick] (5.5,2) .. controls (6,1) .. (6.5,0); % Sanaga/Wouri vers Douala
% Côte
\draw[popblue, very thick] (7,2.5) -- (7.5,0.5);
% Villes pêche
\node[circle, fill=popblue, inner sep=1.5pt, label=right:\tiny Douala] at (6.8,1.2) {};
\node[circle, fill=popblue, inner sep=1.5pt, label=right:\tiny Limbé] at (7.2,1.8) {};
\node[circle, fill=popblue, inner sep=1.5pt, label=right:\tiny Kribi] at (7.1,0.2) {};
% Artisanat
\draw[poporange, fill=poporange] (2.2,1.2) circle (2.5pt) node[above left, font=\tiny, popdark] {Foumban};
\draw[poporange, fill=poporange] (2.5,0.8) circle (2.5pt) node[below right, font=\tiny, popdark] {Bafoussam};
\draw[poporange, fill=poporange] (5.2,3.2) circle (2.5pt) node[above right, font=\tiny, popdark] {Maroua};
% Nord
\node[popdark, rotate=45, font=\Large] at (0.5,4.5) {N};
% Titre
\node[popdark, font=\normalsize\bfseries, align=center] at (3.7,-2.2) {Les grandes activités agropastorales et artisanales du Cameroun};
\end{tikzpicture}
\legende{Croquis schématique : aires d'élevage (vert), zones de pêche (bleu), pôles artisanaux (orange). Localisation approximative.}
\end{popfigure}

\textbf{Commentaire type (5 lignes) :} Le croquis montre une nette opposition nord-sud : l'élevage extensif domine dans les savanes du Nord (Adamaoua, Extrême-Nord), tandis que la pêche artisanale se concentre sur la côte atlantique (Douala, Limbé, Kribi) et les grands hydrosystèmes (lac Tchad, Logone, Wouri, Sanaga). Les pôles artisanaux (Foumban, Bafoussam, Maroua) s'ancrent dans les villes historiques, témoins de savoir-faire anciens (royaumes, chefferies). Cette répartition traduit l'influence du milieu physique (climat, hydrographie, végétation) et de l'histoire des peuplements sur l'organisation de l'espace productif.

\textbf{Points de vigilance :} un croquis sans titre, sans flèche du nord ou sans légende est sanctionné (0 pt) ; les symboles doivent être placés au bon endroit (le lac Tchad à l'extrême nord, Kribi au sud de la côte, Foumban à l'Ouest) ; la légende ne doit contenir que ce qui figure sur le croquis (pas de symbole absent de la carte) ; respecter l'ordre hiérarchique : aires > linéaires > ponctuels.
\end{corrige}

\begin{corrige}[Exercice 11 — Analyse de documents et cas pratique (type Bac)]
\textbf{Barème indicatif (sur 20) :} calculs (4 pts) ; explication de la surexploitation (6 pts) ; mesures de gestion (4 pts) ; cas pratique (6 pts).

\begin{enumerate}
\item \textbf{Calculs des taux de variation (2000–2023).}
Pêche artisanale :
\[
T_{a}=\frac{135-90}{90}\times 100=\frac{45}{90}\times 100=50\,\%.
\]
Pêche industrielle :
\[
T_{i}=\frac{105-60}{60}\times 100=\frac{45}{60}\times 100=75\,\%.
\]
\textbf{Comparaison :} Les deux filières progressent en volume absolu, mais la pêche industrielle croît plus vite en proportion (+75\,\% contre +50\,\%). La pêche artisanale reste toutefois la première pourvoyeuse de poisson pour la consommation nationale (135 000 t contre 105 000 t en 2023).

\begin{popfigure}
\begin{tikzpicture}
\begin{axis}[
    ybar,
    bar width=12pt,
    width=\linewidth,
    height=6cm,
    ylabel={Captures (1000 tonnes)},
    xtick={1,2},
    xticklabels={2000, 2023},
    legend style={at={(0.5,1.05)}, anchor=south, legend columns=2},
    ymin=0, ymax=150,
    nodes near coords,
    nodes near coords align={vertical},
    popblueL,
]
\addplot[popblue, fill=popblue!60] coordinates {(1,90) (2,135)};
\addplot[poporange, fill=poporange!60] coordinates {(1,60) (2,105)};
\legend{Pêche artisanale, Pêche industrielle}
\end{axis}
\end{tikzpicture}
\legende{Évolution des captures de la pêche artisanale et industrielle au Cameroun (2000–2023). Sources : MINEPIA, FAO (données arrondies).}
\end{popfigure}

\item \textbf{Causes de la surexploitation (document 2 et connaissances) :} croissance de la demande (démographie, urbanisation, pouvoir d'achat) ; multiplication des pirogues motorisées et augmentation de l'effort de pêche ; usage d'engins non sélectifs (filets à petites mailles, dynamite, monofilaments — document 2) qui capturent les juvéniles avant reproduction ; concurrence des chalutiers industriels dans les zones de nurserie ; non-respect des périodes de reproduction et des zones de frayère ; gouvernance faible (contrôles insuffisants).
\item \textbf{Deux mesures de gestion durable (appuyées sur le document 2) :} (a) Faire respecter strictement l'interdiction des engins non sélectifs (contrôles aux débarquements, sanctions, mailles réglementaires minimales) pour protéger les juvéniles et le stock reproducteur ; (b) Appliquer les \textbf{périodes de repos biologique} (fermeture de la pêche pendant la saison de reproduction, ex. : juillet–septembre pour la crevette) afin de permettre le renouvellement des stocks.
\item \textbf{Cas pratique — Plan d'action pour la coopérative de potiers de Foumban (4 étapes) :}
\begin{enumerate}
\item \textbf{Diagnostic et sécurisation de la matière première :} Objectif — garantir un approvisionnement durable en argile de qualité ; Moyens — prospection géologique des gisements, tests de plasticité et de cuisson, convention d'exploitation raisonnée avec les propriétaires terriens et la commune, rotation des zones d'extraction.
\item \textbf{Amélioration de la cuisson et de la qualité :} Objectif — réduire le taux de casse (actuellement 30–40\,\%), uniformiser la cuisson, économiser le bois ; Moyens — construction collective d'un four amélioré type four à bois/charbon à tirage forcé ou four à gaz, formation des potiers au contrôle de la température (pyromètres), normalisation des temps de cuisson.
\item \textbf{Obtention d'une Indication Géographique (IG) « Poterie de Foumban » :} Objectif — protéger l'origine, valoriser le prix, fidéliser la clientèle ; Moyens — élaboration d'un cahier des charges collectif (argile locale, formes traditionnelles, décor spécifique), dépôt du dossier auprès de l'OAPI (Organisation Africaine de la Propriété Intellectuelle), mise en place d'un organisme de défense et de gestion (ODG).
\item \textbf{Marketing digital et diversification des débouchés :} Objectif — atteindre les clients nationaux (Yaoundé, Douala) et internationaux (touristes, diaspora, galeries) ; Moyens — création d'un site vitrine / boutique en ligne (photos professionnelles, storytelling), présence sur réseaux sociaux (Instagram, Facebook), partenariats logistiques (DHL, EMS) pour l'expédition, participation aux foires (PROMOTE, SIAO) et plateformes d'artisanat équitable.
\end{enumerate}
\end{enumerate}
\textbf{Points de vigilance :} dans les calculs, diviser par la valeur de l'année de départ (2000) ; chaque mesure proposée doit s'appuyer explicitement sur le document 2 (engins interdits, repos biologique) ; au cas pratique, chaque étape exige un objectif clair ET des moyens concrets (acteurs, outils, délais), sans quoi la réponse est incomplète ; utiliser le vocabulaire précis (repos biologique, juvéniles, IG/OAPI, cahier des charges, ODG).
\end{corrige}

\newpage

\coursec{Fiche bilan}

\begin{popfigure}
\begin{tikzpicture}[
  mindmap,
  concept color=popblue,
  level 1 concept/.append style={level distance=4.5cm, sibling angle=72},
  level 2 concept/.append style={level distance=3cm, sibling angle=45},
  every node/.style={font=\small, text=popdark},
  grow cyclic,
  text width=2.8cm,
  align=center
]
\node[concept, text width=3.5cm] {\textbf{Pratiques agropastorales \& artisanales traditionnelles}}
  child[concept color=popgreen] {
    node[concept] {\textbf{1. Élevage traditionnel}}
    child { node[concept, text width=2.5cm] {Systèmes : nomadisme, transhumance, élevage sédentaire} }
    child { node[concept, text width=2.5cm] {Organisation de l'espace : parcours, couloirs, points d'eau} }
    child { node[concept, text width=2.5cm] {Acteurs : Peuls, Mbororo, Arabes Choa} }
  }
  child[concept color=poporange] {
    node[concept] {\textbf{2. Transformations de l'élevage}}
    child { node[concept, text width=2.5cm] {Modernisation : ranchs, amélioration génétique, alimentation} }
    child { node[concept, text width=2.5cm] {Conflits : agriculteurs/éleveurs, blocage couloirs} }
    child { node[concept, text width=2.5cm] {Politiques : PNDP, CODEPA, sédentarisation} }
  }
  child[concept color=popblue] {
    node[concept] {\textbf{3. Pêche traditionnelle}}
    child { node[concept, text width=2.5cm] {Techniques : nasses, filets, lignes, barrages} }
    child { node[concept, text width=2.5cm] {Espaces : lagunes (Wouri), fleuves (Sanaga, Bénoué), lacs (Tchad, Maga)} }
    child { node[concept, text width=2.5cm] {Mutations : motorisation, surpêche, aquaculture} }
  }
  child[concept color=poppurple] {
    node[concept] {\textbf{4. Artisanat traditionnel}}
    child { node[concept, text width=2.5cm] {Diversité : poterie, vannerie, forge, tissage, sculpture, teinture} }
    child { node[concept, text width=2.5cm] {Organisation : ateliers familiaux, coopératives, marchés locaux/touristiques} }
    child { node[concept, text width=2.5cm] {Évolution : design, e-commerce, formation, label "Made in Cameroon"} }
  }
  child[concept color=poppink] {
    node[concept] {\textbf{5. Enjeux transversaux}}
    child { node[concept, text width=2.5cm] {Développement durable : gestion ressources, savoir-faire} }
    child { node[concept, text width=2.5cm] {Valorisation économique : emploi jeunes, revenus, export} }
    child { node[concept, text width=2.5cm] {Patrimoine culturel : identité, tourisme, transmission} }
  };
\end{tikzpicture}
\legende{Carte mentale de synthèse : les pratiques agropastorales et artisanales traditionnelles au Cameroun.}
\end{popfigure}

\begin{bilan}
\courssub{Définitions clés}
\begin{itemize}
  \item \cle{Transhumance} : Migration saisonnière des troupeaux (et pasteurs) entre zones humides (saison sèche) et zones sèches (saison des pluies) suivant des \cle{couloirs de transhumance} balisés.
  \item \cle{Nomadisme pastoral} : Déplacement permanent et imprévisible du groupe domestique avec le troupeau à la recherche de pâturages et d'eau (ex. : Bororo).
  \item \cle{Élevage sédentaire / agropastoral} : Association culture-élevage sur un terroir fixe ; fumure animale et traction animale valorisées.
  \item \cle{Pêche artisanale} : Activité de capture utilisant des embarcations et engins de faible puissance, main-d'œuvre familiale, commercialisation locale.
  \item \cle{Artisanat de production} : Fabrication d'objets utilitaires (poteries, paniers, outils) vs \cle{artisanat d'art} (objets décoratifs, rituels, touristiques).
  \item \cle{Valeur ajoutée} : Différence entre valeur de vente et coût des matières premières ; enjeu majeur pour l'artisanat (transformation locale vs export brut).
\end{itemize}

\courssub{Repères chiffrés \& Données à connaître (Bac)}
\begin{center}
\begin{tabularx}{\linewidth}{|l|X|X|}\hline
\rowcolor{popblueL}
\textbf{Indicateur} & \textbf{Donnée clé (ordres de grandeur)} & \textbf{Exemple / Localisation} \\ \hline
Part de l'élevage dans le PIB primaire & $\approx 15\text{--}20\,\%$ & Nord, Adamaoua, Extrême-Nord \\ \hline
Effectif bovin national & $\sim 6\text{--}7$ millions de têtes & Majoritairement race \emph{Goudali}, \emph{Red Fulani} \\ \hline
Production halieutique artisanale & $> 70\,\%$ des captures totales & Pêche continentale (Logone, Sanaga, Lac Tchad) + maritime (Kribi, Limbé) \\ \hline
Emploi artisanat (secteur informel) & $\sim 20\,\%$ de la population active urbaine & Foumban (bronze), Baham (tissage), Maroua (cuir), Bafoussam (bois) \\ \hline
Taux d'accroissement demande protéines & $> 3\,\%$/an & Pression sur ressources pastorales \& halieutiques \\ \hline
\end{tabularx}
\end{center}

\courssub{Méthodes express (Bac)}
\begin{methode}[Analyser une carte / croquis d'espace pastoral]
  \begin{enumerate}
    \item Identifier la \cle{légende} (couleurs = types d'élevage / couloirs / points d'eau / zones de culture).
    \item Repérer les \cle{conflits d'usage} : superposition parcours / champs, blocage couloirs par cultures / barrages.
    \item Mesurer la \cle{distance} pâturage / point d'eau (contrainte de rentabilité).
    \item Conclure sur l'\cle{aménagement} : bornage couloirs, aires de repos, forages, marchés à bétail.
  \end{enumerate}
\end{methode}

\begin{methode}[Expliquer la mutation d'une pratique traditionnelle]
  \begin{enumerate}
    \item Facteurs \cle{endogènes} : croissance démographique, quête de revenu, transmission du savoir.
    \item Facteurs \cle{exogènes} : politiques publiques (PNDP, CODEPA), changement climatique (sécheresse), marché (urbanisation, export), technologie (moto-pompe, téléphone, GPS).
    \item Conséquences \cle{positives} (productivité, revenu) et \cle{négatives} (dégradation sols, perte biodiversité, conflits, déperdition savoir-faire).
    \item Pistes \cle{durables} : organisation professionnelle, certification, formation, aménagement concerté.
  \end{enumerate}
\end{methode}
\end{bilan}

\begin{aretenir}[Checklist avant le Bac]
  $\square$ Je sais définir et distinguer \textbf{nomadisme}, \textbf{transhumance} et \textbf{élevage sédentaire} avec exemples camerounais.\par
  $\square$ Je localise sur carte les \textbf{grands espaces pastoraux} (Adamaoua, Nord, Extrême-Nord) et les \textbf{couloirs de transhumance} majeurs.\par
  $\square$ J'explique les \textbf{causes des conflits agriculteurs/éleveurs} et cite 2 mesures d'apaisement (ex. : CODEPA, bornage).\par
  $\square$ Je décris 3 \textbf{techniques de pêche traditionnelle} adaptées au milieu (lagune, fleuve, lac).\par
  $\square$ Je cite les \textbf{principaux bassins de pêche artisanale} (Lac Tchad, Lagune Wouri, Sanaga, Bénoué, Côte).\par
  $\square$ Je relie \textbf{surpêche / destruction frayères} aux mutations techniques (filets mailles fines, barrages).\par
  $\square$ Je classe l'\textbf{artisanat} par matière (bois, cuir, métal, textile, argile) et par fonction (production / art).\par
  $\square$ Je nomme 3 \textbf{pôles artisanaux réputés} et leur spécialité (Foumban, Baham, Maroua, Bafoussam, etc.).\par
  $\square$ J'argumente sur la \textbf{valeur ajoutée} : transformation locale vs export matière première, rôle du design et du numérique.\par
  $\square$ Je mobilise la \textbf{démarche de développement durable} pour une filière (ex. : pastoralisme : gestion concertée parcours + valorisation lait/viande).\par
  $\square$ Je sais rédiger une \textbf{réponse structurée} : définition $\rightarrow$ localisation $\rightarrow$ facteurs $\rightarrow$ conséquences $\rightarrow$ enjeux/pistes.
\end{aretenir}

\findecours

\end{document}
